\documentclass[11pt]{article}

\usepackage[margin=1.1in]{geometry}

\usepackage{amssymb,mathtools,natbib}

% Anchored formal environments for causalsmith papers.
% First mandatory argument = obj_id (machine anchor joining the paper to the
% Lean crosswalk; invisible in the rendered PDF). Optional argument = name.
% Each environment also defines \label{obj:<obj_id>} so prose can use
% \cref{obj:T-1}; cleveref derives the printed kind and number from the target.
\usepackage{amsmath}
\usepackage{mathrsfs}
\usepackage{amsthm}
\usepackage{booktabs}
% Long displays may break across pages; let TeX stretch a little harder before an
% overfull \hbox so wide formulas/tables are less likely to run off the page.
\allowdisplaybreaks
\setlength{\emergencystretch}{3em}
\newtheorem{theorem}{Theorem}
\newtheorem{lemma}{Lemma}
\newtheorem{assumption}{Assumption}
\newtheorem{definition}{Definition}
\newtheorem{proposition}{Proposition}
\newtheorem{remark}{Remark}
\newtheorem{citedresult}{Cited result}
% The amsthm counter is `algorithmbox` (not `algorithm`) so a later `\usepackage{algorithm}` cannot clash.
\newcounter{algorithmbox}
\renewcommand{\thealgorithmbox}{\arabic{algorithmbox}}
\NewDocumentEnvironment{theoremv}{m o s}{\begin{theorem}\label{obj:#1}{\normalfont[\texttt{\detokenize{#1}}]\IfValueT{#2}{\ (#2)}\IfBooleanT{#3}{\verificationfootnotemark}\ }}{\end{theorem}}
\NewDocumentEnvironment{lemmav}{m o s}{\begin{lemma}\label{obj:#1}{\normalfont[\texttt{\detokenize{#1}}]\IfValueT{#2}{\ (#2)}\IfBooleanT{#3}{\verificationfootnotemark}\ }}{\end{lemma}}
\NewDocumentEnvironment{assumptionv}{m o s}{\begin{assumption}\label{obj:#1}{\normalfont[\texttt{\detokenize{#1}}]\IfValueT{#2}{\ (#2)}\IfBooleanT{#3}{\verificationfootnotemark}\ }}{\end{assumption}}
\NewDocumentEnvironment{definitionv}{m o s}{\begin{definition}\label{obj:#1}{\normalfont[\texttt{\detokenize{#1}}]\IfValueT{#2}{\ (#2)}\IfBooleanT{#3}{\verificationfootnotemark}\ }}{\end{definition}}
% A `propositionv` is a proved result tiered as a Proposition (theorem-grade, machine-verified).
\NewDocumentEnvironment{propositionv}{m o s}{\begin{proposition}\label{obj:#1}{\normalfont[\texttt{\detokenize{#1}}]\IfValueT{#2}{\ (#2)}\IfBooleanT{#3}{\verificationfootnotemark}\ }}{\end{proposition}}
% A `remarkv` holds NON-load-bearing interpretive / open-question content (e.g. a causalsmith `oeq:`
% open question). It is neither proved nor assumed; no theorem may depend on it.
\NewDocumentEnvironment{remarkv}{m o s}{\begin{remark}\label{obj:#1}{\normalfont[\texttt{\detokenize{#1}}]\IfValueT{#2}{\ (#2)}\IfBooleanT{#3}{\verificationfootnotemark}\ }}{\end{remark}}
% An `algorithmv` is a DEFINITION-kind object presented as a boxed, numbered-step procedure (an
% estimator / selection rule built in stages). Same anchor contract as `definitionv`; its body is
% ordinary LaTeX whose steps are `\begin{enumerate}` items, so tex2html needs no extra machinery.
% Presented in the CONVENTIONAL algorithm style readers expect from `algorithm`+`algorithmic`:
% a rule above, an "Algorithm N (Title)" caption line, a rule under the caption, the numbered
% steps, and a closing rule. It is NOT a float — the anchor contract requires the object to stay
% where the outline places it, and floats drift. The obj-id tag is suppressed here (it is
% bookkeeping, and a caption line carrying it reads as debug output in a finished paper).
\NewDocumentEnvironment{algorithmv}{m o s}%
  {\par\addvspace{\medskipamount}\noindent\rule{\linewidth}{0.8pt}\par\nobreak\vspace{2pt}%
   \refstepcounter{algorithmbox}\label{obj:#1}%
   \noindent\textbf{Algorithm \thealgorithmbox}\IfValueT{#2}{ \textnormal{(#2)}}%
   \IfBooleanT{#3}{\verificationfootnotemark}\par\nobreak\vspace{2pt}%
   \noindent\rule{\linewidth}{0.4pt}\par\nobreak\vspace{2pt}\normalfont}%
  {\par\nobreak\vspace{2pt}\noindent\rule{\linewidth}{0.8pt}\par\addvspace{\medskipamount}}
% A `citedv` is an IMPORTED external result the paper relies on but does NOT prove
% (a causalsmith `gate_class:"cited"` node). It is machine-checked against the cited
% source at F2.5, never proved here; the fixed provenance note makes that explicit
% so the environment can never be read as a contribution of this paper. The \citep
% attribution is supplied by the surrounding prose (P2), keyed to references.bib.
\NewDocumentEnvironment{citedv}{m o s}{\begin{citedresult}\label{obj:#1}{\normalfont[\texttt{\detokenize{#1}}]\IfValueT{#2}{\ (#2)}\IfBooleanT{#3}{\verificationfootnotemark}\ \textit{(imported result; machine-checked against the cited source, not proved here.)}\ }}{\end{citedresult}}
% \leanref{obj_id}{display text}: an inline first-mention link to a from-note object's
% verified Lean code. In the PDF it renders the display text plainly (the Lean link is a
% web-only affordance); tex2html turns it into a drawer-opening span.
\newcommand{\leanref}[2]{#2}
% For a theorem/lemma whose Lean declaration accepts a source-matched published proposition as an
% input, the starred anchored environment puts the mark in its heading; the generated text remains
% outside the frozen mathematical environment. Keep the combined macro for older paper bundles.
\newcommand{\verificationfootnotemark}{\footnotemark}
\newcommand{\verificationfootnotetext}[1]{\footnotetext{#1}}
\newcommand{\verificationfootnote}[1]{\verificationfootnotemark\verificationfootnotetext{#1}}
% xcolor before hyperref (hyperref would otherwise pull in plain `color` for colorlinks, and a
% later xcolor load clashes). The page-1 identity stamp at the end of this file needs a grey.
\usepackage{xcolor}
\usepackage{graphicx}
\usepackage{eso-pic}
% hyperref follows natbib and the environment macros; cleveref must follow hyperref.
\usepackage[colorlinks=true,linkcolor=blue,citecolor=blue,urlcolor=blue]{hyperref}
\usepackage[capitalise,nameinlink,noabbrev]{cleveref}
% Pin the names explicitly: labels are placed inside the underlying amsthm environments, and these
% declarations make the PDF contract independent of cleveref's theorem-name inference. House style:
% reference names always print with a capital first letter ("Theorem 1", "Definition 2"), in any
% sentence position — hence `capitalise` above and capitalized \crefname forms below.
\crefname{theorem}{Theorem}{Theorems}
\Crefname{theorem}{Theorem}{Theorems}
\crefname{lemma}{Lemma}{Lemmas}
\Crefname{lemma}{Lemma}{Lemmas}
\crefname{assumption}{Assumption}{Assumptions}
\Crefname{assumption}{Assumption}{Assumptions}
\crefname{definition}{Definition}{Definitions}
\Crefname{definition}{Definition}{Definitions}
\crefname{proposition}{Proposition}{Propositions}
\Crefname{proposition}{Proposition}{Propositions}
\crefname{remark}{Remark}{Remarks}
\Crefname{remark}{Remark}{Remarks}
\crefname{citedresult}{Cited result}{Cited results}
\Crefname{citedresult}{Cited result}{Cited results}
\crefname{algorithmbox}{Algorithm}{Algorithms}
\Crefname{algorithmbox}{Algorithm}{Algorithms}
% ---------------------------------------------------------------------------
% Page-1 identity stamp (arXiv style) and the pinned title-block date.
%
% Split of responsibility: THIS file owns the FORMAT (placement, font, colour, punctuation),
% the generated `paper_stamp.tex` owns only the VALUES (number+version, area, revised date,
% DOI, link target). P4 refreshes this template on every emit, so a layout fix reaches every
% bundle from a recompile; and because `paper_stamp.tex` is not one of the P2 assembly
% digest's authored inputs, a DOI that arrives after publication needs only that one small
% file rewritten plus a recompile — never a redraft, never a reassembly.
%
% Defaults first, so a bundle WITHOUT a stamp file compiles exactly as it always did:
% `\cspaperdate` falls back to `\today` and no stamp is drawn.
\providecommand*{\cspaperdate}{\today}  % the \date{} target
\providecommand*{\csstampid}{}          % e.g. CSWP-2026-008v3 (empty ⇒ no stamp at all)
\providecommand*{\csstamparea}{}        % e.g. Stat
\providecommand*{\csstampdate}{}        % e.g. 27 Aug 2026
\providecommand*{\csstampdoi}{}         % e.g. 10.5281/zenodo.123 (empty ⇒ DOI part omitted)
\providecommand*{\csstamplink}{}        % href target (DOI url, else the paper's page; may be empty)
\IfFileExists{paper_stamp.tex}{% GENERATED by CausalSmith P4 — paper identity stamp values. Do not hand-edit.
%
% Field order on the stamp (owned by paper_macros.tex):
%   <number>v<version> · doi:<doi> · [<Area>] <date>   — the DOI is never the last token, so a
%   text extractor cannot run it into whatever follows. Without a DOI that part is omitted.
% Format lives in paper_macros.tex; only the values are here, and this file is NOT one of
% the P2 assembly digest's authored inputs. A DOI arriving after publication therefore needs
% this file rewritten plus a `latexmk` recompile — never a redraft or a reassembly.
\renewcommand*{\csstampid}{CSWP-2026-035v7}
\renewcommand*{\csstamparea}{Stat}
\renewcommand*{\csstampdate}{4 Oct 2026}
\renewcommand*{\csstampdoi}{}
\renewcommand*{\csstamplink}{https://causalsmith.org/papers/stat_twosample_pointcate_annotation_frontier_v1}
\renewcommand*{\cspaperdate}{4 October 2026}
}{}
\makeatletter
% Field order is deliberate: the DOI is NEVER the last token on the line. A trailing identifier
% runs into whatever a text extractor emits next, which makes it unsafe to match mechanically
% (the Zenodo stamp matcher reads exactly this line); the date is a safe terminator.
%   <number>v<version> · doi:<doi> · [<Area>] <date>      — with a DOI
%   <number>v<version> · [<Area>] <date>                  — without
\newcommand{\cs@stampsep}{\,\textperiodcentered\,}
\newcommand{\cs@stamptext}{%
  \csstampid
  \ifx\csstampdoi\@empty\else\cs@stampsep doi:\csstampdoi\fi
  \ifx\csstamparea\@empty\else\cs@stampsep[\csstamparea]\fi
  \ifx\csstampdate\@empty\else\quad\csstampdate\fi
}
\ifx\csstampid\@empty\else
  \definecolor{csstampgrey}{gray}{0.45}
  % Background shipout material: it occupies no space, so the text block and the \thanks
  % footnote are byte-identical to an unstamped compile. The starred form applies to the
  % NEXT page shipped only — i.e. page 1. Placed in the left margin (the text block starts
  % at 1.1in), reading bottom-to-top, in a Times-like face at 8pt.
  \AddToShipoutPictureBG*{%
    \AtPageLowerLeft{%
      \put(\LenToUnit{0.42in},\LenToUnit{1.1in}){%
        \rotatebox{90}{%
          \fontfamily{ptm}\fontsize{8}{10}\selectfont
          \ifx\csstamplink\@empty
            \textcolor{csstampgrey}{\cs@stamptext}%
          \else
            \expandafter\href\expandafter{\csstamplink}{\textcolor{csstampgrey}{\cs@stamptext}}%
          \fi
        }%
      }%
    }%
  }
\fi
\makeatother


\title{Outcome-free records and the minimax precision of conditional treatment effect estimation}

\author{CausalSmith\thanks{Machine-generated by the CausalSmith research pipeline.}}

\date{\cspaperdate}

\begin{document}

\maketitle

\begin{abstract}
We characterize the value of outcome-free covariate-treatment records for estimating a conditional average treatment effect at the center of the unit covariate cube. The experiment combines independent labeled and outcome-free samples from a common population. Under a known uniform covariate design, conditional exchangeability, a public overlap level \(0<\varepsilon<1/2\) with propensity values in \([\varepsilon,1-\varepsilon]\), binary responses with arm means in \([0,1]\), and public Hölder regularity, we establish the sharp minimax expected absolute-error rate uniformly over labeled counts of at least two and every nonnegative auxiliary count. The propensity and control-mean smoothness orders lie in \((0,1]\), the contrast smoothness order lies in \([1,3]\), and the common public Hölder radius satisfies \(L>1/2\). The rate is the maximum of a labeled-sample term and a term involving the product of labeled and total treatment-record counts. A guarded local polynomial estimator with a rectangular projection correction attains this rate. The matching lower bound combines fixed-size mixtures of the original records, which govern the product-of-counts regime, with a two-hypothesis bound under a supplied propensity, which gives the labeled-sample floor. The characterization yields necessary and sufficient sample-growth conditions for attaining the supplied-propensity order and for improving strictly over labeled records alone. It also gives simultaneous acquisition requirements on response labels and total treatment records for a target error scale. The rate exponents do not depend on \(\varepsilon\), while the comparison constants may depend on it together with the other fixed public parameters.
\end{abstract}

\section{Introduction}\label{sec:introduction}

This paper establishes a sharp pointwise precision benchmark for conditional average treatment effect estimation when responses are observed on a smaller collection of records than covariates and treatment assignments. We consider \(n\ge2\), the labeled-record count, and \(m\ge0\), the independent outcome-free covariate-treatment record count. For the specified potential-outcome class in \cref{obj:def:primitive-class}, with known uniform covariate design and the public parameter restrictions stated below, the minimax expected absolute error is
\[
\underbrace{R(n,m)}_{\text{minimax pointwise absolute-error risk}}
\asymp
\max\left\{
n^{-\gamma/(2\gamma+d)},
[n(n+m)]^{-\gamma/(2\gamma+d+\gamma d/(\alpha+\beta))}
\right\},
\]
where \(d\) is the covariate dimension, and \(\alpha,\beta,\gamma\) are the Hölder smoothness orders of the propensity, control response mean, and causal contrast, respectively. The target is the conditional causal contrast at the cube's center. The comparison holds uniformly over both record counts, with constants depending on the fixed public class parameters. A fully specified decision based on the observed records attains this order, as established in \cref{obj:thm:sharp-annotation-frontier}.

The model fixes the covariate distribution to be uniform on \([0,1]^d\), with \(d\ge1\), and assumes conditional treatment exchangeability. Treatment and potential responses are binary. For a fixed public \(\varepsilon\in(0,1/2)\), the propensity lies in \([\varepsilon,1-\varepsilon]\), and both potential-response means lie in \([0,1]\) throughout the cube. The public smoothness orders satisfy \(\alpha,\beta\in(0,1]\) and \(\gamma\in[1,3]\), with \(L\), the common public Hölder radius, satisfying \(L>1/2\). These restrictions define the primitive law class in \cref{obj:def:primitive-class}; the independent, common-population record blocks follow \cref{obj:ass:sample-law}. Admissible conditional-margin witnesses satisfy the restrictions over the whole cube. Their continuous contrast, together with full-support design and the identifying assumptions, determines the point target uniquely from the observed law.

The distinction between record types has a direct acquisition interpretation. Response ascertainment through chart review can yield labeled records containing covariates, treatment, and the true response, while routine records supply covariates and treatment. Surrogate-assisted treatment-effect estimation provides a related electronic-record setting \citep[Sections 2.1, 2.5--2.6, and 4]{ChengAnanthakrishnanCai2021}. Within the present sampling model, adding a new independent labeled record changes \((n,m)\) to \((n+1,m)\). Labeling one already-counted outcome-free record changes the composition to \((n+1,m-1)\) at fixed total count; if the record index is fixed in advance, or chosen by a randomization independent of all records, the resulting dataset follows the stated sampling law at \((n+1,m-1)\). Adding a new outcome-free record changes \((n,m)\) to \((n,m+1)\). The sharp risk formula benchmarks these sample compositions under the stated independent-sampling assumptions.

The estimator combines a coarse local polynomial fit with a finer projection correction. It splits the original records into independent outcome and treatment roles, whose sizes are proportional to the labeled and total treatment-record counts. A rectangular average across these roles corrects the local estimating moments for the unknown nuisance functions. An empirical eigenvalue guard and clipping produce a finite Borel decision for every dataset. Under the stated model conditions, \cref{obj:thm:rectangular-upper} bounds its risk by effect approximation, nuisance approximation, localized response sampling, and rectangular sampling terms. Public tuning balances these terms to attain the sharp rate.

The matching lower bound retains the original fixed-size observation experiment. The construction in \cref{obj:thm:marked-component-certificate} couples propensity and response perturbations while preserving identical auxiliary-only mixture laws across hypotheses. Its information bound supplies the nuisance converse in \cref{obj:lem:nuisance-frontier-converse} over the range where the product-of-counts term governs precision. The supplied-propensity comparison in \cref{obj:thm:supplied-propensity} establishes the labeled-sample order \(n^{-\gamma/(2\gamma+d)}\) over the same primitive class and every auxiliary count. Together, these lower bounds and the attaining estimator give the uniform minimax characterization, including randomized decisions.

The characterization distinguishes strict improvement from attainment of the supplied-propensity order. Write \(S=\alpha+\beta\), the combined nuisance smoothness, and let \(S_{\mathrm{crit}}\) denote the smoothness threshold and \(q_*\) the total-record growth exponent:
\[
S_{\mathrm{crit}}=\frac{\gamma d}{2\gamma+d},
\qquad
q_*=\frac{\gamma d}{S(2\gamma+d)}.
\]
At and above the threshold, labeled records alone attain the supplied-propensity order. Below it, attainment along an auxiliary-size sequence requires total treatment-record counts to remain at least a positive constant times \(n^{q_*}\). Below the smoothness threshold, strict asymptotic improvement over labeled records alone occurs exactly when the auxiliary-to-labeled ratio diverges. \Cref{obj:prop:annotation-threshold} also translates a target absolute-error scale into simultaneous requirements on the labeled count and the product of labeled and total counts. These are risk requirements up to parameter-dependent constants, with exact equivalences for the displayed rate.

The closest supervised comparison is \citet[Section 3, Theorems 1--2 and equation (14)]{KennedyBalakrishnanRobinsWasserman2024}. At zero auxiliary records, our rate has the same numerical exponents and smoothness cutoff, as recorded in \cref{obj:lem:published-cate-benchmark}. Their statistical guarantees concern their bounded-density model and stated attainment conditions; the lower bound for our known-uniform class combines the original-record construction in \cref{obj:thm:marked-component-certificate,obj:lem:nuisance-frontier-converse}, which covers smoothness below \(S_{\mathrm{crit}}\) with total count at most \(n^{q_*}\), with the labeled-sample floor \(n^{-\gamma/(2\gamma+d)}\) from the constant-nuisance subfamily in \cref{obj:thm:supplied-propensity}. Semi-supervised analyses of average and quantile treatment effects characterize nuisance-dependent expansions and inference for marginal targets \citep[Theorem 2.1, Corollary 2.1, Theorem 3.1, and Corollary 3.1]{ChakraborttyDai2022}. Our contribution characterizes worst-case pointwise absolute error and its acquisition requirements for the specified conditional target. The projection construction connects local residualized regression \citep[Theorem 3]{Kennedy2023} with finite-rank second-order correction \citep[Lemma 3]{Robins2009}.



The exposition proceeds through related work, setup and assumptions, the estimator construction, sharp precision and sample requirements, and discussion and limitations. The upper-bound argument is organized in \cref{sec:upper-bound-proofs}, the original-record converse in \cref{sec:original-record-lower-bounds}, and the acquisition algebra in \cref{sec:acquisition-consequences}.


\section{Related work}\label{sec:related-work}

The closest comparison is \citet[Section 3, Theorems 1--2 and equation (14)]{KennedyBalakrishnanRobinsWasserman2024}, who characterize supervised pointwise CATE estimation under H\"older smoothness. Write \(n\) for the labeled sample size, \leanref{sym:d}{\ensuremath{d}} for the covariate dimension, and \leanref{sym:alpha}{\ensuremath{\alpha}}, \leanref{sym:beta}{\ensuremath{\beta}}, and \leanref{sym:gamma}{\ensuremath{\gamma}} for the propensity, control-mean, and treatment-effect smoothness orders. With \(s_{\mathrm{pub}}=(\alpha+\beta)/2\), their absolute-error exponents are
\[
\frac{1}{1+d/(2\gamma)+d/(4s_{\mathrm{pub}})}
\quad\text{when}\quad
s_{\mathrm{pub}}<\frac{d/4}{1+d/(2\gamma)},
\qquad
\frac{\gamma}{2\gamma+d}
\quad\text{otherwise}.
\]
Published attainment requires their matrix regularity, smooth nuisance-estimation errors, boundedness conditions, and control of the third-order remainder involving covariate-distribution estimation. Their localized second-order projection correction balances effect approximation, nuisance approximation, and sampling variation.
% [Supervised pointwise minimax results](https://arxiv.org/pdf/2203.00837)

The algebraic comparison appears in \cref{obj:lem:published-cate-benchmark}. It identifies the supervised exponents obtained by specializing the sample-size expression to zero auxiliary records. Its interpretation follows the published source model, which permits a bounded covariate density and uses a construction-specific design in the lower-bound argument. The present \leanref{S-1}{potential-outcome law class with known uniform design}, specified in \cref{obj:def:primitive-class}, fixes the covariate distribution throughout. Within that class, \cref{obj:thm:marked-component-certificate,obj:lem:nuisance-frontier-converse} supply the original-record testing argument. These results establish the lower bound in the experiment actually observed by the decision maker, while the published benchmark supplies the numerical supervised comparison.

This distinction becomes consequential when information arrives in unequal amounts. The \leanref{S-2}{sampling experiment} combines \(n\) labeled records with \(m\) independent outcome-free covariate-treatment records; its sampling law is given in \cref{obj:ass:sample-law}. The pointwise absolute-error characterization in \cref{obj:thm:sharp-annotation-frontier} retains both record counts, and \cref{obj:thm:supplied-propensity} identifies the precision attainable when the propensity function is supplied. Together, these results connect the amount of treatment information to the accuracy of a conditional causal contrast under the paper's fixed design, overlap, and smoothness conditions.

Related semi-supervised treatment-effect analyses study population functionals through asymptotic expansions. \citet[Theorem 2.1 and Corollary 2.1]{ChakraborttyDai2022} separate labeled residual variation from averaging over the enlarged covariate sample for average treatment effect estimation. Their corollary addresses a vanishing labeled fraction with a correctly specified propensity model, retaining explicit nuisance remainders; asymptotic normality follows under the stated remainder conditions. The corresponding quantile treatment effect expansions additionally account for the density at the target quantile \citep[Theorem 3.1 and Corollary 3.1]{ChakraborttyDai2022}. These expansions explain robustness and efficiency for marginal causal functionals. The present risk characterization instead evaluates a conditional contrast at a fixed interior profile and quantifies worst-case absolute error over the specified H\"older class.
% [Semi-supervised ATE and QTE expansions](https://arxiv.org/pdf/2201.00468)

The information carried by auxiliary records also distinguishes the surrogate-assisted setting of \citet[Sections 2.1, 2.5--2.6, and 4]{ChengAnanthakrishnanCai2021}. Their records include outcome-predictive surrogate variables, potentially measured after treatment, alongside baseline covariates and treatment. An imputation-based ATE estimator exploits these variables, with robustness to imputation-model misspecification and local efficiency under their ideal semi-supervised model; perturbation resampling supports inference. Their electronic health records application illustrates the role of surrogate information when chart-reviewed outcomes are scarce. Here, auxiliary observations contribute covariate-treatment information, and their value is characterized through the pointwise risk and sample requirements in \cref{obj:thm:sharp-annotation-frontier,obj:prop:annotation-threshold}.
% [Surrogate-assisted ATE estimation](https://pmc.ncbi.nlm.nih.gov/articles/PMC7758040/)

A further comparison concerns the statistical criterion used to measure precision. \citet[Theorems 4.2 and 4.4 and Corollary 4.5]{Kato2026} derive two-sample ATE efficiency bounds for regular estimators and establish attainment through cross-fitted estimation under nuisance consistency and product-rate conditions. Their auxiliary sample contains covariates, and their analysis allows distinct labeled and auxiliary covariate distributions subject to support conditions. With coincident distributions and the stated evaluation weighting, the efficiency bound separates outcome residual variation from variation in the conditional effect across covariate profiles. Those conclusions characterize asymptotic variance at a fixed positive sampling fraction. The present minimax criterion, defined in \cref{obj:def:minimax}, measures pointwise absolute error across the primitive class and accommodates every auxiliary-record count in the stated experiment.
% [Regular-estimator ATE efficiency results](https://arxiv.org/html/2511.08303)

Supplied-predictor conditional estimation offers another route for using auxiliary information. \citet[Sections 2.4 and 3 and Theorems 1--3]{SuiZhouZhouDai2026} combine labeled observations, unlabeled covariates, and a supplied predictor to estimate conditional moment targets. Their Sobolev-equivalent reproducing kernel Hilbert space conditions support pointwise error bounds and confidence intervals, with prediction and residual components determining variance. Structure-adaptive CATE estimation uses different restrictions: \citet[Sections 2.3 and 6, Corollary 4]{Kim2026StructureAdaptiveCATE} give a pointwise squared-error result under bounded ambient-RKHS and overlap conditions when the contrast lies in the designated lower-complexity subspace; their broader analysis also treats integrated squared error, source conditions, known lower-dimensional representations, and model selection. Those results use supervised outcomes and RKHS/Sobolev structure. Our experiment instead fixes a H\"older class, measures pointwise absolute error, and studies the extra information supplied by outcome-free covariate-treatment records.
% [Supplied-predictor conditional inference](https://arxiv.org/html/2603.05575)
% [Structure-adaptive CATE estimation](https://arxiv.org/html/2602.18958)

Methodologically, the projection correction belongs to the higher-order influence-function tradition. Finite-rank orthogonal projections provide a tractable representation of second-order corrections in nonlinear functional estimation \citep[Lemma 3]{Robins2009}. Orthogonal regression supplies the conditional-effect estimating structure: residualized treatment and outcome regression underlie the R-learner \citep[Section 1]{Nie2021}, and local polynomial double-residual regression connects this structure to pointwise smoothness-based estimation \citep[Theorem 3]{Kennedy2023}. The \leanref{S-3}{rectangular projection estimator} in \cref{obj:def:moments,obj:def:estimator} combines localization with independent treatment and outcome roles. Its upper bound in \cref{obj:thm:rectangular-upper} makes the separate sampling contributions explicit, linking the projection architecture to the unequal information available from labeled and outcome-free records.
% [Finite-rank projection calculus](https://pmc.ncbi.nlm.nih.gov/articles/PMC3475538/)
% [Residualized conditional-effect regression](https://arxiv.org/pdf/1712.04912)
% [Local polynomial double-residual regression](https://arxiv.org/pdf/2004.14497)


\section{Setup and assumptions}\label{sec:setup}

Throughout, \leanref{sym:c}{\ensuremath{c}} and \leanref{sym:C}{\ensuremath{C}} denote finite positive constants whose values may change between appearances. At fixed public parameters, these constants are independent of the sample sizes. We write \(a\lesssim b\) when \(a\le Cb\), and \(a\asymp b\) when both \(a\lesssim b\) and \(b\lesssim a\). Vector norms are Euclidean unless otherwise indicated. We use \(\mathbb N=\mathbb N_0=\{0,1,\ldots\}\) and write \(\mathbb N_{\ge1}\) for the positive integers. Counts and polynomial multi-indices may be zero; dimensions, resolutions, and sample-index ranges carry their displayed positive lower bounds. The symbol \(N=n+m\) always denotes the total treatment-record count.

Fix the public parameters
\[
d\ge1,\qquad
\alpha,\beta\in(0,1],\qquad
\gamma\in[1,3],\qquad
L>1/2,\qquad
\varepsilon\in(0,1/2).
\]
Here \(d\) is the covariate dimension, \(\alpha,\beta,\gamma\) are the propensity, control-mean, and causal-contrast smoothness orders, respectively, \leanref{sym:L}{\ensuremath{L}} is the common Hölder radius, and \leanref{sym:eps}{\ensuremath{\varepsilon}} is the overlap level. These parameters remain fixed as the sample sizes vary.

\subsection{Primitive variables and identifying restrictions}

We use a coordinate-partial convention for Hölder regularity. For a function \leanref{sym:f}{\ensuremath{f}} and a regularity order \leanref{sym:s}{\ensuremath{s}}, the norm \leanref{sym:normH}{\ensuremath{\|f\|_{H^s}}} controls coordinate derivatives indexed by the nonnegative-integer multi-index \leanref{sym:kappa}{\ensuremath{\kappa}}, together with their highest-order Hölder variation. The following definition specifies the convention, including at integer smoothness orders.

\begin{definitionv}{def:holder-norm}[Coordinate-partial Hölder norm]
The extended-real Hölder norm \(\|f\|_{H^s}\), for \(f:\mathbb R^d\to\mathbb R\), \(d\in\{0,1,\ldots\}\), and \(s\in\mathbb R\), is
\[
\|f\|_{H^s}
=
\begin{cases}
\displaystyle
\max\left\{
\sup_{\substack{\kappa\in\{0,1,\ldots\}^d,\;
|\kappa|\le \max\{\lceil s\rceil-1,0\}\\
x\in[0,1]^d}}
|\partial^\kappa f(x)|,\;
\sup_{\substack{\kappa\in\{0,1,\ldots\}^d,\;
|\kappa|=\max\{\lceil s\rceil-1,0\}\\
x,x'\in[0,1]^d,\;x\ne x'}}
\frac{|\partial^\kappa f(x)-\partial^\kappa f(x')|}
{\|x-x'\|_2^{\,s-\max\{\lceil s\rceil-1,0\}}}
\right\},
& f\in C^{\max\{\lceil s\rceil-1,0\}}([0,1]^d),\\[1ex]
+\infty,
& \text{otherwise}.
\end{cases}
\]
Here \(|\kappa|=\sum_{i=1}^d\kappa_i\), and \(\partial^\kappa f(x)\) is the iterated derivative within \([0,1]^d\), evaluated on coordinate directions with coordinate \(i\) repeated \(\kappa_i\) times. The differentiability condition means continuous differentiability of the indicated order within the cube. Both suprema are taken in the extended nonnegative reals, with an empty supremum equal to zero.
\end{definitionv}

For positive \(s\), this convention uses derivatives through order \(\lceil s\rceil-1\) and a highest-order Hölder exponent of \(s-\lceil s\rceil+1\). Thus integer orders use a Lipschitz condition on derivatives of one lower order. Every regularity restriction below is evaluated within the unit cube.

A full record consists of a covariate vector \leanref{sym:X}{\ensuremath{X}}, a binary treatment indicator \leanref{sym:A}{\ensuremath{A}}, and binary potential responses \leanref{sym:Yzero}{\ensuremath{Y(0)}} and \leanref{sym:Yone}{\ensuremath{Y(1)}} under control and treatment. Let \leanref{sym:P}{\ensuremath{P=\mathcal L(X,A,Y(0),Y(1))}} denote their primitive Borel law. The observed response \leanref{sym:Y}{\ensuremath{Y}} selects the potential response corresponding to the realized treatment.

\begin{definitionv}{synth_5}[Observed response selection]
The observed response \(Y\) for a full record \((X,A,Y(0),Y(1))\), where \(X\in\mathbb R^d\) is the covariate vector for a nonnegative integer \(d\), \(A\in\{0,1\}\) is the treatment indicator, and \(Y(0),Y(1)\in\{0,1\}\) are the potential responses under control and treatment, is defined by
\[
Y=
\begin{cases}
Y(0), & A=0,\\
Y(1), & A=1.
\end{cases}
\]
The corresponding observed record is \((X,A,Y)\).
\end{definitionv}

The evaluation profile is the fixed interior point \leanref{sym:x0}{\ensuremath{x_0=(1/2,\ldots,1/2)}}. We write \leanref{sym:e}{\ensuremath{e_P}} for an admissible propensity witness, \leanref{sym:mu0}{\ensuremath{\mu_{0,P}}} and \leanref{sym:mu1}{\ensuremath{\mu_{1,P}}} for admissible control and treated conditional-mean witnesses, and \leanref{sym:tau}{\ensuremath{\tau_P}} for the canonical continuous causal contrast. Their precise admissibility requirements are collected in \cref{obj:def:primitive-class}. The identifying restrictions begin with the covariate design and treatment assignment.

\begin{assumptionv}{ass:uniform-design}[Uniform covariate design]
The covariate vector \(X\) satisfies
\[
\mathcal L_P(X)=\mathrm{Unif}([0,1]^d).
\]
\end{assumptionv}

The known uniform design fixes the distribution used for localization and integration and gives full support on the cube. It is the maintained uniform-design restriction of the bounded-density framework considered by \citep{KennedyBalakrishnanRobinsWasserman2024}. Treatment assignment is linked to the potential responses through conditional exchangeability.

\begin{assumptionv}{ass:exchangeability}[Conditional treatment exchangeability]
The treatment indicator \(A\) and the potential responses \(Y(0),Y(1)\) satisfy
\[
A\perp(Y(0),Y(1))\mid X,
\]
where \(X\) is the covariate vector.
\end{assumptionv}

Conditional exchangeability allows treatment-specific observed means to identify the corresponding potential-response means after conditioning on the covariates \citep{ChakraborttyDai2022}. Both treatment arms must also be represented throughout the evaluation domain.

\begin{assumptionv}{ass:overlap}[Pointwise overlap on the cube]
The overlap condition for a primitive law \(P\), at the overlap level \(\varepsilon\in\mathbb R\), is \(0<\varepsilon\) together with the requirement that the raw Borel propensity witness carried by \(P\) takes values in \([\varepsilon,1-\varepsilon]\) at every \(x\in[0,1]^d\).
\end{assumptionv}

Bounded overlap keeps both conditional treatment probabilities away from zero, as in the pointwise conditional-effect framework of \citep{KennedyBalakrishnanRobinsWasserman2024}. Here the endpoints are the public values \(\varepsilon\) and \(1-\varepsilon\), with \(0<\varepsilon<1/2\), and the restriction applies at every point of the cube.

The regularity conditions distinguish the two nuisance functions from the causal contrast. We first impose isotropic Hölder regularity on the propensity.

\begin{assumptionv}{ass:propensity-holder}[Hölder smoothness of the propensity]
The designated Borel propensity witness \(e_P\) satisfies
\[
\|e_P\|_{H^\alpha}\le L.
\]
\end{assumptionv}

Isotropic Hölder propensity regularity controls how treatment probabilities vary across nearby covariate profiles \citep{KennedyBalakrishnanRobinsWasserman2024}. The order \(\alpha\) is common across coordinates, and the radius \(L\) bounds both the function and its Hölder variation. The control response mean has its own smoothness order.

\begin{assumptionv}{ass:control-holder}[Hölder smoothness of the control mean]
The designated Borel control-mean witness \(\mu_{0,P}\) satisfies
\[
\|\mu_{0,P}\|_{H^\beta}\le L.
\]
\end{assumptionv}

Isotropic Hölder control-mean regularity governs variation in the baseline response surface \citep{KennedyBalakrishnanRobinsWasserman2024}. Allowing \(\beta\) to differ from \(\alpha\) separates outcome regularity from treatment-assignment regularity. The target surface is controlled directly by the following restriction.

\begin{assumptionv}{ass:effect-holder}[Hölder smoothness of the causal contrast]
The canonical continuous causal contrast \(\tau_P\) satisfies
\[
\|\tau_P\|_{H^\gamma}\le L.
\]
\end{assumptionv}

Isotropic Hölder contrast regularity controls the local variation of the conditional treatment effect itself \citep{KennedyBalakrishnanRobinsWasserman2024}. The separate order \(\gamma\in[1,3]\) permits a smoother causal contrast than either nuisance function and specifies the regularity used to estimate its value at \(x_0\).

The binary-response model also requires pointwise probability bounds on the designated arm means. We state these restrictions explicitly because admissible witnesses are functions defined throughout the cube.

\begin{assumptionv}{ass:control-interior}[Control response bounds]
The designated Borel control-mean witness \(\mu_{0,P}\) satisfies
\[
\mu_{0,P}(x)\in[0,1]
\qquad\text{for every }x\in[0,1]^d.
\]
\end{assumptionv}

This pointwise control-response restriction is specific to the present analysis and keeps the designated control mean within the Bernoulli probability range, including its endpoints. The corresponding treated mean obeys the same numerical bounds.

\begin{assumptionv}{ass:treated-interior}[Treated response bounds]
The designated Borel treated-mean witness \(\mu_{1,P}\) satisfies
\[
\mu_{1,P}(x)\in[0,1]
\qquad\text{for every }x\in[0,1]^d.
\]
\end{assumptionv}

The treated-response restriction is likewise specific to this analysis. Together, the two arm bounds ensure that the designated contrast takes values in \([-1,1]\) everywhere on the cube.

Membership in the primitive class is existential: a law belongs when it admits conditional-mean witnesses satisfying the stated restrictions. Let \leanref{sym:M}{\ensuremath{\mathcal M(d,\alpha,\beta,\gamma,L,\varepsilon)}} denote this class, abbreviated to \(\mathcal M\) at the fixed public parameters.

\begin{definitionv}{def:primitive-class}[Primitive law class \(\mathcal M(d,\alpha,\beta,\gamma,L,\varepsilon)\)]
\(\mathcal M(d,\alpha,\beta,\gamma,L,\varepsilon)\) is the class of primitive Borel laws \(P=\mathcal L(X,A,Y(0),Y(1))\), where \(X\) is the covariate vector, \(A\) is the binary treatment indicator, and \(Y(0)\) and \(Y(1)\) are the binary potential responses under control and treatment. For public parameters \(L>1/2\) and \(\varepsilon\in(0,1/2)\), membership is defined by the following properties:
\begin{itemize}
\item \textbf{(Uniform design.)} The law satisfies \cref{obj:ass:uniform-design}.
\item \textbf{(Exchangeability.)} The law satisfies \cref{obj:ass:exchangeability}.
\item \textbf{(Conditional means.)} There exist Borel functions
\[
g_e,g_0,g_1:[0,1]^d\to[0,1]
\]
such that
\[
\begin{aligned}
g_e(X)&=\mathbb E_P[A\mid X] &&P\text{-almost surely},\\
g_0(X)&=\mathbb E_P[Y(0)\mid X] &&P\text{-almost surely},\\
g_1(X)&=\mathbb E_P[Y(1)\mid X] &&P\text{-almost surely}.
\end{aligned}
\]
\item \textbf{(Overlap and smoothness.)} These functions satisfy \cref{obj:ass:overlap,obj:ass:propensity-holder,obj:ass:control-holder,obj:ass:effect-holder}, with \(g_e\), \(g_0\), and \(g_1-g_0\) as the propensity, control mean, and causal contrast.
\item \textbf{(Response bounds.)} The arm means satisfy \cref{obj:ass:control-interior,obj:ass:treated-interior}.
\end{itemize}
For each admitted law, \(e_P,\mu_{0,P},\mu_{1,P}\) designate admissible functions \(g_e,g_0,g_1\), respectively. The function \(\tau_P\) is the canonical continuous representative of the almost-sure conditional causal contrast, and the designated functions satisfy
\[
\mu_{1,P}(x)-\mu_{0,P}(x)=\tau_P(x)
\qquad\text{for every }x\in[0,1]^d.
\]
\end{definitionv}

The witness formulation separates almost-sure conditional-expectation identities from restrictions imposed at every covariate profile. Under the positive smoothness orders fixed above, \(e_P\), \(\mu_{0,P}\), and \(\tau_P\) are continuous; the identity in \cref{obj:def:primitive-class} then gives continuity of \(\mu_{1,P}\) as well.

Observed response selection, conditional exchangeability, and overlap yield
\[
\mathbb E_P[Y\mid A=a,X]=\mu_{a,P}(X)
\qquad P\text{-almost surely},\qquad a\in\{0,1\}.
\]
Consequently, the observed-record law identifies the causal contrast almost everywhere under the covariate distribution. Uniform design assigns positive probability to every nonempty relatively open subset of the cube. Continuous representatives that agree almost everywhere therefore agree everywhere, which makes \(\tau_P(x_0)\) a uniquely defined, identified point target. The continuous representative gives substantive meaning to evaluation at a fixed profile even though conditional expectations are initially defined through almost-sure identities.

\subsection{Sampling experiment and risk}

Write \leanref{sym:Po}{\ensuremath{P^o=\mathcal L_P(X,A,Y)}} for the observed-record law induced by \cref{obj:synth_5}, and \leanref{sym:PXA}{\ensuremath{P_{XA}=\mathcal L_P(X,A)}} for the covariate-treatment marginal. For \(n\ge2\) and \(m\ge0\), the original dataset \leanref{sym:D}{\ensuremath{\mathcal D_{n,m}}} consists of \(n\) independent labeled records with law \(P^o\) and \(m\) independent outcome-free records with law \(P_{XA}\):
\[
\mathcal D_{n,m}
=
\left(
((X_i,A_i,Y_i))_{i=1}^{n},
((\widetilde X_j,\widetilde A_j))_{j=1}^{m}
\right).
\]
The two blocks are independent. A uniform randomizer \leanref{sym:U}{\ensuremath{U}} on \([0,1]\), independent of both blocks, is available to decisions. Thus the product sampling law \leanref{sym:E}{\ensuremath{\mathsf E_{n,m}(P)}} in \cref{obj:ass:sample-law} is
\[
\mathsf E_{n,m}(P)
=
(P^o)^{\otimes n}
\otimes P_{XA}^{\otimes m}
\otimes\mathrm{Unif}([0,1]).
\]
The following condition makes this sampling experiment explicit.

\begin{assumptionv}{ass:sample-law}[Original experiment sampling law]
The original dataset \(\mathcal D_{n,m}\) and independent uniform randomizer \(U\) have joint law
\[
\mathcal L_P(\mathcal D_{n,m},U)=\mathsf E_{n,m}(P),
\]
where \(\mathsf E_{n,m}(P)\) is the product sampling law of the original dataset and randomizer.
\end{assumptionv}

Independent fixed-size labeled and outcome-free samples from a shared law provide the two-channel sampling structure used in semi-supervised causal estimation \citep{ChakraborttyDai2022}; the independent randomizer permits randomized decisions. In this experiment, each labeled record supplies covariates, treatment, and a response, while each outcome-free record supplies covariates and treatment. The common marginal \(P_{XA}\) makes both blocks informative about treatment assignment.

A decision \leanref{sym:T}{\ensuremath{T}} is a Borel real-valued function of the dataset and randomizer. Its absolute-error risk \leanref{sym:risk}{\ensuremath{\mathcal R_{n,m}(T,P)}} measures expected error in estimating the identified point contrast.

\begin{definitionv}{def:risk}[Absolute-error risk $\mathcal R_{n,m}(T,P)$]
The absolute-error risk is
\[
\mathcal R_{n,m}(T,P)
=
\int
\left|T(\mathcal D_{n,m},U)-\tau_P(x_0)\right|
\,d\mathsf E_{n,m}(P).
\]
Here \(T\) is a Borel real-valued decision using the original dataset \(\mathcal D_{n,m}\) and independent uniform randomizer \(U\), \(x_0\) is the fixed interior covariate profile, and \(\mathsf E_{n,m}(P)\) is the product sampling law of the original dataset and randomizer. The target \(\tau_P(x_0)\) is the canonical conditional causal contrast from \cref{obj:def:primitive-class}.
\end{definitionv}

Absolute loss expresses precision on the scale of the conditional response difference. To assess precision uniformly across admissible primitive laws, the minimax risk \leanref{sym:R}{\ensuremath{R(n,m)}} optimizes the worst-case value of this loss over all Borel decisions.

\begin{definitionv}{def:minimax}[Minimax risk \(R(n,m)\)]
\[
R(n,m)=\inf_{T\ {\rm Borel}}\sup_{P\in\mathcal M}\mathcal R_{n,m}(T,P).
\]
Here \(\mathcal M\) is the primitive law class at the fixed public parameters, and \(\mathcal R_{n,m}(T,P)\) is the expected absolute error for the point contrast under law \(P\). The infimum ranges over Borel real-valued decisions \(T\) using the original dataset of \(n\) labeled records and \(m\) independent outcome-free treatment records, together with an independent uniform randomizer.
\end{definitionv}

The sample sizes enter separately because response information is supplied by the labeled block and treatment-assignment information is supplied by both blocks. A complementary experiment supplies the entire admissible propensity function to the decision. Denote such a decision by \leanref{sym:Te}{\ensuremath{T^e}}, and its minimax risk by \leanref{sym:Re}{\ensuremath{R^e(n,m)}}.

\begin{definitionv}{def:oracle-risk}[Supplied propensity minimax risk]
The supplied propensity minimax risk \(R^e(n,m)\), for \(d,n,m\in\mathbb N\) and \(\alpha,\beta,\gamma,L,\varepsilon\in\mathbb R\), is the extended nonnegative quantity
\[
R^e(n,m)=\inf_{T^e}\sup_{P\in\mathcal M}
\int
\left|T^e(\mathcal D_{n,m},U,e_P)-\tau_P(x_0)\right|
\,d\mathsf E_{n,m}(P).
\]
Here \(\mathcal M\) is the primitive law class at these parameters, and \(e_P\) and \(\tau_P\) are its designated admissible propensity function and canonical causal contrast, as defined in \cref{obj:def:primitive-class}. The evaluation profile is \(x_0=(1/2,\ldots,1/2)\). The original dataset \(\mathcal D_{n,m}\) comprises \(n\) labeled records and \(m\) independent outcome-free treatment records; \(U\) is the independent uniform randomizer on \([0,1]\), and \(\mathsf E_{n,m}(P)\) is their product sampling law from \cref{obj:ass:sample-law}. The infimum ranges over real-valued decisions \(T^e\) receiving the dataset, randomizer, and a supplied function from the covariate space to \(\mathbb R\), with Borel measurability in \((\mathcal D_{n,m},U)\) for every fixed supplied function.
\end{definitionv}

Supplying \(e_P\) gives a reference experiment in which treatment probabilities are available throughout the cube. Measurability is required in the observed data and randomizer for each fixed supplied function. Comparing \cref{obj:def:minimax,obj:def:oracle-risk} quantifies how the original sampling experiment approaches the precision available with supplied treatment probabilities.

\subsection{Rate notation}

We introduce the rate notation as a single family. Let \(N=n+m\) be the total treatment-record count and \(S=\alpha+\beta\) the sum of the two nuisance smoothness orders. Define the tuning denominator \leanref{sym:Delta}{\ensuremath{\Delta}}, nuisance-smoothness cutoff \leanref{sym:Scrit}{\ensuremath{S_{\mathrm{crit}}}}, and total-record comparison exponent \leanref{sym:qstar}{\ensuremath{q_*}} by
\[
\Delta=2\gamma+d+\frac{\gamma d}{S},
\qquad
S_{\mathrm{crit}}=\frac{\gamma d}{2\gamma+d},
\qquad
q_*=\frac{\gamma d}{S(2\gamma+d)}
=\frac{S_{\mathrm{crit}}}{S}.
\]
The supplied-propensity reference order \leanref{sym:ro}{\ensuremath{r_o(n)}} and the unequal-channel upper-rate expression \leanref{sym:rup}{\ensuremath{r_{\mathrm{up}}(n,m)}} are
\[
r_o(n)=n^{-\gamma/(2\gamma+d)},
\qquad
r_{\mathrm{up}}(n,m)
=
\max\left\{
r_o(n),\,
(nN)^{-\gamma/\Delta}
\right\}.
\]
These expressions distinguish a term indexed by the response-bearing sample size from a term indexed by the product of the labeled and total treatment-record counts. Their statistical interpretation is supplied by \cref{obj:thm:rectangular-upper,obj:thm:sharp-annotation-frontier,obj:thm:supplied-propensity}.

The notation also makes the sample-size comparison transparent. Algebraically,
\[
(nN)^{-\gamma/\Delta}\le r_o(n)
\quad\Longleftrightarrow\quad
N\ge n^{q_*}.
\]
When \(S\ge S_{\mathrm{crit}}\), the exponent satisfies \(q_*\le1\), so \(N\ge n\) meets this comparison. For \(S<S_{\mathrm{crit}}\), one has \(q_*>1\), and the comparison requires total treatment-record growth beyond the labeled count. This division organizes the precision and sample-requirement statements that follow.


\section{Estimation with outcome-free records}\label{sec:estimation}

The estimator allocates labeled and outcome-free records to two independent statistical roles. One role supplies response information; the other supplies treatment-assignment information. A local polynomial fit estimates the contrast at the target profile \leanref{sym:x0}{\ensuremath{x_0}}, while a finer projection corrects its empirical moments for the unknown propensity and control mean. Separate localization and projection scales allow the construction to use the smoothness of the effect and the combined smoothness of the nuisance functions.

We first specify the neighborhood associated with the localization scale \leanref{sym:h}{\ensuremath{h}} and its probability measure.

\begin{definitionv}{synth_3}[Localization cube and uniform localization measure]
For an integer \(d\ge1\), let \(x_0=(1/2,\ldots,1/2)\in[0,1]^d\). For \(0<h\le1/2\), define the localization cube of side length \(h\), centered at \(x_0\), by
\[
C_h=x_0+[-h/2,h/2]^d.
\]
The uniform probability measure \(\nu_h\) on \(C_h\) has density \(w_h\) with respect to \(d\)-dimensional Lebesgue measure \(dx\) on \([0,1]^d\):
\[
w_h(x)=h^{-d}\mathbf1\{x\in C_h\},
\qquad
d\nu_h(x)=w_h(x)\,dx.
\]
Equivalently, for every Borel set \(E\subseteq[0,1]^d\),
\[
\nu_h(E)=h^{-d}\int_{E\cap C_h}dx.
\]
\end{definitionv}

The localization cube \leanref{sym:Ch}{\ensuremath{C_h}} lies inside the covariate domain. Its uniform probability measure \leanref{sym:nuh}{\ensuremath{\nu_h}} defines the inner product for both polynomial spaces below. Under \cref{obj:ass:uniform-design}, weighting an observed covariate by the localization density \leanref{sym:wh}{\ensuremath{w_h}} converts a population average into an integral under this local probability measure.

The coarse basis \leanref{sym:r}{\ensuremath{r}} consists of scaled tensor-product Legendre polynomials of total degree at most two. This degree accommodates local polynomial approximation throughout the stipulated effect-smoothness range.

\begin{definitionv}{synth_1}[Scaled coarse polynomial basis]
The coarse polynomial basis vector \(r\), for dimension \(d\in\mathbb N\) and scale \(h\in\mathbb R\), is defined at each \(x\in\mathbb R^d\) by
\[
r(x)=
\left(
\prod_{i=1}^{d}
\begin{cases}
1, & \kappa_i=0,\\
\sqrt{12}\,\dfrac{x_i-\tfrac12}{h}, & \kappa_i=1,\\
\sqrt{5}\left(6\left(\dfrac{x_i-\tfrac12}{h}\right)^2-\tfrac12\right),
& \kappa_i=2
\end{cases}
\right)_{\kappa\in\mathbb N_0^d:\,\sum_{i=1}^{d}\kappa_i\le 2}.
\]
Division by zero is interpreted as zero, so this definition also specifies \(r\) at \(h=0\).
\end{definitionv}

\begin{definitionv}{synth_7}[Localized cellwise polynomial basis]
For \(d\ge1\), \(0<h\le1/2\), and an integer \(J\ge1\), let \(x_0=(1/2,\ldots,1/2)\), \(C_h=x_0+[-h/2,h/2]^d\), and \(\delta=h/J\). Partition each coordinate interval of \(C_h\) into \(J\) intervals of length \(\delta\), taking each interval half-open on the right except for the last interval, which includes the outer right endpoint. Their Cartesian products form the partition \(\{C_{\mathbf j}:\mathbf j\in\{0,\ldots,J-1\}^d\}\). The center of \(C_{\mathbf j}\) is
\[
c_{\mathbf j}=x_0-\frac h2\mathbf1+\delta\left(\mathbf j+\frac12\mathbf1\right).
\]
Write
\[
P_0(u)=1,\qquad P_1(u)=u,\qquad P_2(u)=\frac{3u^2-1}{2}.
\]
For each cell index \(\mathbf j\) and multi-index \(\kappa\in\mathbb N_0^d\) with \(|\kappa|\le2\), define
\[
b_{\mathbf j,\kappa}(x)
=
J^{d/2}\mathbf1\{x\in C_{\mathbf j}\}
\prod_{v=1}^d
\sqrt{2\kappa_v+1}\,
P_{\kappa_v}\!\left(\frac{2(x_v-c_{\mathbf j,v})}{\delta}\right).
\]
The vector \(\mathbf b_J(x)\) collects these functions, ordered lexicographically by cell index and then by multi-index. Its dimension is
\[
k=J^d\binom{d+2}{2}.
\]
These functions form the positive-leading-coefficient cellwise scaled Legendre orthonormal basis, under \(d\nu_h(x)=h^{-d}\mathbf1\{x\in C_h\}\,dx\), of the space \(V_J\) of polynomials of total degree at most two on each partition cell. Set \(\mathbf b_J(x)=0\) for \(x\notin C_h\). Thus the associated kernel is uniquely specified by
\[
K_J(x,x')=\mathbf b_J(x)^\top\mathbf b_J(x').
\]
\end{definitionv}



The vector \(r\) is the coarse orthonormal basis and has \(q=\binom{d+2}{2}\) coordinates. Write \leanref{sym:r0}{\ensuremath{r_0=r(x_0)}} for its value at the target profile. A coefficient vector \leanref{sym:v}{\ensuremath{v\in\mathbb R^q}} represents the coarse polynomial \leanref{sym:pv}{\ensuremath{p_v(x)=r(x)^\top v}}; evaluating the fitted contrast therefore amounts to evaluating its coefficient vector against \(r_0\). The prescribed fine basis is \leanref{sym:bvec}{\ensuremath{\mathbf b_J}} from the preceding definition. It has \(k=qJ^d\) coordinates and spans the already specified cellwise space \(V_J\); the associated projection is defined next.

\begin{definitionv}{def:projection}[Projection \(\Pi_J\) and kernel \(K_J\)]
\[
\delta=\frac hJ,\qquad
K_J(x,x')=\mathbf b_J(x)^\top \mathbf b_J(x'),\qquad
\Pi_Jf(x)=\int K_J(x,x')f(x')\,d\nu_h(x').
\]
Here \(\delta\) is the fine-cell side length, \(h\) is the localization scale, \(J\) is the grid resolution, \(\mathbf b_J\) is the cellwise orthonormal fine polynomial basis vector, and \(\nu_h\) is the uniform probability measure on the localization cube \(C_h\). Weighted basis expressions are zero outside \(C_h\).
\end{definitionv}

The projection kernel \leanref{sym:K}{\ensuremath{K_J}} represents the orthogonal projection \leanref{sym:Proj}{\ensuremath{\Pi_J}} onto \(V_J\). The coarse fit operates across the entire localization cube, whereas the projection resolves nuisance variation on cells of side length \(\delta\). Refining the grid from \(J\) to a multiple \(J'\) of \(J\) subdivides every cell, so \(V_J\subseteq V_{J'}\) and the correction space is enriched; at every resolution \(J\), \(V_J\) contains the coarse polynomials, so the projection reproduces them.

We next assign the original observations to independent roles under the sampling law in \cref{obj:ass:sample-law}.

The concatenated treatment sequence \leanref{sym:DT}{\ensuremath{\mathcal D^T}} retains the covariate coordinate \leanref{sym:XT}{\ensuremath{X_i^T}} and treatment coordinate \leanref{sym:AT}{\ensuremath{A_i^T}} of each record. The outcome-role count \leanref{sym:ell}{\ensuremath{\ell}} is proportional to the labeled sample size, while the treatment-role count \leanref{sym:t}{\ensuremath{t}} is proportional to the total treatment-record count. Their index sets, \leanref{sym:IL}{\ensuremath{I_L}} and \leanref{sym:IT}{\ensuremath{I_T}}, select disjoint records. Thus additional outcome-free observations enlarge the treatment role while the response information remains tied to the labeled sample. The split also makes products of empirical averages across roles suitable for estimating products of population moments.

To interpret the response correction, write
\[
\eta_P(x)=\mathbb E_P[Y\mid X=x]
=\mu_{0,P}(x)+e_P(x)\tau_P(x),
\]
using the designated witnesses from \cref{obj:def:primitive-class}. At the population level, residualizing treatment against its propensity gives
\[
\mathbb E_P[(A-e_P(X))Y\mid X=x]
=e_P(x)\{1-e_P(x)\}\tau_P(x).
\]
The empirical construction implements this local moment relation through a finite projection and independent roles. Its outcome projection uses the observed response, and its treatment role supplies the assignment factor in the projection correction.

\begin{algorithmv}{def:roles}[Record roles \(\mathcal D^T,I_L,I_T\)]
Given the original dataset \(\mathcal D_{n,m}\) of \(n\) labeled records and \(m\) independent outcome-free treatment records, form the two roles as follows.
\begin{enumerate}
\item Concatenate the covariate-treatment pairs, with total treatment-record count \(N\):
\[
N=n+m,\qquad
\mathcal D^T=((X_i^T,A_i^T))_{i=1}^N,
\qquad
(X_i^T,A_i^T)=
\begin{cases}
(X_i,A_i),&1\le i\le n,\\
(\widetilde X_{i-n},\widetilde A_{i-n}),&n<i\le N.
\end{cases}
\]
Here \(X\) denotes the covariate vector and \(A\) the binary treatment indicator; tildes denote auxiliary records.
\item Define the outcome-role count \(\ell\), treatment-role count \(t\), and their index sets:
\[
\ell=\lfloor n/2\rfloor,\qquad
t=N-\ell,\qquad
I_L=\{1,\ldots,\ell\},\qquad
I_T=\{\ell+1,\ldots,N\}.
\]
\item Output the outcome-role labeled records and treatment-role covariate-treatment pairs:
\[
\bigl((X_j,A_j,Y_j)\bigr)_{j\in I_L},
\qquad
\bigl((X_i^T,A_i^T)\bigr)_{i\in I_T}.
\]
Here \(Y_j\) is the observed response. Outcomes from treatment-role records are discarded.
\end{enumerate}
\end{algorithmv}

The empirical outcome projection \leanref{sym:etahat}{\ensuremath{\widehat\eta_J^L}} estimates \(\Pi_J\eta_P\). In the response vector \leanref{sym:Rhat}{\ensuremath{\widehat R_{h,J}}}, the first average estimates the treated-response moment and the rectangular average supplies its projected treatment-response correction. The raw matrix \leanref{sym:Qraw}{\ensuremath{\widehat Q^{\mathrm{raw}}_{h,J}}} applies the corresponding correction to the local polynomial design. Symmetrization produces the matrix \leanref{sym:Qhat}{\ensuremath{\widehat Q_{h,J}}} used for inversion.

Define the population matrix \leanref{sym:Q}{\ensuremath{Q_{P,h,J}}} and population response vector \(\bar R\) by
\[
Q_{P,h,J}=\mathbb E_P[\widehat Q_{h,J}],
\qquad
\bar R=\mathbb E_P[\widehat R_{h,J}],
\]
where expectations are taken under the original sampling experiment. These are the population counterparts of the projected local estimating equations. Polynomial reproduction links those equations to approximation of the contrast, and overlap supplies the population positivity used to control inversion; the respective statements are \cref{obj:lem:local-polynomial-projection-facts,obj:lem:population-positivity}.

For the sample decision, an eigenvalue guard supplies a criterion for inversion at the public overlap level. We use the clipping operator
\[
\mathrm{clip}_{[-1,1]}(z)
=\max\{-1,\min\{1,z\}\},
\]
which projects a real number onto the contrast range implied by the probability-valued arm means in \cref{obj:def:primitive-class}. The guard level is
\[
g(\varepsilon)=\frac{\varepsilon(1-\varepsilon)}{2},
\]
half the population Gram margin supplied by overlap. The complete decision combines the empirical moments, guard, and clipping.

\begin{algorithmv}{def:moments}[Empirical moments \(\widehat R_{h,J},\widehat Q_{h,J}\)]
Given the original dataset \(\mathcal D_{n,m}\), localization scale \(h\), grid resolution \(J\), localization cube \(C_h\), localization density \(w_h\), coarse polynomial basis vector \(r\), and fine projection kernel \(K_J\), compute the moments as follows.
\begin{enumerate}
\item Form the independent outcome and treatment roles, with counts and index sets
\[
N=n+m,\qquad
\ell=\lfloor n/2\rfloor,\qquad
t=N-\ell,\qquad
I_L=\{1,\ldots,\ell\},\qquad
I_T=\{\ell+1,\ldots,N\}.
\]
The treatment sequence consists of the covariate-treatment pairs
\[
(X_i^T,A_i^T)=
\begin{cases}
(X_i,A_i),&1\le i\le n,\\
(\widetilde X_{i-n},\widetilde A_{i-n}),&n<i\le N.
\end{cases}
\]
Use labeled records \((X_j,A_j,Y_j)\) for \(j\in I_L\) in the outcome role and pairs \((X_i^T,A_i^T)\) for \(i\in I_T\) in the treatment role. Treatment-role outcomes are discarded. Weighted basis expressions are zero outside \(C_h\).
\item Define the outcome-role empirical projection:
\[
\widehat\eta_J^L(x)=
\begin{cases}
\displaystyle\frac1\ell\sum_{j\in I_L}
w_h(X_j)K_J(x,X_j)Y_j,&x\in C_h,\\
0,&x\notin C_h.
\end{cases}
\]
\item Define the empirical response vector:
\[
\begin{aligned}
\widehat R_{h,J}
&=\frac1\ell\sum_{j\in I_L}w_h(X_j)r(X_j)A_jY_j
-\frac1t\sum_{i\in I_T}w_h(X_i^T)r(X_i^T)A_i^T
\widehat\eta_J^L(X_i^T)\\
&=\frac1\ell\sum_{j\in I_L}w_h(X_j)r(X_j)A_jY_j
-\frac1{t\ell}\sum_{i\in I_T}\sum_{j\in I_L}
w_h(X_i^T)w_h(X_j)r(X_i^T)A_i^T
K_J(X_i^T,X_j)Y_j.
\end{aligned}
\]
The response factor in the rectangular correction is \(Y_j\), and the first summand uses \(A_jY_j\). The nested and rectangular expressions define the same vector.
\item Define the unsymmetrized empirical matrix:
\[
\begin{aligned}
\widehat Q^{\mathrm{raw}}_{h,J}
&=\frac1\ell\sum_{j\in I_L}w_h(X_j)r(X_j)r(X_j)^\top A_j\\
&\quad-\frac1{t\ell}\sum_{i\in I_T}\sum_{j\in I_L}
w_h(X_i^T)w_h(X_j)r(X_i^T)r(X_j)^\top
A_i^TA_jK_J(X_i^T,X_j).
\end{aligned}
\]
\item Output the response vector and symmetrized matrix:
\[
\left(\widehat R_{h,J},\widehat Q_{h,J}\right),
\qquad
\widehat Q_{h,J}
=\frac{\widehat Q^{\mathrm{raw}}_{h,J}
+(\widehat Q^{\mathrm{raw}}_{h,J})^\top}{2}.
\]
All inputs are observed in \(\mathcal D_{n,m}\), including when \(m=0\).
\end{enumerate}
\end{algorithmv}

The guarded decision \leanref{sym:That}{\ensuremath{\widehat T_{h,J}}} evaluates the fitted polynomial when the empirical matrix meets the eigenvalue cutoff and returns zero otherwise. Clipping bounds the decision on both branches. Together, these operations permit a risk guarantee over every admissible localization scale and grid resolution, including choices with sparse local observations.

The public choices of localization scale \leanref{sym:hstar}{\ensuremath{h_*}} and grid resolution \leanref{sym:Jstar}{\ensuremath{J_*}} balance approximation and sampling error. They determine the tuned decision \leanref{sym:Tup}{\ensuremath{\widehat T_{\mathrm{up}}}} using the sample sizes and the public class parameters.

\begin{definitionv}{def:estimator}[Guarded local polynomial decision \(\widehat T_{h,J}\)]
The local polynomial decision \(\widehat T_{h,J}\) is defined by
\[
\widehat T_{h,J}=
\begin{cases}
\mathrm{clip}_{[-1,1]}
\bigl(r_0^\top\widehat Q_{h,J}^{-1}\widehat R_{h,J}\bigr),
&\lambda_{\min}(\widehat Q_{h,J})\ge g(\varepsilon),\\
0,
&\lambda_{\min}(\widehat Q_{h,J})<g(\varepsilon).
\end{cases}
\]
Here \(r_0\) is the coarse polynomial basis vector evaluated at the target profile, \(\widehat Q_{h,J}\) is the symmetrized corrected empirical matrix, and \(\widehat R_{h,J}\) is the corrected empirical response vector. The operator \(\mathrm{clip}_{[-1,1]}\) projects a real value onto \([-1,1]\), and \(\lambda_{\min}\) denotes the smallest eigenvalue. For the public overlap level \(\varepsilon\), the inversion threshold is
\[
g(\varepsilon)=\frac{\varepsilon(1-\varepsilon)}{2}.
\]
\end{definitionv}

When \(S\ge\gamma\), the grid choice is \(J_*=1\). When \(S<\gamma\), the grid becomes finer as the localization scale shrinks, making \(\delta^S\) comparable in order to \(h^\gamma\). Thus nuisance approximation can use a finer spatial scale than the local effect fit. The localization choice then balances these approximation terms against the two sampling terms stated below.

For use in the decision-theoretic comparison, the same estimator is recorded as the original-record decision \leanref{sym:Tstar}{\ensuremath{T_*}}.

\begin{definitionv}{def:tuning}[Public tuning \(h_*,J_*,\widehat T_{\mathrm{up}}\)]
The localization scale \(h_*\), grid resolution \(J_*\), and tuned decision \(\widehat T_{\mathrm{up}}\) are
\[
h_*=\frac12\max\left\{n^{-1/(2\gamma+d)},(nN)^{-1/\Delta}\right\},
\qquad
J_*=\left\lceil h_*^{-\max\{0,\gamma/S-1\}}\right\rceil,
\qquad
\widehat T_{\mathrm{up}}=\widehat T_{h_*,J_*},
\]
where \(N\) is the total treatment-record count, \(S\) is the sum of the propensity and control-mean Hölder orders, and \(\Delta\) is the denominator governing the unequal-channel rate exponent.
\end{definitionv}

\begin{definitionv}{def:sharp-decision}[Sharp-rate decision \(T_*\)]
\[
T_*(\mathcal D_{n,m},U)
=\widehat T_{h_*,J_*}(\mathcal D_{n,m})
=\widehat T_{\mathrm{up}}(\mathcal D_{n,m}).
\]
For every original dataset \(\mathcal D_{n,m}\) and independent uniform randomizer \(U\), this decision uses the prescribed empirical moments, empirical eigenvalue cutoff, clipping, role split, and public tuning parameters \(h_*\) and \(J_*\). Its value is independent of \(U\).
\end{definitionv}

The following guarantee gives the construction's four statistical terms once: effect approximation \(h^\gamma\), nuisance approximation \(\delta^{\alpha+\beta}\), localized response sampling \((nh^d)^{-1/2}\), and rectangular sampling \((nNh^d\delta^d)^{-1/2}\). The public tuning balances them. The approximation, positivity, and sampling ingredients are isolated in \cref{obj:lem:local-polynomial-projection-facts,obj:lem:population-positivity,obj:lem:rectangular-sampling-bound}; their proofs are organized in \cref{sec:upper-bound-proofs}.

\begin{theoremv}{thm:rectangular-upper}[Rectangular estimation upper bounds]
Fix public parameters satisfying
\begin{itemize}
\item \textbf{(Dimension.)} \(d\) is an integer with \(d\ge1\).
\item \textbf{(Smoothness.)} \(0<\alpha\le1\), \(0<\beta\le1\), and \(1\le\gamma\le3\).
\item \textbf{(Radius.)} \(L>1/2\).
\item \textbf{(Overlap.)} \(0<\varepsilon<1/2\).
\end{itemize}
Let \(\mathcal M\) be the primitive class at these fixed parameters from \cref{obj:def:primitive-class}, and use the absolute-error risk from \cref{obj:def:risk}. There exists a finite constant \(C>0\), depending only on these public parameters, with both of the following guarantees.

For every pair of integers \(n\ge2\) and \(m\ge0\), every \(h\in(0,1/2]\), and every integer \(J\ge1\), the estimator \(\widehat T_{h,J}\) from \cref{obj:def:estimator}, viewed as a decision on the original dataset and randomizer that ignores the randomizer, is Borel measurable and satisfies, for every \(P\in\mathcal M\),
\[
\mathcal R_{n,m}(\widehat T_{h,J},P)
\le
C\min\left\{
1,\,
h^\gamma+
(h/J)^{\alpha+\beta}+
(nh^d)^{-1/2}+
\bigl(n(n+m)h^d(h/J)^d\bigr)^{-1/2}
\right\}.
\]

Define the supplied-propensity point-estimation order and the denominator governing the unequal-channel rate exponent by
\[
r_o(n)=n^{-\gamma/(2\gamma+d)},
\qquad
\Delta=2\gamma+d+\frac{\gamma d}{\alpha+\beta},
\]
and define the upper rate by
\[
r_{\mathrm{up}}(n,m)
=
\max\left\{
r_o(n),\,
\bigl(n(n+m)\bigr)^{-\gamma/\Delta}
\right\}.
\]
For every pair of integers \(n\ge2\) and \(m\ge0\), the publicly tuned decision \(\widehat T_{\mathrm{up}}\) from \cref{obj:def:tuning}, likewise viewed as a decision that ignores the randomizer, is Borel measurable and satisfies, for every \(P\in\mathcal M\),
\[
\mathcal R_{n,m}(\widehat T_{\mathrm{up}},P)
\le C r_{\mathrm{up}}(n,m).
\]
\end{theoremv}

This theorem completes the main-body estimator construction. The approximation, sampling, and inversion arguments are organized in \cref{sec:upper-bound-proofs}.


\section{Sharp precision and sample requirements}\label{sec:sharp-precision}

The sharp precision bound separates the information supplied by response labels from that supplied by all covariate-treatment records. The rate and attaining decision are defined once in \cref{obj:def:sharp-rate,obj:def:sharp-decision}. The following result gives their uniform minimax comparison, explicit sample-size regimes, and extension to randomized decisions.

\begin{definitionv}{def:sharp-rate}[Sharp rate \(r_*\)]
The sharp rate at the public class parameters is
\[
r_*(n,m;d,\alpha,\beta,\gamma)
=
\max\left\{
n^{-\gamma/(2\gamma+d)},
[n(n+m)]^{-\gamma/(2\gamma+d+\gamma d/(\alpha+\beta))}
\right\}
=
r_{\mathrm{up}}(n,m),
\]
where \(r_{\mathrm{up}}(n,m)\) is the maximum of the supplied-propensity and unequal-channel rate orders.
\end{definitionv}

\begin{theoremv}{thm:sharp-annotation-frontier}[Sharp annotation frontier]
Fix the covariate dimension \(d\in\mathbb N\), the propensity and control-mean smoothness orders \(\alpha,\beta\), the effect smoothness order \(\gamma\), the Hölder radius \(L\), and the overlap level \(\varepsilon\), satisfying:
\begin{itemize}
\item \textbf{(Dimension.)} \(d\ge1\).
\item \textbf{(Nuisance smoothness.)} \(0<\alpha\le1\) and \(0<\beta\le1\).
\item \textbf{(Effect smoothness.)} \(1\le\gamma\le3\).
\item \textbf{(Radius.)} \(L>1/2\).
\item \textbf{(Overlap.)} \(0<\varepsilon<1/2\).
\end{itemize}
Let \(\mathcal M\) denote the primitive class at these fixed parameters, and let \(r_*\), \(T_*\), and \(R(n,m)\) be the sharp rate, total data-only decision, and minimax risk of \cref{obj:def:sharp-rate,obj:def:sharp-decision,obj:def:minimax}, respectively. There exist real constants \(0<c\le C<\infty\), depending only on the fixed parameters, such that, simultaneously for every \(n,m\in\mathbb N\) with \(n\ge2\), \(T_*\) is Borel measurable and
\[
c\,r_*
\le R(n,m)
\le \sup_{P\in\mathcal M}\mathcal R_{n,m}(T_*,P)
\le C\,r_*.
\]
Moreover, every Borel real-valued decision \(T\) using the original dataset and the independent uniform randomizer satisfies
\[
c\,r_*\le \sup_{P\in\mathcal M}\mathcal R_{n,m}(T,P).
\]

Define the total treatment-record count \(N\), the combined nuisance smoothness \(S\), the exponent denominator \(\Delta\), the smoothness threshold \(S_{\mathrm{crit}}\), the total-record growth exponent \(q_*\), and the supplied-propensity order \(r_o(n)\) by
\[
N=n+m,\qquad S=\alpha+\beta,\qquad
\Delta=2\gamma+d+\frac{\gamma d}{\alpha+\beta},
\]
\[
S_{\mathrm{crit}}=\frac{\gamma d}{2\gamma+d},
\qquad
q_*=\frac{\gamma d}{(\alpha+\beta)(2\gamma+d)},
\qquad
r_o(n)=n^{-\gamma/(2\gamma+d)}.
\]
For these sample sizes, the sharp rate has the branches
\[
r_*=
\begin{cases}
r_o(n), & S\ge S_{\mathrm{crit}},\\
(nN)^{-\gamma/\Delta}, & S<S_{\mathrm{crit}},\quad N\le n^{q_*},\\
r_o(n), & S<S_{\mathrm{crit}},\quad N\ge n^{q_*}.
\end{cases}
\]
Whenever \(N=n^{q_*}\), the two rate formulas coincide:
\[
(nN)^{-\gamma/\Delta}=r_o(n).
\]
\end{theoremv}

The two terms express complementary precision requirements. The supplied-propensity term reflects the labeled-sample requirement for estimating the smooth contrast at a point; the product term reflects the joint requirement on labeled and treatment-record counts when the nuisance functions are unknown. If \(S\ge S_{\mathrm{crit}}\), the labeled count already supports the supplied-propensity order. Equality \(S=S_{\mathrm{crit}}\) belongs to this regime: then \(q_*=1\), and the necessary inequality \(N\ge n\) places every admissible sample pair at or beyond the saturation threshold. If \(S<S_{\mathrm{crit}}\), the exponent \(q_*>1\) creates an intermediate regime \(n\le N\le n^{q_*}\), in which precision depends on the product \(nN\). At \(N=n^{q_*}\), the two terms coincide; for larger total counts, the supplied-propensity term governs the rate.

The lower bound concerns the original records and includes every Borel decision using the independent uniform randomizer. Thus the comparison applies to both deterministic and randomized estimation procedures. The original-record construction in \cref{obj:thm:marked-component-certificate} underlies the sparse-range nuisance converse in \cref{obj:lem:nuisance-frontier-converse}. That converse supplies the product-of-counts lower bound in the range where it governs precision. Combined with the supplied-propensity lower bound below, it yields the sharp comparison across all sample imbalances. The matching upper bound is attained by the publicly tuned decision through \cref{obj:thm:rectangular-upper}.

To interpret saturation, the next result evaluates the same primitive class when the designated propensity function is supplied to the decision. This benchmark identifies the labeled-sample precision against which the value of outcome-free records is measured. Its attaining known-design decision uses
\[
h=n^{-1/(2\gamma+d)},\qquad
T^e(\mathcal D_{n,m},U;e_P)
=\mathrm{clip}_{[-1,1]}\!\left[
\frac1n\sum_{i=1}^n
w_h(X_i)\,r_0^\top r(X_i)
\left\{
\frac{A_iY_i}{e_P(X_i)}-
\frac{(1-A_i)Y_i}{1-e_P(X_i)}
\right\}
\right].
\]
Here \(r_0=r(x_0)\), and \(w_h\) and the degree-two basis \(r\) use the centered localization measure from \cref{obj:synth_3,obj:synth_1}. These formulas are defined for every bandwidth \(h>0\); the supplied-propensity smoother below uses only \(h_n=n^{-1/(2\gamma+d)}\), which lies in \((0,1)\) for \(n\ge2\) and can exceed \(1/2\) for small \(n\). This is a known-design polynomial evaluation average, not an empirical local least-squares fit.

\begin{theoremv}{thm:supplied-propensity}[Supplied-propensity point estimation rate]
Fix public parameters satisfying
\begin{itemize}
\item \textbf{(Dimension.)} \(d\in\mathbb N\) and \(d\ge1\).
\item \textbf{(Nuisance smoothness.)} \(0<\alpha\le1\) and \(0<\beta\le1\).
\item \textbf{(Effect smoothness.)} \(1\le\gamma\le3\).
\item \textbf{(Radius.)} \(L>1/2\).
\item \textbf{(Overlap.)} \(0<\varepsilon<1/2\).
\end{itemize}
Let \(\mathcal M\) be the primitive class at these parameters in \cref{obj:def:primitive-class}, and let \(R^e(n,m)\) be the supplied-propensity minimax absolute-error risk in \cref{obj:def:oracle-risk}. Define the supplied-propensity rate and the evaluation profile by
\[
r_o(n)=n^{-\gamma/(2\gamma+d)},
\qquad
x_0=(1/2,\ldots,1/2).
\]
There exist real constants \(0<c\le C<\infty\), depending only on the fixed public parameters, such that, for every \(n,m\in\mathbb N\) with \(n\ge2\),
\[
c\,r_o(n)\le R^e(n,m)\le C\,r_o(n).
\]

The upper bound is attained by the explicit degree-two smoother under the known uniform design, applied to the inverse-propensity pseudo-outcomes
\[
\frac{A_iY_i}{e_P(X_i)}
-
\frac{(1-A_i)Y_i}{1-e_P(X_i)},
\qquad i=1,\ldots,n,
\]
at bandwidth
\[
h=n^{-1/(2\gamma+d)},
\]
with final clipping to \([-1,1]\). Here the smoother uses its localization weight and coarse polynomial basis at this bandwidth for every \(h>0\), including \(h>1/2\). For every \(P\in\mathcal M\), the resulting decision \(T^e\) is measurable as a function of the dataset and randomizer when supplied with the designated admissible propensity \(e_P\), satisfies
\[
\left|T^e(\mathcal D_{n,m},U;e_P)\right|\le1
\]
for every dataset and randomizer value, and obeys
\[
\mathbb E_{\mathsf E_{n,m}(P)}
\left[
\left|T^e(\mathcal D_{n,m},U;e_P)-\tau_P(x_0)\right|
\right]
\le C\,r_o(n).
\]
The sampling law \(\mathsf E_{n,m}(P)\) is the product law of \(n\) independent labeled records, \(m\) independent outcome-free covariate-treatment records, and an independent uniform randomizer \(U\), as specified in \cref{obj:ass:sample-law}.

The same lower bound holds for the supplied-propensity minimax risk over the same-class subfamily with constant propensity and control mean:
\[
c\,r_o(n)\le
\inf_{T^e}
\sup_{\substack{P\in\mathcal M\\ e_P=\mu_{0,P}=1/2}}
\mathbb E_{\mathsf E_{n,m}(P)}
\left[
\left|T^e(\mathcal D_{n,m},U;e_P)-\tau_P(x_0)\right|
\right],
\]
where the constant-witness equalities hold on the unit cube and the infimum ranges over the supplied-propensity decisions of \cref{obj:def:oracle-risk}.
\end{theoremv}

The upper-bound localization formula uses the same centered cube and uniform localization measure as \cref{obj:synth_3}, evaluated at the single bandwidth \(h_n=n^{-1/(2\gamma+d)}\in(0,1)\), which may exceed \(1/2\). Because the cube has side length \(h_n<1\) and is centered at \(x_0=(1/2,\ldots,1/2)\), it lies inside \([0,1]^d\). The rectangular estimator in \cref{obj:thm:rectangular-upper} retains its stated restriction \(0<h\le1/2\).

With the propensity supplied, inverse-propensity weighting gives a response whose conditional mean is the causal contrast. Local polynomial estimation can therefore target the contrast's smoothness directly, yielding the order \(r_o(n)\). The same-class lower bound with constant propensity and control mean isolates the remaining response information requirement: even within this subfamily, the minimax risk is bounded below by a constant times \(r_o(n)\) for every auxiliary count. Since a supplied-propensity decision can also implement any decision based on the original data, this lower bound provides the labeled-sample floor in \cref{obj:thm:sharp-annotation-frontier}.

The sample requirements for attaining this benchmark can be expressed as a region of admissible counts. For a tolerance \leanref{sym:zeta}{\ensuremath{\zeta}}, the annotation region \leanref{sym:Bstar}{\ensuremath{\mathcal B_*(\zeta)}} records the pairs whose sharp rate lies within that tolerance of the supplied-propensity order.

\begin{definitionv}{def:annotation-region}[Annotation region \(\mathcal B_*(\zeta)\)]
The annotation region at tolerance \(\zeta\) is
\[
\mathcal B_*(\zeta)=
\left\{
(n,m)\in\mathbb N_{\ge2}\times\mathbb N_0:
r_*(n,m;d,\alpha,\beta,\gamma)\le\zeta r_o(n)
\right\},
\]
where \(r_*\) is the sharp rate and \(r_o(n)\) is the supplied-propensity point-estimation order.
\end{definitionv}

For \(\zeta\ge1\), the explicit rate gives the equivalent count condition
\[
(n,m)\in\mathcal B_*(\zeta)
\quad\Longleftrightarrow\quad
n+m\ge \zeta^{-\Delta/\gamma}n^{q_*},
\qquad n\ge2,\quad m\ge0.
\]
In particular, \(\zeta=1\) gives the exact rate-level saturation region. Allowing a fixed finite tolerance gives attainment of the supplied-propensity order up to constants. The complete sequence equivalence, strict-improvement criterion, and target-error conditions are collected in \cref{obj:prop:annotation-threshold} and proved in \cref{sec:acquisition-consequences}.

Attainment of the supplied-propensity order and strict improvement over the labeled-only order impose distinct count requirements. Under \(S<S_{\mathrm{crit}}\), strict improvement occurs exactly when \(m_n/n\to+\infty\), whereas supplied-propensity attainment requires the total count to remain at least a positive constant times \(n^{q_*}\). Consequently, diverging auxiliary-to-labeled ratios can improve the minimax order while leaving precision in the intermediate product-of-counts regime. A bounded auxiliary-to-labeled ratio preserves the labeled-only order up to constants. At and above the smoothness threshold, every auxiliary-size sequence attains the supplied-propensity order.

For acquisition planning, \cref{obj:prop:annotation-threshold} gives two simultaneous requirements. For a target absolute-error scale \(0<\eta<1\), the labeled count must support point-estimation precision at scale \(\eta^{-(2+d/\gamma)}\), and the product of labeled and total counts must support nuisance adjustment at scale \(\eta^{-(2+d/\gamma+d/S)}\). These requirements are necessary up to a constant depending on \((d,\alpha,\beta,\gamma,L,\varepsilon)\) and sufficient up to another such constant for uniform expected absolute error at most \(\eta\). Their rate-level form is exact. Thus increasing either count can satisfy the product requirement, while the labeled-count requirement specifies the response information needed at the target precision.

A one-dimensional example makes the threshold concrete. Set \(d=1\), \(\gamma=1\), and \(\alpha=\beta=1/8\), with any fixed \(L>1/2\) and \(\varepsilon\in(0,1/2)\). Then
\[
S=\frac14,\qquad
S_{\mathrm{crit}}=\frac13,\qquad
\Delta=7,\qquad
q_*=\frac43,
\]
and
\[
r_*(n,m)=\max\left\{n^{-1/3},[n(n+m)]^{-1/7}\right\}.
\]
The labeled-only rate is \(n^{-2/7}\), and the exact branch boundary is \(N=n^{4/3}\). Total counts of order \(n^{4/3}\) therefore attain the supplied-propensity order up to constants. For example, with \(n=1000\), the branch boundary is \(N=10000\), corresponding to \(m=9000\) outcome-free records. These counts describe equality of the rate formulas; the minimax risk retains constants that may depend on \(L\) and \(\varepsilon\). In the same example, the exact rate-level requirements for target scale \(\eta\) are
\[
n\ge\eta^{-3},
\qquad
nN\ge\eta^{-7}.
\]
Taking the labeled count at order \(\eta^{-3}\) consequently requires a total count at order \(\eta^{-4}\), illustrating how response labels and routine covariate-treatment records jointly determine attainable precision.


\section{Discussion and limitations}\label{sec:discussion}

The acquisition region in \cref{obj:def:annotation-region} provides a benchmark for allocating response labels and routine covariate-treatment records under the model in \cref{obj:def:primitive-class,obj:ass:sample-law}. It identifies sample-size pairs for which the sharp rate is within a specified factor of the supplied-propensity order. Through \cref{obj:thm:sharp-annotation-frontier}, this rate comparison translates into a minimax precision comparison up to constants depending on the public class parameters. For an absolute-error target, \cref{obj:prop:annotation-threshold} instead gives simultaneous requirements on the labeled count and the product of labeled and total treatment-record counts. These requirements offer a basis for comparing acquisition plans once their costs are specified.

A chart-review setting illustrates the distinction between the two record types. A labeled record contains covariates, treatment, and a true response ascertained through review; an outcome-free record contains the covariates and treatment available from routine records. Adding a new independent labeled record changes \((n,m)\) to \((n+1,m)\) and increases the total count. Reviewing a response for one record already present in the auxiliary block changes the composition to \((n+1,m-1)\) at fixed total count; if the record index is fixed in advance, or chosen by a randomization independent of all records, the resulting dataset follows the stated sampling law at \((n+1,m-1)\). Adding a new routine record changes \((n,m)\) to \((n,m+1)\). The sample requirements benchmark these alternative compositions under the specified independent, common-population experiment; acquisition costs determine which transition is relevant in an application.

The two channels address distinct sources of statistical difficulty. Response labels carry information about the probability-valued conditional response means and the causal contrast. The arm-mean conditions in \cref{obj:ass:control-interior,obj:ass:treated-interior} place both means in \([0,1]\), and \cref{obj:thm:supplied-propensity} characterizes the labeled-sample precision when the propensity function is supplied. Outcome-free records carry information about treatment assignment conditional on covariates. Their contribution enters the second component of the sharp rate in \cref{obj:def:sharp-rate}, where the product of labeled and total counts governs the cost of nuisance adjustment. Increasing the routine-record count can reduce that component until the supplied-propensity order determines the overall rate.

Public nuisance regularity determines how much routine information is needed to reach this benchmark. The rate exponent depends on the sum of the propensity and control-mean H\"older orders, while the effect H\"older order governs the supplied-propensity precision. Greater nuisance regularity can reduce the total-record growth needed at a given labeled count. In the sufficiently smooth regime characterized by \cref{obj:thm:sharp-annotation-frontier}, labeled records alone attain the supplied-propensity order; in the lower-regularity regime, additional outcome-free records can improve the minimax order until saturation. These comparisons concern rates, with constants depending on the fixed class parameters. Because the regularity parameters are public, they also determine the estimator's localization and projection choices in \cref{obj:def:tuning}. Acquisition planning and estimator tuning thus use the same description of the statistical model.

\subsection{Limitations and future work}

The analysis uses a known uniform covariate distribution on the unit cube, a fixed public overlap level \(\varepsilon\in(0,1/2)\), binary responses with arm means in \([0,1]\), and independent record blocks from a common population. Extending the acquisition benchmark to general, potentially unknown covariate distributions is a direction for future work. Covariate-dependent or adaptive selection of records for review would also change the labeled-record law and requires a sampling or missingness model beyond the experiment studied here. Covariate shift between reviewed and routine records would introduce an additional transport problem. These extensions would help connect the benchmark to populations with uneven support and heterogeneous record availability.

Unknown smoothness raises a separate question of adaptation. The present tuning and acquisition comparisons use public H\"older orders and a public radius. Procedures that select tuning parameters from the data would need risk guarantees over multiple regularity classes, together with an assessment of any adaptation cost. Uncertainty quantification also requires further analysis: the expected absolute-error bounds characterize point-estimation precision, whereas confidence intervals require coverage guarantees and control of approximation bias.

Alternative losses and causal targets may change the allocation requirements. Squared-error loss, integrated estimation of the treatment-effect function, average treatment effects, and policy-value targets each call for their own precision benchmark. Richer auxiliary information is another direction, including response proxies, repeated measurements, and additional variables predictive of the response. Such information would enlarge the observation model beyond covariate-treatment pairs and require an explicit account of how the auxiliary measurements relate to the true response.


\appendix

\section*{Appendices}

\section{Upper-bound proofs}\label{sec:upper-bound-proofs}

The upper bounds in \cref{obj:thm:rectangular-upper,obj:thm:supplied-propensity} combine local polynomial approximation with control of the corresponding empirical moments. For the rectangular estimator, the two-scale construction in \cref{obj:def:projection} and the public role split in \cref{obj:def:roles} are statistical inputs: the coarse polynomial targets the smooth contrast, the fine projection controls nuisance approximation, and the independent roles separate response information from treatment-assignment information. The auxiliary results below follow this dependency order.

\subsection{Local approximation and projection}

Fix a primitive law \(P\) from \cref{obj:def:primitive-class} and admissible localization and resolution inputs \(h,J\). Retain the notation \(\bar R\) for the expected response vector and \(Q_{P,h,J}\) for the expected symmetrized matrix:
\[
\bar R=\mathbb E_{\mathsf E_{n,m}(P)}\widehat R_{h,J},
\qquad
Q_{P,h,J}=\mathbb E_{\mathsf E_{n,m}(P)}\widehat Q_{h,J}.
\]
Both expectations concern the empirical moments in \cref{obj:def:moments}. Let \(p\) denote the Taylor polynomial of \(\tau_P\) at \(x_0\), with degree \(p_*=\lceil\gamma\rceil-1\), and let \(\vartheta\) denote its coefficient vector in the coarse basis. Since \(1\le\gamma\le3\), this Taylor degree lies within the degree-two polynomial space used by the estimator. The following result collects the approximation and projection properties needed for the population estimating equations and their sampling bounds.

\begin{lemmav}{lem:local-polynomial-projection-facts}[Local polynomial projection bounds]
Fix public parameters \(d\in\mathbb N\) and \(\alpha,\beta,\gamma,L,\varepsilon\in\mathbb R\) satisfying:
\begin{itemize}
\item \textbf{(Dimension.)} \(d\ge1\).
\item \textbf{(Smoothness.)} \(0<\alpha\le1\), \(0<\beta\le1\), and \(1\le\gamma\le3\).
\item \textbf{(Radius.)} \(L>1/2\).
\item \textbf{(Overlap.)} \(0<\varepsilon<1/2\).
\end{itemize}
There exists a finite constant \(C=C(d,\alpha,\beta,\gamma,L,\varepsilon)>0\) such that, for every primitive law \(P\in\mathcal M(d,\alpha,\beta,\gamma,L,\varepsilon)\) as defined in \cref{obj:def:primitive-class}, every \(h\in\mathbb R\) with \(0<h\le1/2\), and every integer \(J\ge1\), the following conclusions hold for the bases and projection of \cref{obj:def:projection}.

Write \(x_0=(1/2,\ldots,1/2)\) for the evaluation center, \(C_h=[1/2-h/2,1/2+h/2]^d\) for the localization cube, and \(\nu_h\) for its uniform probability measure. Let \(p\) be the Taylor polynomial of the canonical contrast \(\tau_P\) at \(x_0\), with Taylor degree \(p_*=\lceil\gamma\rceil-1\):
\[
p(x)=
\sum_{\substack{\kappa\in\mathbb N^d\\ \sum_{i=1}^d\kappa_i\le p_*}}
\frac{\partial^\kappa\tau_P(x_0)}
     {\prod_{i=1}^d\kappa_i!}
\prod_{i=1}^d(x_i-1/2)^{\kappa_i}.
\]
There exists a coarse coefficient vector \(\vartheta\) whose represented polynomial
\[
p_\vartheta(x)=r(x)^\top\vartheta
\]
satisfies
\[
\sup_{x\in C_h}|\tau_P(x)-p_\vartheta(x)|\le Ch^\gamma,
\qquad
\|\vartheta\|_2\le C,
\qquad
p_\vartheta(x_0)=\tau_P(x_0),
\]
and
\[
p_\vartheta(x)=p(x)\qquad\text{for every }x\in C_h.
\]

The coarse basis vector \(r\) and fine projection kernel \(K_J\) satisfy
\[
\sup_{x\in C_h}\|r(x)\|_2\le C,
\qquad
\sup_{x\in C_h}\int |K_J(x,x')|\,d\nu_h(x')\le C,
\]
and
\[
\iint K_J(x,x')^2\,d\nu_h(x)\,d\nu_h(x')
=qJ^d,
\]
where \(q\) is the number of coarse multi-indices of total degree at most two:
\[
q=\#\left\{\kappa\in\mathbb N^d:\sum_{i=1}^d\kappa_i\le2\right\}.
\]

Finally, with fine-cell side length \(\delta=h/J\), every coarse basis coordinate \(r_u\) satisfies
\[
\|e_Pr_u-\Pi_J(e_Pr_u)\|_{L_2(\nu_h)}
\le C\delta^\alpha,
\]
where \(e_P\) is the designated admissible propensity witness, and the designated admissible control-mean witness \(\mu_{0,P}\) satisfies
\[
\|\mu_{0,P}-\Pi_J\mu_{0,P}\|_{L_2(\nu_h)}
\le C\delta^\beta.
\]
\end{lemmav}

\begin{proof}[Proof of \cref{obj:lem:local-polynomial-projection-facts}]
All constants constructed below depend only on the fixed public parameters. Write
\[
\begin{gathered}
\mathcal U_{\mathrm{coarse}}
 =\left\{\kappa\in\mathbb N^d:|\kappa|\le2\right\},
\qquad
|\kappa|=\sum_{i=1}^d\kappa_i,
\qquad
\kappa!=\prod_{i=1}^d\kappa_i!,
\qquad
q=|\mathcal U_{\mathrm{coarse}}|,\\
\mathsf L_0(\upsilon)=1,\qquad
\mathsf L_1(\upsilon)=\sqrt{12}\,\upsilon,\qquad
\mathsf L_2(\upsilon)=\sqrt5\,(6\upsilon^2-1/2)
\quad(\upsilon\in\mathbb R).
\end{gathered}
\]
Thus the coarse coordinates are
\[
r_\kappa(x)
 =\prod_{i=1}^d
   \mathsf L_{\kappa_i}\!\left(\frac{x_i-1/2}{h}\right),
\qquad
p_\vartheta(x)=\sum_{\kappa\in\mathcal U_{\mathrm{coarse}}}
r_\kappa(x)\vartheta_\kappa
\quad(x\in\mathbb R^d).
\]
We use the designated functions supplied by \cref{obj:def:primitive-class} and the localized measure and projection of \cref{obj:def:projection}.

\begin{enumerate}
\item \emph{Taylor coefficients and approximation.}
We first construct a positive constant \(C_{\mathrm T}\), uniform over admitted laws and localization scales.

If \(\gamma=1\), then \(p_*=0\). Let
\[
\mathbf0=(0,\ldots,0),\qquad
\vartheta_\kappa=
\begin{cases}
\tau_P(x_0),&\kappa=\mathbf0,\\
0,&\kappa\ne\mathbf0,
\end{cases}
\qquad
C_{\mathrm T}=L(\sqrt d+1).
\]
Since \(r_{\mathbf0}=1\), this gives
\[
p_\vartheta(x)=\tau_P(x_0)=p(x).
\]
By \cref{obj:def:holder-norm,obj:def:primitive-class},
\[
|\tau_P(x_0)|\le L,
\qquad
|\tau_P(x)-\tau_P(x')|\le L\|x-x'\|_2
\quad(x,x'\in[0,1]^d).
\]
For equal arguments the second inequality follows from its zero left-hand side; for distinct arguments it is the order-zero Hölder bound. Moreover, for \(x\in C_h\),
\[
\|x-x_0\|_2^2
 =\sum_{i=1}^d|x_i-1/2|^2
 \le dh^2.
\]
The inclusion \(C_h\subset[0,1]^d\) follows from \(h\le1/2\). Consequently,
\[
|\tau_P(x)-p_\vartheta(x)|
 \le L\sqrt d\,h
 \le C_{\mathrm T}h,
\qquad
\|\vartheta\|_2=|\tau_P(x_0)|\le L\le C_{\mathrm T}.
\]
The representation also gives \(p_\vartheta(x_0)=\tau_P(x_0)\).

Suppose now that \(\gamma>1\). Define
\[
p_*=\lceil\gamma\rceil-1,\qquad
s_{\mathrm T}=\gamma-p_*,
\qquad
D_{\mathrm{dim}}=d+1.
\]
The public range of \(\gamma\) implies
\[
p_*\in\{1,2\},\qquad
0<s_{\mathrm T}\le1,\qquad
p_*+s_{\mathrm T}=\gamma.
\]
In particular, at \(\gamma=2\) and \(\gamma=3\), the pairs
\((p_*,s_{\mathrm T})\) are \((1,1)\) and \((2,1)\), respectively.
The same two definitions supply continuous derivatives through order \(p_*\) and the bounds
\[
\begin{aligned}
|\partial^\kappa\tau_P(x)|&\le L
&&(|\kappa|\le p_*,\ x\in[0,1]^d),\\
|\partial^\kappa\tau_P(x)-\partial^\kappa\tau_P(x')|
 &\le L\|x-x'\|_2^{s_{\mathrm T}}
&&(|\kappa|=p_*,\ x,x'\in[0,1]^d).
\end{aligned}
\]
The second inequality includes equal arguments, where both sides vanish.

For \(x\in C_h\), set
\[
z_x=x-x_0,\qquad
F_x(\upsilon)=\tau_P(x_0+\upsilon z_x)
\quad(0\le\upsilon\le1).
\]
The segment lies in the design cube. The chain rule, together with symmetry of the continuous derivatives, yields
\[
F_x^{(p_*)}(\upsilon)
 =
 \sum_{(i_1,\ldots,i_{p_*})\in\{1,\ldots,d\}^{p_*}}
 \left(\prod_{j=1}^{p_*}z_{x,i_j}\right)
 (\partial_{i_1}\cdots\partial_{i_{p_*}}\tau_P)
 (x_0+\upsilon z_x).
\]
Each coordinate of \(z_x\) has absolute value at most \(\|z_x\|_2\).
There are \(d^{p_*}\le D_{\mathrm{dim}}^{p_*}\) terms, so the top-derivative modulus gives
\[
\left|F_x^{(p_*)}(\upsilon)-F_x^{(p_*)}(\upsilon')\right|
 \le D_{\mathrm{dim}}^{p_*}L
       \|z_x\|_2^{p_*+s_{\mathrm T}}
       |\upsilon-\upsilon'|^{s_{\mathrm T}}
\quad(0\le\upsilon,\upsilon'<1).
\]
Grouping the coordinate derivatives by their multiplicities gives
\[
p(x)
 =\sum_{j=0}^{p_*}\frac{F_x^{(j)}(0)}{j!}
 =\sum_{|\kappa|\le p_*}
   \frac{\partial^\kappa\tau_P(x_0)}{\kappa!}
   \prod_{i=1}^d(x_i-1/2)^{\kappa_i}.
\]
The Lagrange remainder for the Taylor polynomial of degree \(p_*-1\)
provides \(\upsilon_x\in(0,1)\) such that, after subtracting the degree-\(p_*\) term,
\[
\tau_P(x)-p(x)
 =\frac{F_x^{(p_*)}(\upsilon_x)-F_x^{(p_*)}(0)}{p_*!}.
\]
Since \(p_*!\ge1\) and \(\upsilon_x^{s_{\mathrm T}}\le1\), this implies
\[
|\tau_P(x)-p(x)|
 \le D_{\mathrm{dim}}^{p_*}L\|z_x\|_2^\gamma.
\]
To obtain a constant uniform in the fractional order, observe that
\[
\|z_x\|_2\le\sqrt d\,\frac h2\le D_{\mathrm{dim}}h,
\qquad
D_{\mathrm{dim}}\ge1,
\qquad
\gamma=p_*+s_{\mathrm T}\le p_*+1.
\]
Hence, with
\[
C_{\mathrm{rem}}
 =D_{\mathrm{dim}}^{p_*}D_{\mathrm{dim}}^{p_*+1},
\]
we have
\[
|\tau_P(x)-p(x)|
 \le C_{\mathrm{rem}}Lh^\gamma
 \quad(x\in C_h).
\]

We next represent \(p\) in the coarse basis by an explicit change of coefficients.
Define the matrix, whose rows and columns are indexed by \(0,1,2\),
\[
\mathsf M_h=
\begin{pmatrix}
1&0&0\\
0&h/\sqrt{12}&0\\
h^2/12&0&h^2/(6\sqrt5)
\end{pmatrix}.
\]
Direct substitution gives
\[
(h\upsilon)^j
 =\sum_{\ell=0}^2(\mathsf M_h)_{j\ell}\mathsf L_\ell(\upsilon)
 \quad(j\in\{0,1,2\},\ \upsilon\in\mathbb R).
\]
Its entries satisfy
\[
(\mathsf M_h)_{j\ell}=0\quad(\ell>j),
\qquad
|(\mathsf M_h)_{j\ell}|\le h^j.
\]
Put
\[
\begin{gathered}
\mathcal U_{\mathrm T}
 =\{\kappa\in\mathbb N^d:|\kappa|\le p_*\},
\qquad
N_{\mathrm T}=|\mathcal U_{\mathrm T}|,\\
a_{\mathrm T,\kappa}
 =\frac{\partial^\kappa\tau_P(x_0)}{\kappa!}
 \quad(\kappa\in\mathcal U_{\mathrm T}),\\
\mathsf M_h[\kappa,\kappa']
 =\prod_{i=1}^d(\mathsf M_h)_{\kappa_i\kappa'_i}
 \quad(\kappa,\kappa'\in\mathcal U_{\mathrm T}).
\end{gathered}
\]
Since \(\kappa!\ge1\),
\[
|a_{\mathrm T,\kappa}|\le L,
\qquad
|\mathsf M_h[\kappa,\kappa']|
 \le h^{|\kappa|}\le1.
\]
Expanding the product of the one-dimensional identities gives
\[
\prod_{i=1}^d(x_i-1/2)^{\kappa_i}
 =\sum_{\kappa'\in\mathcal U_{\mathrm T}}
   \mathsf M_h[\kappa,\kappa']r_{\kappa'}(x).
\]
Indeed, a nonzero term has \(\kappa'_i\le\kappa_i\) in every coordinate,
so its total degree remains at most \(p_*\).

Define the coarse vector by zero extension:
\[
\vartheta_{\kappa'}=
\begin{cases}
\displaystyle\sum_{\kappa\in\mathcal U_{\mathrm T}}
 a_{\mathrm T,\kappa}\mathsf M_h[\kappa,\kappa'],
 &\kappa'\in\mathcal U_{\mathrm T},\\
0,&\kappa'\in\mathcal U_{\mathrm{coarse}}\setminus\mathcal U_{\mathrm T}.
\end{cases}
\]
Exchanging the finite sums proves
\[
p_\vartheta(x)=p(x)\quad(x\in\mathbb R^d).
\]
Furthermore,
\[
|\vartheta_{\kappa'}|\le N_{\mathrm T}L,
\qquad
\sum_{\kappa'\in\mathcal U_{\mathrm{coarse}}}
 \vartheta_{\kappa'}^2
 \le N_{\mathrm T}(N_{\mathrm T}L)^2.
\]
Thus, defining
\[
C_{\mathrm{coef}}=\sqrt{N_{\mathrm T}}\,N_{\mathrm T}+1,
\qquad
C_{\mathrm T}=L\max\{C_{\mathrm{rem}},C_{\mathrm{coef}}\},
\]
we obtain
\[
\|\vartheta\|_2\le C_{\mathrm{coef}}L\le C_{\mathrm T},
\qquad
\sup_{x\in C_h}|\tau_P(x)-p_\vartheta(x)|
 \le C_{\mathrm T}h^\gamma.
\]
At \(x_0\), every positive-degree centered monomial vanishes and the
zero-degree coefficient is \(\tau_P(x_0)\); therefore
\(p_\vartheta(x_0)=\tau_P(x_0)\).
Both cases construct a finite positive \(C_{\mathrm T}\).
% lean: CausalSmith.Stat.TwosamplePointcateAnnotationFrontier.taylor_coefficient_bounds

\item \emph{Basis and kernel bounds.}
For \(|\upsilon|\le1/2\), the three reference polynomials satisfy
\[
|\mathsf L_0(\upsilon)|=1,\qquad
|\mathsf L_1(\upsilon)|\le\frac{\sqrt{12}}2\le4,\qquad
|\mathsf L_2(\upsilon)|\le\sqrt5\le4.
\]
For the last inequality we used
\(0\le\upsilon^2\le1/4\), which gives
\(|6\upsilon^2-1/2|\le1\). It follows that
\[
|r_\kappa(x)|\le4^d,
\qquad
\|r(x)\|_2^2\le q(4^d)^2
\quad(x\in C_h).
\]

Here is the fine-cell apparatus used for the kernel calculations.
Define
\[
\delta=\frac hJ>0,\qquad
\mathcal Z_J=\{0,\ldots,J-1\}^d,
\qquad
c_{\mathbf j,i}=1/2-h/2+(j_i+1/2)\delta
\quad(\mathbf j\in\mathcal Z_J).
\]
For each coordinate, let
\[
I_{j_i}=
\begin{cases}
[1/2-h/2+j_i\delta,\;1/2-h/2+(j_i+1)\delta),
 &j_i<J-1,\\
[1/2-h/2+(J-1)\delta,\;1/2+h/2],
 &j_i=J-1,
\end{cases}
\qquad
D_{\mathbf j}=\prod_{i=1}^d I_{j_i}.
\]
These cells partition \(C_h\), with every terminal face assigned to the
last interval. To see coverage explicitly, for \(x\in C_h\) define
\[
t_{x,i}=\frac{x_i-(1/2-h/2)}{\delta}\in[0,J],
\qquad
j_i(x)=
\begin{cases}
J-1,&t_{x,i}=J,\\
\lfloor t_{x,i}\rfloor,&t_{x,i}<J.
\end{cases}
\]
Then \(x\in D_{\mathbf j(x)}\). Distinct indices give disjoint cells:
in a coordinate where their indices differ, the smaller interval's right
endpoint is at or before the larger interval's left endpoint, and the
smaller interval excludes its right endpoint.
These conventions also apply when \(J=1\).

The fine coordinates are
\[
b_{\mathbf j,\kappa}(x)=
\begin{cases}
\displaystyle
\sqrt{J^d}\prod_{i=1}^d
 \mathsf L_{\kappa_i}\!\left(\frac{x_i-c_{\mathbf j,i}}{\delta}\right),
 &x\in D_{\mathbf j},\\
0,&x\notin D_{\mathbf j},
\end{cases}
\quad
(\mathbf j\in\mathcal Z_J,\ \kappa\in\mathcal U_{\mathrm{coarse}}).
\]
In particular, all fine coordinates vanish outside \(C_h\).
The orthonormality specified in \cref{obj:def:projection} reads
\[
\int b_{\mathbf j,\kappa}(x)b_{\mathbf j',\kappa'}(x)\,d\nu_h(x)
 =
\begin{cases}
1,&(\mathbf j,\kappa)=(\mathbf j',\kappa'),\\
0,&(\mathbf j,\kappa)\ne(\mathbf j',\kappa').
\end{cases}
\]
Its normalization can also be seen directly: the reference integrals satisfy
\[
\int_{-1/2}^{1/2}\mathsf L_j(\upsilon)\mathsf L_\ell(\upsilon)\,d\upsilon
 =\mathbf1_{\{j=\ell\}}
 \quad(j,\ell\in\{0,1,2\}),
\]
by integrating the displayed polynomials. On a common cell, the
change of variables \(x_i=c_{\mathbf j,i}+\delta\upsilon_i\) contributes
\[
h^{-d}J^d\delta^d=1.
\]
Different cells have disjoint supports, and the assigned faces have zero
Lebesgue measure.

The coordinate bound gives, for every \(x\in\mathbb R^d\),
\[
|b_{\mathbf j,\kappa}(x)|\le\sqrt{J^d}\,4^d.
\]
The constant coordinate on a cell is
\[
b_{\mathbf j,\mathbf0}(x)=\sqrt{J^d}\,\mathbf1_{D_{\mathbf j}}(x).
\]
Consequently,
\[
|b_{\mathbf j,\kappa}(x)|
 \le\frac{4^d}{\sqrt{J^d}}b_{\mathbf j,\mathbf0}(x)^2,
\qquad
\int|b_{\mathbf j,\kappa}(x)|\,d\nu_h(x)
 \le\frac{4^d}{\sqrt{J^d}},
\]
where the second inequality uses
\(\int b_{\mathbf j,\mathbf0}^2\,d\nu_h=1\).

If \(x\in D_{\mathbf j}\), disjointness reduces the kernel to
\[
K_J(x,x')
 =\sum_{\kappa\in\mathcal U_{\mathrm{coarse}}}
 b_{\mathbf j,\kappa}(x)b_{\mathbf j,\kappa}(x').
\]
The triangle inequality and the preceding bounds therefore give
\[
\begin{aligned}
\int|K_J(x,x')|\,d\nu_h(x')
&\le
 \sum_{\kappa\in\mathcal U_{\mathrm{coarse}}}
 |b_{\mathbf j,\kappa}(x)|
 \int|b_{\mathbf j,\kappa}(x')|\,d\nu_h(x')\\
&\le q\bigl(\sqrt{J^d}\,4^d\bigr)
          \frac{4^d}{\sqrt{J^d}}
 =q(4^d)^2.
\end{aligned}
\]
For \(x\) outside all cells, \(K_J(x,x')=0\), so the same bound holds there.

There are \(|\mathcal Z_J|=J^d\) cells and \(q\) coordinates per cell;
hence the fine index set has \(qJ^d\) elements. Expanding the squared
finite kernel and factoring the product integrals yields
\[
\begin{aligned}
\iint K_J(x,x')^2\,d\nu_h(x)\,d\nu_h(x')
&=
 \sum_{\mathbf j,\kappa}\sum_{\mathbf j',\kappa'}
 \left(\int
 b_{\mathbf j,\kappa}(x)b_{\mathbf j',\kappa'}(x)\,d\nu_h(x)\right)^2\\
&=qJ^d=k.
\end{aligned}
\]
All terms are integrable by the square integrability of the basis coordinates.
Define
\[
C_{\mathrm B}
 =\max\{\sqrt q\,4^d+1,\ q(4^d)^2+1\}>0.
\]
We have proved
\[
\sup_{x\in C_h}\|r(x)\|_2\le C_{\mathrm B},
\qquad
\sup_{x\in C_h}\int|K_J(x,x')|\,d\nu_h(x')\le C_{\mathrm B},
\]
together with the exact squared-kernel identity.
% lean: CausalSmith.Stat.TwosamplePointcateAnnotationFrontier.kernel_basis_bounds

\item \emph{Control-mean residual.}
For square-integrable functions, write
\[
\|f\|_h=\left(\int f(x)^2\,d\nu_h(x)\right)^{1/2}.
\]
The coarse constant coordinate equals one, so coarse orthonormality gives
\[
\nu_h(C_h)=\int1\,d\nu_h
 =\int r_{\mathbf0}(x)^2\,d\nu_h(x)=1.
\]
Thus \(\nu_h\) is a probability measure supported on \(C_h\).

The centers satisfy \(c_{\mathbf j}\in D_{\mathbf j}\).
For two points in the same cell, coordinatewise interval bounds give
\[
|x_i-x_i'|\le\delta,
\qquad
\|x-x'\|_2^2\le d\delta^2
\quad(x,x'\in D_{\mathbf j}).
\]
By \cref{obj:def:holder-norm,obj:def:primitive-class} and
\(0<\beta\le1\),
\[
|\mu_{0,P}(x)|\le L,
\qquad
|\mu_{0,P}(x)-\mu_{0,P}(x')|
 \le L\|x-x'\|_2^\beta
\quad(x,x'\in C_h).
\]
Indeed, \(\lceil\beta\rceil-1=0\), so the Hölder norm bounds the
function and its order-zero modulus; equal arguments satisfy the modulus
inequality directly.

Define the fine coefficient vector and its expansion by
\[
v^{(0,J)}_{\mathbf j,\kappa}=
\begin{cases}
\mu_{0,P}(c_{\mathbf j})/\sqrt{J^d},&\kappa=\mathbf0,\\
0,&\kappa\ne\mathbf0,
\end{cases}
\qquad
f_{0,J}(x)=
 \sum_{\mathbf j,\kappa}
 b_{\mathbf j,\kappa}(x)v^{(0,J)}_{\mathbf j,\kappa}.
\]
The cell assignments imply
\[
f_{0,J}(x)=
\begin{cases}
\mu_{0,P}(c_{\mathbf j}),&x\in D_{\mathbf j},\\
0,&x\notin C_h.
\end{cases}
\]
Set
\[
E_{0,J}=L(\sqrt d)^\beta\delta^\beta\ge0.
\]
For \(x\in D_{\mathbf j}\), the modulus and distance bounds give
\[
|\mu_{0,P}(x)-f_{0,J}(x)|
 \le L\|x-c_{\mathbf j}\|_2^\beta
 \le L(\sqrt d\,\delta)^\beta
 =E_{0,J}.
\]
Because \(\nu_h\) is a probability measure,
\[
\|\mu_{0,P}-f_{0,J}\|_h^2
 \le\int E_{0,J}^2\,d\nu_h=E_{0,J}^2.
\]

We now derive the projection comparison used here and in the next step.
For any \(f\in L_2(\nu_h)\) and any real fine coefficient vector
\(v_J=(v_{J,\mathbf j,\kappa})\), put
\[
f_J(x)=\sum_{\mathbf j,\kappa}
 b_{\mathbf j,\kappa}(x)v_{J,\mathbf j,\kappa}.
\]
Expanding the kernel integral in \cref{obj:def:projection} gives
\[
\Pi_Jf(x)=
 \sum_{\mathbf j,\kappa}b_{\mathbf j,\kappa}(x)
 \int b_{\mathbf j,\kappa}(x')f(x')\,d\nu_h(x').
\]
Orthonormality then implies
\[
\int b_{\mathbf j,\kappa}(x)f_J(x)\,d\nu_h(x)
 =v_{J,\mathbf j,\kappa},
\qquad
\Pi_J(f-f_J)=\Pi_Jf-f_J,
\]
and
\[
\int b_{\mathbf j,\kappa}(x)
       \bigl(f(x)-\Pi_Jf(x)\bigr)\,d\nu_h(x)=0.
\]
The last identity makes the cross term in the squared norm vanish, giving
\[
\|f\|_h^2=\|\Pi_Jf\|_h^2+\|f-\Pi_Jf\|_h^2.
\]
Apply this identity to \(f-f_J\) and use the displayed projection identity:
\[
\|f-f_J\|_h^2
 =\|\Pi_Jf-f_J\|_h^2+\|f-\Pi_Jf\|_h^2
 \ge\|f-\Pi_Jf\|_h^2.
\]
The bounded control mean and its finite-basis approximant are square
integrable. Taking \(f=\mu_{0,P}\) and \(f_J=f_{0,J}\), and then taking
square roots, proves
\[
\|\mu_{0,P}-\Pi_J\mu_{0,P}\|_h
 \le E_{0,J}.
\]
Therefore the positive constant
\[
C_{\mathrm m}=L(\sqrt d)^\beta+1
\]
satisfies
\[
\|\mu_{0,P}-\Pi_J\mu_{0,P}\|_h
 \le C_{\mathrm m}\delta^\beta.
\]
% lean: CausalSmith.Stat.TwosamplePointcateAnnotationFrontier.control_projection_bounds

\item \emph{Propensity-weighted residuals.}
Similarly, \(0<\alpha\le1\) and
\cref{obj:def:holder-norm,obj:def:primitive-class} give
\[
|e_P(x)|\le L,
\qquad
|e_P(x)-e_P(x')|\le L\|x-x'\|_2^\alpha
\quad(x,x'\in C_h).
\]
Fix \(u\in\mathcal U_{\mathrm{coarse}}\), and define
\[
f_{e,u}(x)=e_P(x)r_u(x)
\quad(x\in\mathbb R^d).
\]
On \(C_h\), \(|f_{e,u}(x)|\le L4^d\), so it is square integrable
under \(\nu_h\).

We construct a fine-basis expansion that equals
\(e_P(c_{\mathbf j})r_u\) on each cell. Define
\[
\zeta_{\mathbf j,i}=\frac{c_{\mathbf j,i}-1/2}{h},
\qquad
\sigma_J=\frac1J,
\qquad
\mathsf A_{\mathbf j,i}=
\begin{pmatrix}
1&0&0\\
\sqrt{12}\,\zeta_{\mathbf j,i}&\sigma_J&0\\
\sqrt5(6\zeta_{\mathbf j,i}^2+\sigma_J^2/2-1/2)
 &\sqrt5\sqrt{12}\,\zeta_{\mathbf j,i}\sigma_J&\sigma_J^2
\end{pmatrix}.
\]
Rows and columns are indexed by \(0,1,2\). Substitution into the three
reference polynomials verifies
\[
\mathsf L_j(\zeta_{\mathbf j,i}+\sigma_J\upsilon_{\mathrm{aff}})
 =
 \sum_{\ell=0}^j
 (\mathsf A_{\mathbf j,i})_{j\ell}
 \mathsf L_\ell(\upsilon_{\mathrm{aff}})
\quad(j\in\{0,1,2\},\ \upsilon_{\mathrm{aff}}\in\mathbb R).
\]
For \(x\in D_{\mathbf j}\),
\[
\frac{x_i-1/2}{h}
 =\zeta_{\mathbf j,i}
   +\sigma_J\frac{x_i-c_{\mathbf j,i}}{\delta}.
\]
Define the tensor coefficients
\[
\mathsf A_{\mathbf j}[u,\kappa]
 =\prod_{i=1}^d(\mathsf A_{\mathbf j,i})_{u_i\kappa_i}
 \quad(\kappa\in\mathcal U_{\mathrm{coarse}}).
\]
Expanding the product of the affine identities gives
\[
r_u(x)=
 \sum_{\kappa\in\mathcal U_{\mathrm{coarse}}}
 \mathsf A_{\mathbf j}[u,\kappa]
 \prod_{i=1}^d
 \mathsf L_{\kappa_i}\!\left(\frac{x_i-c_{\mathbf j,i}}{\delta}\right).
\]
Only indices with \(\kappa_i\le u_i\) contribute, and hence every
contributing index has total degree at most two.

Now set
\[
v^{(e,u,J)}_{\mathbf j,\kappa}
 =\frac{e_P(c_{\mathbf j})\mathsf A_{\mathbf j}[u,\kappa]}
        {\sqrt{J^d}},
\qquad
f_{e,u,J}(x)=
 \sum_{\mathbf j,\kappa}
 b_{\mathbf j,\kappa}(x)v^{(e,u,J)}_{\mathbf j,\kappa}.
\]
Disjointness of the cells and the tensor identity imply the everywhere
defined formula
\[
f_{e,u,J}(x)=
\begin{cases}
e_P(c_{\mathbf j})r_u(x),&x\in D_{\mathbf j},\\
0,&x\notin C_h.
\end{cases}
\]
Define
\[
E_{e,J}=L(\sqrt d)^\alpha4^d\delta^\alpha\ge0.
\]
For \(x\in D_{\mathbf j}\),
\[
\begin{aligned}
|f_{e,u}(x)-f_{e,u,J}(x)|
 &=|e_P(x)-e_P(c_{\mathbf j})|\,|r_u(x)|\\
 &\le L\|x-c_{\mathbf j}\|_2^\alpha4^d\\
 &\le L(\sqrt d\,\delta)^\alpha4^d
 =E_{e,J}.
\end{aligned}
\]
The probability normalization therefore gives
\[
\|f_{e,u}-f_{e,u,J}\|_h^2\le E_{e,J}^2.
\]
Applying the projection comparison from the preceding step to this
fine-basis expansion yields
\[
\|e_Pr_u-\Pi_J(e_Pr_u)\|_h
 \le\|f_{e,u}-f_{e,u,J}\|_h
 \le E_{e,J}.
\]
Thus the positive constant
\[
C_{\mathrm e}=L(\sqrt d)^\alpha4^d+1
\]
satisfies, uniformly in \(u\),
\[
\|e_Pr_u-\Pi_J(e_Pr_u)\|_h
 \le C_{\mathrm e}\delta^\alpha.
\]
% lean: CausalSmith.Stat.TwosamplePointcateAnnotationFrontier.propensity_projection_bounds

\item \emph{A common residual constant.}
Define
\[
C_{\mathrm R}=\max\{C_{\mathrm e},C_{\mathrm m}\}>0.
\]
Since \(\delta^\alpha,\delta^\beta\ge0\), the preceding estimates imply
\[
\|e_Pr_u-\Pi_J(e_Pr_u)\|_h
 \le C_{\mathrm R}\delta^\alpha
 \quad(u\in\mathcal U_{\mathrm{coarse}}),
\qquad
\|\mu_{0,P}-\Pi_J\mu_{0,P}\|_h
 \le C_{\mathrm R}\delta^\beta.
\]
% lean: CausalSmith.Stat.TwosamplePointcateAnnotationFrontier.residual_projection_bounds

\item \emph{Combination of the bounds.}
Finally, choose
\[
C=\max\{C_{\mathrm T},\max\{C_{\mathrm B},C_{\mathrm R}\}\}.
\]
This is finite, positive, and depends only on the public parameters.
The inequalities
\[
C_{\mathrm T}\le C,\qquad
C_{\mathrm B}\le C,\qquad
C_{\mathrm R}\le C,
\qquad
h^\gamma\ge0,\qquad
\delta^\alpha\ge0,\qquad
\delta^\beta\ge0
\]
allow every bound above to be enlarged to the claimed constant \(C\).
The same Taylor coefficient vector retains its exact center value and
its equality with \(p\) on \(C_h\), while the squared-kernel integral
retains the exact value \(qJ^d=k\). These conclusions hold for every
admitted \(P\), every \(0<h\le1/2\), and every integer \(J\ge1\).
% lean: CausalSmith.Stat.TwosamplePointcateAnnotationFrontier.local_polynomial_projection_facts
\end{enumerate}
\end{proof}

The Taylor approximation controls variation of the contrast across the localization cube while reproducing its value exactly at the evaluation profile. The bounded coefficient vector and coarse basis provide uniform control of polynomial evaluation. For the fine projection, the integral bound controls its action on bounded approximation errors, whereas the squared-kernel identity records the sampling cost of its rank. The two nuisance residual bounds concern different functions: the propensity multiplied by a coarse basis coordinate, and the control mean. Their respective orders account for the combined nuisance approximation term in \cref{obj:thm:rectangular-upper}.

\subsection{Population positivity}

The corrected empirical matrix has a population counterpart whose stability is determined by overlap and orthogonal projection geometry. The next identity separates its quadratic form into an overlap-weighted polynomial norm and a squared projection residual.

\begin{lemmav}{lem:population-positivity}[Population quadratic form positivity]
Suppose the following conditions hold:
\begin{itemize}
\item \textbf{(Public parameters.)} The dimension \(d\) is an integer with \(d\ge1\), the smoothness parameters satisfy \(0<\alpha,\beta\le1\) and \(1\le\gamma\le3\), the radius satisfies \(L>1/2\), and the overlap parameter satisfies \(0<\varepsilon<1/2\).
\item \textbf{(Primitive law.)} The law \(P\) belongs to the primitive class \(\mathcal M\) at these public parameters, as specified in \cref{obj:def:primitive-class}. Let \(e_P\) be its designated admissible propensity witness.
\item \textbf{(Sample sizes.)} The sample counts \(n,m\) are nonnegative integers with \(n\ge2\).
\item \textbf{(Localization and projection.)} The localization scale satisfies \(0<h\le1/2\), and the grid resolution \(J\) is an integer with \(J\ge1\). Use the projection \(\Pi_J\) from \cref{obj:def:projection}.
\end{itemize}
Let \(\nu_h\) be the uniform localization measure on the localization cube \(C_h\), with
\[
d\nu_h(x)=h^{-d}\mathbf 1_{C_h}(x)\,dx.
\]
For every coefficient vector \(v\in\mathbb R^q\), indexed by the multi-indices of total degree at most two, define its coarse-basis polynomial and the population matrix by
\[
p_v(x)=\sum_u r_u(x)v_u,
\qquad
Q_{P,h,J}
=\mathbb E_{\mathsf E_{n,m}(P)}[\widehat Q_{h,J}],
\]
where \(r_u\) are the scaled orthonormal coarse polynomial basis functions and \(\widehat Q_{h,J}\) is the symmetrized corrected empirical matrix. Then
\[
\begin{aligned}
v^\top Q_{P,h,J}v
&=\int e_P(x)\bigl(1-e_P(x)\bigr)p_v(x)^2\,d\nu_h(x)\\
&\quad+\int
\bigl[e_P(x)p_v(x)-\Pi_J(e_Pp_v)(x)\bigr]^2\,d\nu_h(x)\\
&\ge \varepsilon(1-\varepsilon)\sum_u v_u^2.
\end{aligned}
\]
\end{lemmav}

\begin{proof}[Proof of \cref{obj:lem:population-positivity}]
Fix \(v\in\mathbb R^q\). We compute the population entries, contract them with \(v\), and then obtain the energy decomposition and lower bound.

\begin{enumerate}
\item The localization density and fine-basis coordinates are
\[
w_h(x)=h^{-d}\mathbf 1_{C_h}(x),
\qquad
k=qJ^d,
\qquad
b_{J,z_{\mathrm f}}(x)=[\mathbf b_J(x)]_{z_{\mathrm f}},
\quad z_{\mathrm f}\in\{1,\ldots,k\}.
\]
Here \(x\in\mathbb R^d\), and the density is as in
\cref{obj:synth_3}. The role counts from \cref{obj:def:roles} satisfy
\[
\ell=\lfloor n/2\rfloor\ge1,
\qquad
t=n+m-\ell\ge n-\ell\ge1.
\]
For fixed \(h,J\), the localized basis products are bounded and measurable, so all expectations and integrals below are finite.

For coarse indices \(u,w\in\{1,\ldots,q\}\), expand the kernel in
\cref{obj:def:projection} inside the raw matrix of
\cref{obj:def:moments}. This gives
\[
\begin{aligned}
\widehat Q^{\mathrm{raw}}_{h,J;uw}
&=\frac1\ell\sum_{j\in I_L}
 w_h(X_j)r_u(X_j)r_w(X_j)A_j\\
&\quad-\sum_{z_{\mathrm f}=1}^{k}\frac1{t\ell}
 \sum_{i\in I_T}\sum_{j\in I_L}
 \bigl[w_h(X_i^T)r_u(X_i^T)b_{J,z_{\mathrm f}}(X_i^T)A_i^T\bigr]
 \bigl[w_h(X_j)r_w(X_j)b_{J,z_{\mathrm f}}(X_j)A_j\bigr].
\end{aligned}
\]
Under the product sampling law of \cref{obj:ass:sample-law}, each covariate-treatment pair has marginal law \(P_{XA}\). Each treatment-role record is independent of each outcome-role record: for a labeled treatment-role record their indices are distinct, and for an auxiliary record independence follows from the product sampling law. Taking expectations term by term therefore yields
\[
\begin{aligned}
\mathbb E_{\mathsf E_{n,m}(P)}
 [\widehat Q^{\mathrm{raw}}_{h,J;uw}]
&=\frac{\ell}{\ell}\,
 \mathbb E_P[w_h(X)r_u(X)r_w(X)A]\\
&\quad-\sum_{z_{\mathrm f}=1}^{k}\frac{t\ell}{t\ell}\,
 \mathbb E_P[w_h(X)r_u(X)b_{J,z_{\mathrm f}}(X)A]\,
 \mathbb E_P[w_h(X)r_w(X)b_{J,z_{\mathrm f}}(X)A].
\end{aligned}
\]
For any of the functions
\[
f_{\mathrm{loc}}\in
\{r_ur_w,\ r_ub_{J,z_{\mathrm f}},\ r_wb_{J,z_{\mathrm f}}\},
\qquad
f_{\mathrm{loc}}:\mathbb R^d\longrightarrow\mathbb R,
\]
the designated conditional treatment mean and uniform design in
\cref{obj:def:primitive-class} give
\[
\begin{aligned}
\mathbb E_P[w_h(X)f_{\mathrm{loc}}(X)A]
&=\mathbb E_P[w_h(X)f_{\mathrm{loc}}(X)e_P(X)]\\
&=h^{-d}\int_{C_h}f_{\mathrm{loc}}(x)e_P(x)\,dx\\
&=\int f_{\mathrm{loc}}(x)e_P(x)\,d\nu_h(x).
\end{aligned}
\]
The localization cube lies in \([0,1]^d\), so these identities use the designated witness on its specified domain. Substitution gives a formula symmetric in \(u,w\). Consequently, the symmetrization in \cref{obj:def:moments} preserves its expectation, and
\[
\begin{aligned}
(Q_{P,h,J})_{uw}
&=\int r_u(x)e_P(x)r_w(x)\,d\nu_h(x)\\
&\quad-\sum_{z_{\mathrm f}=1}^{k}
 \left(\int b_{J,z_{\mathrm f}}(x)e_P(x)r_u(x)\,d\nu_h(x)\right)
 \left(\int b_{J,z_{\mathrm f}}(x)e_P(x)r_w(x)\,d\nu_h(x)\right).
\end{aligned}
\]
% lean: CausalSmith.Stat.TwosamplePointcateAnnotationFrontier.population_gram_entries

\item Define the weighted coarse-basis coefficients by
\[
c_{z_{\mathrm f},u}
=\int b_{J,z_{\mathrm f}}(x)e_P(x)r_u(x)\,d\nu_h(x),
\qquad
1\le z_{\mathrm f}\le k,\quad 1\le u\le q.
\]
Multiplying the entry formula by \(v_uv_w\), summing, and using
\(p_v(x)=\sum_{u=1}^{q}r_u(x)v_u\), we obtain
\[
\begin{aligned}
v^\top Q_{P,h,J}v
&=\int e_P(x)
 \left(\sum_{u=1}^{q}r_u(x)v_u\right)^2\,d\nu_h(x)
 -\sum_{z_{\mathrm f}=1}^{k}
 \left(\sum_{u=1}^{q}c_{z_{\mathrm f},u}v_u\right)^2\\
&=\int e_P(x)p_v(x)^2\,d\nu_h(x)
 -\sum_{z_{\mathrm f}=1}^{k}
 \left(\sum_{u=1}^{q}c_{z_{\mathrm f},u}v_u\right)^2.
\end{aligned}
\]
Indeed, the integral term follows by expanding the finite double sum, and the correction follows from
\[
\sum_{u,w=1}^{q}v_uv_w
 \sum_{z_{\mathrm f}=1}^{k}c_{z_{\mathrm f},u}c_{z_{\mathrm f},w}
=
\sum_{z_{\mathrm f}=1}^{k}
 \left(\sum_{u=1}^{q}c_{z_{\mathrm f},u}v_u\right)^2.
\]

The finite kernel expansion in \cref{obj:def:projection} and linearity of integration give
\[
\begin{aligned}
\Pi_J(e_Pp_v)(x)
&=\sum_{z_{\mathrm f}=1}^{k}b_{J,z_{\mathrm f}}(x)
 \int b_{J,z_{\mathrm f}}(y)e_P(y)p_v(y)\,d\nu_h(y)\\
&=\sum_{z_{\mathrm f}=1}^{k}b_{J,z_{\mathrm f}}(x)
 \sum_{u=1}^{q}c_{z_{\mathrm f},u}v_u.
\end{aligned}
\]
The fine basis in that definition is orthonormal. Expanding the square and integrating therefore gives
\[
\begin{aligned}
\int \bigl[\Pi_J(e_Pp_v)(x)\bigr]^2\,d\nu_h(x)
&=\sum_{z_{\mathrm f},z'_{\mathrm f}=1}^{k}
 \left(\sum_{u=1}^{q}c_{z_{\mathrm f},u}v_u\right)
 \left(\sum_{w=1}^{q}c_{z'_{\mathrm f},w}v_w\right)
 \int b_{J,z_{\mathrm f}}b_{J,z'_{\mathrm f}}\,d\nu_h\\
&=\sum_{z_{\mathrm f}=1}^{k}
 \left(\sum_{u=1}^{q}c_{z_{\mathrm f},u}v_u\right)^2.
\end{aligned}
\]
Thus
\[
v^\top Q_{P,h,J}v
=\int e_Pp_v^2\,d\nu_h
 -\int \bigl[\Pi_J(e_Pp_v)\bigr]^2\,d\nu_h.
\]
% lean: CausalSmith.Stat.TwosamplePointcateAnnotationFrontier.population_gram_quadratic_of_entries

\item The finite coarse expansion makes \(p_v\) square-integrable, and bounded overlap makes \(e_Pp_v\) square-integrable. The projection expansion just established also yields
\[
\begin{aligned}
\int e_Pp_v\,\Pi_J(e_Pp_v)\,d\nu_h
&=\sum_{z_{\mathrm f}=1}^{k}
 \left(\int b_{J,z_{\mathrm f}}e_Pp_v\,d\nu_h\right)^2\\
&=\int \bigl[\Pi_J(e_Pp_v)\bigr]^2\,d\nu_h.
\end{aligned}
\]
Expanding the residual square and using this equality gives the Pythagorean identity
\[
\int (e_Pp_v)^2\,d\nu_h
=
\int \bigl[\Pi_J(e_Pp_v)\bigr]^2\,d\nu_h
+\int \bigl[e_Pp_v-\Pi_J(e_Pp_v)\bigr]^2\,d\nu_h.
\]
Meanwhile, the pointwise identity
\(e_P(1-e_P)p_v^2=e_Pp_v^2-(e_Pp_v)^2\) gives
\[
\int e_P(1-e_P)p_v^2\,d\nu_h
=
\int e_Pp_v^2\,d\nu_h-\int(e_Pp_v)^2\,d\nu_h.
\]
Combining these two identities with the quadratic-form formula proves
\[
v^\top Q_{P,h,J}v
=
\int e_P(1-e_P)p_v^2\,d\nu_h
+\int\bigl[e_Pp_v-\Pi_J(e_Pp_v)\bigr]^2\,d\nu_h.
\]
% lean: CausalSmith.Stat.TwosamplePointcateAnnotationFrontier.population_projection_energy_identity

\item The scaled coarse polynomials of \cref{obj:synth_1} satisfy
\[
\int r_u(x)r_w(x)\,d\nu_h(x)
=
\begin{cases}
1,&u=w,\\
0,&u\ne w.
\end{cases}
\]
To see this, integrate their tensor products over the localization cube and rescale each coordinate to \([-1/2,1/2]\). Direct integration of the three normalized one-dimensional polynomials in that definition gives unit squared norms and zero pairwise inner products. When \(u=w\), every coordinate factor is one; when \(u\ne w\), some multi-index coordinate differs and its factor is zero. Hence finite expansion gives
\[
\int p_v(x)^2\,d\nu_h(x)
=\sum_{u,w=1}^{q}v_uv_w\int r_ur_w\,d\nu_h
=\sum_{u=1}^{q}v_u^2.
\]

Class membership supplies \cref{obj:ass:overlap}. Since \(\nu_h\) is supported on \(C_h\subseteq[0,1]^d\), for \(\nu_h\)-almost every \(x\),
\[
e_P(x)(1-e_P(x))-\varepsilon(1-\varepsilon)
=
\left(e_P(x)-\varepsilon\right)
\left(1-\varepsilon-e_P(x)\right)
\ge0.
\]
Multiplying by \(p_v(x)^2\) and integrating therefore yields
\[
\int e_P(1-e_P)p_v^2\,d\nu_h
\ge\varepsilon(1-\varepsilon)\int p_v^2\,d\nu_h
=\varepsilon(1-\varepsilon)\sum_{u=1}^{q}v_u^2.
\]
Finally,
\[
\int\bigl[e_Pp_v-\Pi_J(e_Pp_v)\bigr]^2\,d\nu_h\ge0,
\]
so the decomposition from the preceding step implies
\[
v^\top Q_{P,h,J}v\ge\varepsilon(1-\varepsilon)\sum_{u=1}^{q}v_u^2.
\]
% lean: CausalSmith.Stat.TwosamplePointcateAnnotationFrontier.population_expression_positivity
\end{enumerate}
\end{proof}

The overlap restriction in \cref{obj:ass:overlap} supplies the lower bound through the first integral. The second integral is nonnegative and makes the contribution of finite projection explicit. Thus the population matrix retains a fixed positivity margin at every admissible resolution. This margin supplies the population reference for the empirical eigenvalue guard in \cref{obj:def:estimator}.

\subsection{Rectangular sampling}

The public split assigns disjoint observations to the outcome and treatment roles under \cref{obj:ass:sample-law}. Their counts in \cref{obj:def:roles} are proportional to \(n\) and \(N=n+m\), respectively, including when \(m=0\). Consequently, the sampling analysis distinguishes ordinary localized averages from the rectangular averages formed across the two roles. The next result controls both the response vector and the matrix used for inversion.

\begin{lemmav}{lem:rectangular-sampling-bound}[Rectangular sampling bound]
Fix the covariate dimension \(d\in\mathbb N\), the smoothness parameters \(\alpha,\beta,\gamma\in\mathbb R\), the Hölder radius \(L\in\mathbb R\), and the overlap parameter \(\varepsilon\in\mathbb R\), satisfying
\begin{itemize}
\item \textbf{(Dimension.)} \(d\ge1\).
\item \textbf{(Nuisance smoothness.)} \(0<\alpha\le1\) and \(0<\beta\le1\).
\item \textbf{(Effect smoothness.)} \(1\le\gamma\le3\).
\item \textbf{(Radius.)} \(L>1/2\).
\item \textbf{(Overlap parameter.)} \(0<\varepsilon<1/2\).
\end{itemize}
There exists a finite constant \(C=C(d,\alpha,\beta,\gamma,L,\varepsilon)>0\) such that the following bound holds for every primitive law
\(P\in\mathcal M(d,\alpha,\beta,\gamma,L,\varepsilon)\) from \cref{obj:def:primitive-class}, and every \(n,m,J\in\mathbb N\) and \(h\in\mathbb R\) satisfying
\begin{itemize}
\item \textbf{(Sample sizes.)} \(n\ge2\) and \(m\ge0\).
\item \textbf{(Localization.)} \(0<h\le1/2\).
\item \textbf{(Resolution.)} \(J\ge1\).
\end{itemize}
Take expectations under \(\mathsf E_{n,m}(P)\), the product sampling law specified in \cref{obj:ass:sample-law}. Write \(N\) for the total treatment-record count, \(q\) for the number of polynomial multi-indices of total degree at most two, and \(k\) for the fine projection rank:
\[
N=n+m,\qquad q=\binom{d+2}{2},\qquad k=qJ^d.
\]
For the empirical response vector \(\widehat R_{h,J}\) and symmetrized corrected empirical matrix \(\widehat Q_{h,J}\) from \cref{obj:def:moments}, define their population expectations by
\[
\bar R=\mathbb E_{\mathsf E_{n,m}(P)}\widehat R_{h,J},
\qquad
Q_{P,h,J}=\mathbb E_{\mathsf E_{n,m}(P)}\widehat Q_{h,J}.
\]
Then
\[
\mathbb E_{\mathsf E_{n,m}(P)}
 \bigl\|\widehat R_{h,J}-\bar R\bigr\|_2^2
+
\mathbb E_{\mathsf E_{n,m}(P)}
 \bigl\|\widehat Q_{h,J}-Q_{P,h,J}\bigr\|_F^2
\le
C\left\{\frac{1}{nh^d}+\frac{k}{nNh^{2d}}\right\}.
\]
\end{lemmav}

\begin{proof}[Proof of \cref{obj:lem:rectangular-sampling-bound}]
\begin{enumerate}
\item Define the dimension-dependent constants
\[
M_d=4^{2d},\qquad
D_d=q4^{2d},\qquad
C_{\mathrm{ker},d}=M_d^4(D_d^2+1),
\]
\[
C_{\mathrm{coord},d}=6M_d^2+12C_{\mathrm{ker},d},
\qquad
C=(q+q^2+1)C_{\mathrm{coord},d}.
\]
These constants are finite and positive. Fix an admitted law \(P\) and admissible \(n,m,h,J\), and set
\[
V_{n,m,h,J}
=\frac{1}{nh^d}+\frac{k}{nNh^{2d}}\ge0.
\]
We will bound each response coordinate and each raw-matrix coordinate by
\(C_{\mathrm{coord},d}V_{n,m,h,J}\), then sum the bounds after symmetrization.
% lean: CausalSmith.Stat.TwosamplePointcateAnnotationFrontier.rectangular_sampling_bound

\item We first establish the explicit basis bounds used in the coordinate estimates. Write
\[
\mathcal U_d
=\left\{u\in\mathbb N^d:\sum_{i=1}^d u_i\le2\right\},
\qquad |\mathcal U_d|=q,
\]
and define the reference polynomial factors, for \(\xi_{\mathrm{ref}}\in\mathbb R\), by
\[
\begin{aligned}
p_{\mathrm{ref},0}(\xi_{\mathrm{ref}})&=1,\\
p_{\mathrm{ref},1}(\xi_{\mathrm{ref}})&=\sqrt{12}\,\xi_{\mathrm{ref}},\\
p_{\mathrm{ref},2}(\xi_{\mathrm{ref}})
&=\sqrt5\left(6\xi_{\mathrm{ref}}^2-\frac12\right).
\end{aligned}
\]
For \(|\xi_{\mathrm{ref}}|\le1/2\),
\[
0\le\xi_{\mathrm{ref}}^2\le\frac14,
\qquad
\left|6\xi_{\mathrm{ref}}^2-\frac12\right|\le1,
\]
and consequently
\[
|p_{\mathrm{ref},0}(\xi_{\mathrm{ref}})|=1\le4,\qquad
|p_{\mathrm{ref},1}(\xi_{\mathrm{ref}})|
\le\frac{\sqrt{12}}2\le4,\qquad
|p_{\mathrm{ref},2}(\xi_{\mathrm{ref}})|\le\sqrt5\le3\le4.
\]
The product formula in \cref{obj:synth_1} therefore gives
\[
|r_u(x)|\le4^d
\qquad(x\in C_h,\ u\in\mathcal U_d).
\]

For completeness, the same factors give an explicit row-integral bound for the fine kernel. Set
\[
\delta=\frac hJ,\qquad
\mathcal Z_J=\{0,\ldots,J-1\}^d,\qquad
x_{\mathrm{cell},\mathbf z,i}
=\frac12-\frac h2+\left(z_i+\frac12\right)\delta
\quad(\mathbf z\in\mathcal Z_J).
\]
Let \(C_{\mathrm{cell},\mathbf z}\) be the product of the intervals
\[
\left[\frac12-\frac h2+z_i\delta,\,
      \frac12-\frac h2+(z_i+1)\delta\right),
\]
with the right endpoint included when \(z_i=J-1\). Thus these cells partition \(C_h\).
The fine basis coordinates of \cref{obj:def:projection} are
\[
b_{\mathbf z,u}(x)
=\sqrt{J^d}\,\mathbf1_{C_{\mathrm{cell},\mathbf z}}(x)
 \prod_{i=1}^d
 p_{\mathrm{ref},u_i}
 \left(\frac{x_i-x_{\mathrm{cell},\mathbf z,i}}{\delta}\right),
\qquad
x\in\mathbb R^d,\quad
(\mathbf z,u)\in\mathcal Z_J\times\mathcal U_d.
\]
In particular, their values are zero outside their supporting cells, and
\[
|b_{\mathbf z,u}(x)|\le\sqrt{J^d}\,4^d.
\]
With \(\boldsymbol0=(0,\ldots,0)\), the constant coordinate satisfies
\[
b_{\mathbf z,\boldsymbol0}(x)
=\sqrt{J^d}\,\mathbf1_{C_{\mathrm{cell},\mathbf z}}(x),
\qquad
\int b_{\mathbf z,\boldsymbol0}(x)^2\,d\nu_h(x)=1,
\]
where the last identity is its orthonormal normalization. Pointwise,
\[
|b_{\mathbf z,u}(x)|
\le \frac{4^d}{\sqrt{J^d}}\,
      b_{\mathbf z,\boldsymbol0}(x)^2,
\]
so integration yields
\[
\int |b_{\mathbf z,u}(x)|\,d\nu_h(x)
\le\frac{4^d}{\sqrt{J^d}}.
\]
If \(x\in C_{\mathrm{cell},\mathbf z}\), only that cell contributes to the kernel, and hence
\[
K_J(x,x')
=\sum_{u\in\mathcal U_d}
 b_{\mathbf z,u}(x)b_{\mathbf z,u}(x').
\]
It follows that
\[
\begin{aligned}
\int |K_J(x,x')|\,d\nu_h(x')
&\le\sum_{u\in\mathcal U_d}
 |b_{\mathbf z,u}(x)|
 \int |b_{\mathbf z,u}(x')|\,d\nu_h(x')\\
&\le q\bigl(\sqrt{J^d}\,4^d\bigr)
          \frac{4^d}{\sqrt{J^d}}
=D_d.
\end{aligned}
\]
If \(x\) belongs to no fine cell, every fine basis coordinate at \(x\) is zero, and the same bound holds. Thus
\[
\int |K_J(x,x')|\,d\nu_h(x')\le D_d
\qquad(x\in\mathbb R^d).
\]
The basis-product formula also gives
\[
K_J(x,x')=K_J(x',x).
\]
% lean: CausalSmith.Stat.TwosamplePointcateAnnotationFrontier.fineKernel_abs_row_bound

\item We now express the two coordinate types through a common marked kernel. Define the role counts
\[
\ell=\lfloor n/2\rfloor,\qquad t=N-\ell.
\]
Both are positive because \(n\ge2\) and \(N\ge n\). Let \(W\) and \(Z\) be independent generic treatment-role and outcome-role records, respectively:
\[
W=(X^T,A^T)\in[0,1]^d\times\{0,1\},
\qquad
Z=(X^L,A^L,Y^L)\in[0,1]^d\times\{0,1\}^2,
\]
\[
\mathcal L(W)=P_{XA},
\qquad
\mathcal L(Z)=\mathcal L_P(X,A,Y).
\]
Their covariate marginals are uniform by \cref{obj:def:primitive-class}. The product sampling law and the disjoint roles in \cref{obj:def:roles} give independent records
\[
W_i=(X_i^T,A_i^T)\quad(i\in I_T),
\qquad
Z_j=(X_j,A_j,Y_j)\quad(j\in I_L),
\]
independent within and between the two roles.

Use the coordinate tags
\[
\diamond\in
\bigl(\{\mathrm R\}\times\mathcal U_d\bigr)
\sqcup
\bigl(\{\mathrm Q\}\times\mathcal U_d\times\mathcal U_d\bigr).
\]
For each tag, define the three mark functions by
\[
\begin{array}{c|ccc}
\diamond
& f_{\diamond,0}(Z)
& f_{\diamond,T}(W)
& f_{\diamond,L}(Z)\\ \hline
(\mathrm R,u)
& r_u(X^L)A^LY^L
& r_u(X^T)A^T
& Y^L\\
(\mathrm Q,u,v)
& r_u(X^L)r_v(X^L)A^L
& r_u(X^T)A^T
& r_v(X^L)A^L
\end{array}
\]
on the full record domains displayed above. Binary marks have absolute value at most one. The coarse-coordinate bound from the preceding step therefore gives
\[
\begin{aligned}
|f_{\diamond,0}(Z)|&\le M_d &&(X^L\in C_h),\\
|f_{\diamond,T}(W)|&\le M_d &&(X^T\in C_h),\\
|f_{\diamond,L}(Z)|&\le M_d &&(X^L\in C_h).
\end{aligned}
\]
Define
\[
H_\diamond(W,Z)
=w_h(X^T)w_h(X^L)K_J(X^T,X^L)
 f_{\diamond,T}(W)f_{\diamond,L}(Z),
\]
\[
\widehat S_\diamond
=\frac1\ell\sum_{j\in I_L}
 w_h(X_j)f_{\diamond,0}(Z_j),
\qquad
\widehat H_\diamond
=\frac1{t\ell}\sum_{i\in I_T}\sum_{j\in I_L}
 H_\diamond(W_i,Z_j),
\qquad
F_\diamond=\widehat S_\diamond-\widehat H_\diamond.
\]
These expressions vanish termwise whenever a localization weight is zero. The rectangular formulas in \cref{obj:def:moments} give exactly
\[
F_{(\mathrm R,u)}=(\widehat R_{h,J})_u,
\qquad
F_{(\mathrm Q,u,v)}
=(\widehat Q^{\mathrm{raw}}_{h,J})_{uv}.
\]
In the response identity, expanding the empirical projection gives the rectangular formula on \(C_h\); outside \(C_h\), its defining zero value and the zero localization weight give the same identity.
% lean: CausalSmith.Stat.TwosamplePointcateAnnotationFrontier.rectangular_r_coordinate_bound

\item The localized mark and kernel energies admit explicit bounds. Since
\[
w_h(x)=h^{-d}\mathbf1_{C_h}(x),
\qquad
\mathbb P(X^L\in C_h)=\mathbb P(X^T\in C_h)=h^d,
\]
uniform design gives
\[
\mathbb E\,w_h(X^L)^2
=\mathbb E\,w_h(X^T)^2=h^{-d}.
\]
Thus
\[
\mathbb E\!\left[
 \bigl(w_h(X^L)f_{\diamond,0}(Z)\bigr)^2
 \right]\le \frac{M_d^2}{h^d}.
\]
Similarly, bounding both marks and integrating over the two independent covariates gives
\[
\begin{aligned}
\mathbb E\,H_\diamond(W,Z)^2
&\le M_d^4h^{-4d}
 \iint_{C_h\times C_h}K_J(x,x')^2\,dx\,dx'\\
&=\frac{M_d^4}{h^{2d}}
 \iint K_J(x,x')^2\,d\nu_h(x)\,d\nu_h(x')\\
&=\frac{M_d^4k}{h^{2d}},
\end{aligned}
\]
using the kernel-square identity in
\cref{obj:lem:local-polynomial-projection-facts}.

Define the integrated slices on the full record domains by
\[
H_{\diamond,T}(W)=\mathbb E_Z H_\diamond(W,Z),
\qquad
H_{\diamond,L}(Z)=\mathbb E_W H_\diamond(W,Z),
\]
where each expectation integrates an independent record with the law specified above. The row-integral estimate gives
\[
\begin{aligned}
\left|
 \mathbb E_Z\!\left[
 w_h(X^L)K_J(x,X^L)f_{\diamond,L}(Z)
 \right]\right|
&\le M_d\,\mathbb E_Z\!
 \left[w_h(X^L)|K_J(x,X^L)|\right]\\
&=M_d\int |K_J(x,x')|\,d\nu_h(x')\\
&\le M_dD_d.
\end{aligned}
\]
Consequently,
\[
|H_{\diamond,T}(W)|
\le M_d^2D_d\,w_h(X^T).
\]
Kernel symmetry gives the corresponding bound with the roles reversed:
\[
|H_{\diamond,L}(Z)|
\le M_d^2D_d\,w_h(X^L).
\]
Squaring and integrating yields
\[
\mathbb E\,H_{\diamond,T}(W)^2
\le\frac{M_d^4D_d^2}{h^d},
\qquad
\mathbb E\,H_{\diamond,L}(Z)^2
\le\frac{M_d^4D_d^2}{h^d}.
\]
All the displayed features and kernels are measurable and square-integrable. In particular, their means, centered squares, and finite-sum integrals are well defined.
% lean: CausalSmith.Stat.TwosamplePointcateAnnotationFrontier.rectangular_marked_kernel_row_energy

\item We derive the rectangular variance formula used with these energies. Define
\[
\begin{aligned}
\mu_\diamond&=\mathbb E\,H_\diamond(W,Z),\\
H_{\diamond,0}(W,Z)&=H_\diamond(W,Z)-\mu_\diamond,\\
H_{\diamond,1}(W)&=H_{\diamond,T}(W)-\mu_\diamond,\\
H_{\diamond,2}(Z)&=H_{\diamond,L}(Z)-\mu_\diamond,\\
H_{\diamond,12}(W,Z)
&=H_{\diamond,0}(W,Z)-H_{\diamond,1}(W)-H_{\diamond,2}(Z).
\end{aligned}
\]
Fubini's theorem gives
\[
\mathbb E H_{\diamond,0}
=\mathbb E H_{\diamond,1}
=\mathbb E H_{\diamond,2}=0,
\]
and
\[
\begin{aligned}
\mathbb E\!\left[H_{\diamond,0}H_{\diamond,1}(W)\right]
&=\mathbb E H_{\diamond,1}(W)^2,\\
\mathbb E\!\left[H_{\diamond,0}H_{\diamond,2}(Z)\right]
&=\mathbb E H_{\diamond,2}(Z)^2,\\
\mathbb E\!\left[H_{\diamond,1}(W)H_{\diamond,2}(Z)\right]
&=0.
\end{aligned}
\]
Indeed, the first two identities follow by integrating the other record first, and the third follows from independence and centering. Expanding the residual square therefore gives
\[
\mathbb E H_{\diamond,12}^2
=\mathbb E H_{\diamond,0}^2
 -\mathbb E H_{\diamond,1}^2
 -\mathbb E H_{\diamond,2}^2.
\]
Also,
\[
\mathbb E H_{\diamond,0}^2
=\mathbb E H_\diamond^2-\mu_\diamond^2
\le\mathbb E H_\diamond^2.
\]
Hence
\[
\mathbb E H_{\diamond,12}^2\le\mathbb E H_\diamond^2,
\]
and centering each slice gives
\[
\mathbb E H_{\diamond,1}^2
\le\mathbb E H_{\diamond,T}^2,
\qquad
\mathbb E H_{\diamond,2}^2
\le\mathbb E H_{\diamond,L}^2.
\]

To compute the rectangular second moment, recall the independent-sample identity. For \(n_{\mathrm{av}}\ge1\) independent identically distributed square-integrable real variables \(Z_{\mathrm{av},j}\), define
\[
\overline Z_{\mathrm{av}}
=\frac1{n_{\mathrm{av}}}
 \sum_{j=1}^{n_{\mathrm{av}}}Z_{\mathrm{av},j}.
\]
Expanding the square and factoring the mixed expectations gives
\[
\mathbb E\,\overline Z_{\mathrm{av}}^2
=\frac1{n_{\mathrm{av}}}\mathbb E Z_{\mathrm{av},1}^2
 +\left(1-\frac1{n_{\mathrm{av}}}\right)
   (\mathbb E Z_{\mathrm{av},1})^2.
\]
Apply this identity first to the outcome-record average, with the treatment records fixed, and then to the treatment-record averages. Finite-sum integration commutes with taking the slice means. Since \(H_{\diamond,0}\) has mean zero, the resulting identity is
\[
\begin{aligned}
\mathbb E(\widehat H_\diamond-\mu_\diamond)^2
&=\frac{\mathbb E H_{\diamond,0}^2}{t\ell}\\
&\quad+\frac1t\left(1-\frac1\ell\right)
       \mathbb E H_{\diamond,1}^2\\
&\quad+\left(1-\frac1t\right)\frac1\ell
       \mathbb E H_{\diamond,2}^2.
\end{aligned}
\]
Substituting the residual-energy identity above reduces this to
\[
\mathbb E(\widehat H_\diamond-\mu_\diamond)^2
=\frac{\mathbb E H_{\diamond,1}^2}{t}
 +\frac{\mathbb E H_{\diamond,2}^2}{\ell}
 +\frac{\mathbb E H_{\diamond,12}^2}{t\ell}.
\]
The three preceding energy comparisons imply
\[
\mathbb E(\widehat H_\diamond-\mu_\diamond)^2
\le
\frac{\mathbb E H_{\diamond,T}^2}{t}
+\frac{\mathbb E H_{\diamond,L}^2}{\ell}
+\frac{\mathbb E H_\diamond^2}{t\ell}.
\]
% lean: CausalSmith.Stat.TwosamplePointcateAnnotationFrontier.rectangular_ustat_variance

\item We convert this inequality to the sampling envelope and then include the first-order term. The floor split satisfies
\[
\ell\ge1,\qquad n\le2\ell+1\le3\ell,
\qquad 2\ell\le n\le N.
\]
Therefore
\[
\ell\ge\frac n3,\qquad
t=N-\ell\ge\frac N2,
\]
and
\[
\frac1\ell\le\frac3n,\qquad
\frac1t\le\frac2N\le\frac2n,\qquad
\frac1{t\ell}\le\frac6{nN}.
\]
These inequalities apply for every \(m\ge0\), including \(m=0\).

By the definition of \(C_{\mathrm{ker},d}\),
\[
M_d^4D_d^2\le C_{\mathrm{ker},d},
\qquad
M_d^4\le C_{\mathrm{ker},d}.
\]
The slice and kernel bounds thus give
\[
\begin{aligned}
\mathbb E(\widehat H_\diamond-\mu_\diamond)^2
&\le \frac{5C_{\mathrm{ker},d}}{nh^d}
     +\frac{6C_{\mathrm{ker},d}k}{nNh^{2d}}\\
&\le6C_{\mathrm{ker},d}V_{n,m,h,J}.
\end{aligned}
\]
The independent-sample mean identity also gives
\[
\mathbb E\widehat S_\diamond
=\mathbb E\!\left[w_h(X^L)f_{\diamond,0}(Z)\right],
\qquad
\mathbb E\widehat H_\diamond=\mu_\diamond,
\]
and
\[
\begin{aligned}
\mathbb E(\widehat S_\diamond-\mathbb E\widehat S_\diamond)^2
&=\frac1\ell
 \left\{
 \mathbb E\!\left[
   \bigl(w_h(X^L)f_{\diamond,0}(Z)\bigr)^2
 \right]
 -\left(\mathbb E[w_h(X^L)f_{\diamond,0}(Z)]\right)^2
 \right\}\\
&\le\frac{3M_d^2}{nh^d}.
\end{aligned}
\]
Linearity of expectation gives the centered-difference identity
\[
F_\diamond-\mathbb EF_\diamond
=(\widehat S_\diamond-\mathbb E\widehat S_\diamond)
 -(\widehat H_\diamond-\mu_\diamond).
\]
The pointwise inequality
\[
(\xi_{\mathrm{diff}}-\zeta_{\mathrm{diff}})^2
\le2\xi_{\mathrm{diff}}^2+2\zeta_{\mathrm{diff}}^2
\qquad(\xi_{\mathrm{diff}},\zeta_{\mathrm{diff}}\in\mathbb R)
\]
follows from \((\xi_{\mathrm{diff}}+\zeta_{\mathrm{diff}})^2\ge0\). Applying it to the centered difference yields
\[
\begin{aligned}
\mathbb E(F_\diamond-\mathbb EF_\diamond)^2
&\le\frac{6M_d^2}{nh^d}
  +12C_{\mathrm{ker},d}V_{n,m,h,J}\\
&\le C_{\mathrm{coord},d}V_{n,m,h,J},
\end{aligned}
\]
because \(1/(nh^d)\le V_{n,m,h,J}\). This argument accommodates the shared outcome records in the two terms.
% lean: CausalSmith.Stat.TwosamplePointcateAnnotationFrontier.rectangular_marked_moment_bound

\item Applying the preceding estimate to the response tags gives
\[
\mathbb E\left[
 \bigl((\widehat R_{h,J})_u-\bar R_u\bigr)^2
 \right]
\le C_{\mathrm{coord},d}V_{n,m,h,J}
\qquad(u\in\mathcal U_d).
\]
Applying it to the raw-matrix tags gives
\[
\mathbb E\left[
 \bigl((\widehat Q^{\mathrm{raw}}_{h,J})_{uv}
       -\mathbb E(\widehat Q^{\mathrm{raw}}_{h,J})_{uv}\bigr)^2
 \right]
\le C_{\mathrm{coord},d}V_{n,m,h,J}
\qquad(u,v\in\mathcal U_d).
\]
The square-integrability established above applies to every coordinate in these two displays.
% lean: CausalSmith.Stat.TwosamplePointcateAnnotationFrontier.rectangular_q_coordinate_bound

\item Define the centered raw matrix by
\[
B_{\mathrm{raw},uv}
=(\widehat Q^{\mathrm{raw}}_{h,J})_{uv}
 -\mathbb E(\widehat Q^{\mathrm{raw}}_{h,J})_{uv},
\qquad u,v\in\mathcal U_d.
\]
The symmetrization formula in \cref{obj:def:moments} and linearity of expectation give
\[
(Q_{P,h,J})_{uv}
=\frac{
 \mathbb E(\widehat Q^{\mathrm{raw}}_{h,J})_{uv}
 +\mathbb E(\widehat Q^{\mathrm{raw}}_{h,J})_{vu}}2,
\]
and therefore
\[
(\widehat Q_{h,J}-Q_{P,h,J})_{uv}
=\frac{B_{\mathrm{raw},uv}+B_{\mathrm{raw},vu}}2.
\]
Every symmetrized coordinate is square-integrable as an average of two square-integrable raw coordinates.

For each sample,
\[
\left(\frac{B_{\mathrm{raw},uv}+B_{\mathrm{raw},vu}}2\right)^2
\le\frac{B_{\mathrm{raw},uv}^2+B_{\mathrm{raw},vu}^2}{2},
\]
since the difference between the right and left sides is
\((B_{\mathrm{raw},uv}-B_{\mathrm{raw},vu})^2/4\). Summing and exchanging the two indices gives the pointwise contraction
\[
\begin{aligned}
\|\widehat Q_{h,J}-Q_{P,h,J}\|_F^2
&\le\frac12\sum_{u,v\in\mathcal U_d}
 \left(B_{\mathrm{raw},uv}^2+B_{\mathrm{raw},vu}^2\right)\\
&=\sum_{u,v\in\mathcal U_d}B_{\mathrm{raw},uv}^2.
\end{aligned}
\]
% lean: CausalSmith.Stat.TwosamplePointcateAnnotationFrontier.rectangular_symmetrization_energy_le

\item Integrability permits exchanging expectation and each finite sum. Summing the response-coordinate bounds gives
\[
\begin{aligned}
\mathbb E\|\widehat R_{h,J}-\bar R\|_2^2
&=\sum_{u\in\mathcal U_d}
 \mathbb E\left[
   \bigl((\widehat R_{h,J})_u-\bar R_u\bigr)^2
 \right]\\
&\le qC_{\mathrm{coord},d}V_{n,m,h,J}.
\end{aligned}
\]
Integrating the pointwise matrix contraction and summing the raw-coordinate bounds gives
\[
\begin{aligned}
\mathbb E\|\widehat Q_{h,J}-Q_{P,h,J}\|_F^2
&\le\sum_{u,v\in\mathcal U_d}
 \mathbb E B_{\mathrm{raw},uv}^2\\
&\le q^2C_{\mathrm{coord},d}V_{n,m,h,J}.
\end{aligned}
\]
Adding these inequalities, and using
\(C_{\mathrm{coord},d}V_{n,m,h,J}\ge0\), yields
\[
\begin{aligned}
&\mathbb E\|\widehat R_{h,J}-\bar R\|_2^2
+\mathbb E\|\widehat Q_{h,J}-Q_{P,h,J}\|_F^2\\
&\qquad\le(q+q^2)C_{\mathrm{coord},d}V_{n,m,h,J}\\
&\qquad\le(q+q^2+1)C_{\mathrm{coord},d}V_{n,m,h,J}\\
&\qquad=C\left\{\frac1{nh^d}+\frac{k}{nNh^{2d}}\right\}.
\end{aligned}
\]
The chosen constant depends on \(d\), and hence has the permitted dependence on the fixed public parameters. All estimates hold uniformly over the admitted law and the admissible sample sizes and tuning inputs.
% lean: CausalSmith.Stat.TwosamplePointcateAnnotationFrontier.rectangular_sampling_bound
\end{enumerate}
\end{proof}

The first term reflects labeled information within the localization cube. The second couples the two record counts and increases with the fine projection rank. Because \(\delta=h/J\), its square-root order is the rectangular sampling term in \cref{obj:thm:rectangular-upper}. The same bound covers fluctuations of the response moments and of the empirical matrix, so it supplies the sampling control needed for the guarded local fit.

\subsection{Guarded inversion and public tuning}

The decision in \cref{obj:def:estimator} uses the empirical cutoff \(g(\varepsilon)=\varepsilon(1-\varepsilon)/2\) against the population positivity margin \(\varepsilon(1-\varepsilon)\) in \cref{obj:lem:population-positivity}. This half-margin gap connects matrix fluctuations to the stability of the inversion branch. The complementary branch returns zero, and clipping confines the decision to the contrast range \([-1,1]\) supplied by \cref{obj:ass:control-interior,obj:ass:treated-interior}. These operations give the total, finite, Borel decision covered by \cref{obj:thm:rectangular-upper}; the resulting constants may depend on the fixed overlap level.

The roles of the auxiliary results are distinct. \Cref{obj:lem:local-polynomial-projection-facts} supplies the effect approximation order \(h^\gamma\), the nuisance projection orders, and bounded polynomial coefficients. \Cref{obj:lem:population-positivity} supplies population stability, while \cref{obj:lem:rectangular-sampling-bound} controls empirical departures from the population moments. Together, these ingredients underlie the four contributions in the uniform risk bound of \cref{obj:thm:rectangular-upper}: effect approximation, combined nuisance approximation, labeled sampling, and rectangular sampling. Clipping supplies the constant bound represented by the minimum with one.

The public tuning in \cref{obj:def:tuning} assigns the fine resolution according to \(S=\alpha+\beta\), the sum of nuisance smoothness orders. When \(S<\gamma\), the resolution increases as the localization scale decreases, placing nuisance approximation at the effect approximation order. When \(S\ge\gamma\), the choice \(J_*=1\) provides the required approximation control. The ceiling preserves an integer resolution while changing its order by a bounded factor.

The localization scale balances the labeled sampling requirement and the rectangular sampling requirement. Its two candidate orders correspond to the two terms in \cref{obj:def:sharp-rate}, with denominator \(\Delta=2\gamma+d+\gamma d/S\) for the unequal-count term. Taking their maximum yields the publicly tuned risk guarantee in \cref{obj:thm:rectangular-upper}. The resulting precision has a direct statistical interpretation: response labels govern the ordinary local sampling term, while the total treatment-record count reduces the sampling cost of the fine projection correction.

\subsection{The supplied-propensity upper bound}

The decision displayed in \cref{obj:thm:supplied-propensity} uses the designated propensity from \cref{obj:def:primitive-class} to form an inverse-propensity-weighted response. Under exchangeability and overlap, this response has conditional mean equal to the canonical contrast and bounded magnitude. The known uniform covariate distribution then supplies the design measure for the degree-two local polynomial average.

Polynomial reproduction makes the contrast smoothness the relevant approximation order for this decision. Its bandwidth \(n^{-1/(2\gamma+d)}\), specified in \cref{obj:thm:supplied-propensity}, balances local approximation with labeled sampling. Clipping preserves the bounded contrast range. These are the statistical ingredients of the upper-bound component of that result, which attains the order \(r_o(n)=n^{-\gamma/(2\gamma+d)}\) uniformly over the primitive class and every auxiliary-record count. This benchmark identifies the labeled-sample precision attained by the rectangular estimator when its combined approximation and sampling costs are controlled at the same order.


\section{Original-record lower bounds}\label{sec:original-record-lower-bounds}

The lower bounds use a \leanref{S-4}{marked-component lower-bound experiment} within the primitive class of \cref{obj:def:primitive-class}. Under each hypothesis, a finite prior on primitive laws induces a mixture of the original sampling experiment: each primitive law generates exactly \(n\) labeled records and \(m\) outcome-free treatment records before mixing. Write \leanref{sym:Mtheta}{\ensuremath{M_\theta}} for this fixed-size mixture law under hypothesis \leanref{sym:theta}{\ensuremath{\theta\in\{-1,+1\}}}. The construction separates the point contrasts while controlling the statistical distance between the two complete-data mixtures. We use the following normalization of that distance.

\begin{definitionv}{synth_4}[Treatment-record marginal]
For a primitive Borel probability law
\(P=\mathcal L(X,A,Y(0),Y(1))\) on
\([0,1]^d\times\{0,1\}^3\), define its treatment-record marginal by
\[
P_{XA}=\mathcal L_P(X,A).
\]
Thus \(P_{XA}\) is the probability law on \([0,1]^d\times\{0,1\}\) satisfying
\[
P_{XA}(E)=P\bigl((X,A)\in E\bigr)
\]
for every Borel set \(E\subseteq[0,1]^d\times\{0,1\}\).
\end{definitionv}

\begin{definitionv}{synth_6}[Fixed-size hypothesis mixtures]
For \(n\in\mathbb N\) with \(n\ge2\), \(m\in\mathbb N_0\), and the finite priors \(\varpi_\theta\) produced by \cref{obj:def:marked-handle} on its probability-valued domain, define
\[
M_\theta
:=\int\bigl(P_o^{\otimes n}\otimes P_{XA}^{\otimes m}\bigr)\,d\varpi_\theta(P),
\qquad \theta\in\{-1,+1\},
\]
where \(P_o=\mathcal L_P(X,A,Y)\), \(Y=(1-A)Y(0)+AY(1)\), and \(P_{XA}=\mathcal L_P(X,A)\). Thus \(M_{+1}\) and \(M_{-1}\) are the positive- and negative-hypothesis probability laws of the original-record dataset \(\mathcal D_{n,m}\), with \(n\) labeled records and \(m\) outcome-free treatment records. Under each mixture, one primitive law \(P\) is drawn from the corresponding prior and all \(n+m\) records are drawn conditionally independently from that law through their respective observation channels.
\end{definitionv}

The unhalved squared Hellinger distance \leanref{sym:Htwo}{\ensuremath{H^2(M_{+1},M_{-1})}} measures the distinguishability of the two hypotheses using both record types. The finite-prior construction \leanref{sym:lowerHandle}{\ensuremath{\mathfrak C}} of \cref{obj:def:marked-handle} below couples propensity and response perturbations through correlated sign pairs. A smooth localization bump controls the contrast near the evaluation profile, while a finer smooth frame controls the spatial dependence of the perturbations.

\begin{definitionv}{def:hellinger}[Squared Hellinger distance \(H^2\)]
The squared Hellinger distance between the positive- and negative-hypothesis mixture laws \(M_{+1}\) and \(M_{-1}\) is
\[
H^2(M_{+1},M_{-1})
=
\int
\left(
\sqrt{\frac{dM_{+1}}{d(M_{+1}+M_{-1})}}
-
\sqrt{\frac{dM_{-1}}{d(M_{+1}+M_{-1})}}
\right)^2
\,d(M_{+1}+M_{-1}).
\]
\end{definitionv}

\begin{definitionv}{synth_2}[Constructed arm-mean functions]
The real-valued arm-mean functions \(\mu_{0,\theta,\eta}\) and \(\mu_{1,\theta,\eta}\) are defined for every dimension \(d\in\mathbb N\), real scales \(h,\delta\), real propensity perturbation amplitude \(a\), real response perturbation amplitude \(b\), hypothesis sign \(\theta\in\{-1,+1\}\), and response sign array \(\eta\in\{-1,+1\}^{\mathcal I}\) by
\[
\begin{aligned}
\mu_{0,\theta,\eta}(x)
&=\frac12+b\sum_{\mathbf z\in\mathcal I}
  \eta_{\mathbf z}\phi_{\mathbf z}(x)
  +\theta ab\,B_h(x)^2,\\
\mu_{1,\theta,\eta}(x)
&=\frac12+b\sum_{\mathbf z\in\mathcal I}
  \eta_{\mathbf z}\phi_{\mathbf z}(x)
  -\theta ab\,B_h(x)^2,
\qquad x\in\mathbb R^d.
\end{aligned}
\]
Here \(\mathcal I\), the finite frame index set, \(B_h\), the localization bump, and \(\phi_{\mathbf z}\), the localized frame entries, use the defining formulas in \cref{obj:def:marked-handle}, evaluated at dimension \(d\) and the stated real parameter values. The hypothesis sign and response signs encode Boolean values by assigning \(+1\) to true and \(-1\) to false. Each \(\eta_{\mathbf z}\) is the encoded second component of the Boolean sign pair at \(\mathbf z\); the first component is arbitrary. Control and treatment correspond to false and true arm values, respectively.
\end{definitionv}



In \cref{obj:def:marked-handle}, \leanref{sym:a}{\ensuremath{a}} is the propensity perturbation amplitude and \(b\) is the scalar response-mean perturbation amplitude. The arm means are centered at the shared function \(1/2+bG_\eta(x)\), with offsets \(-\tau_\theta(x)/2\) under control and \(+\tau_\theta(x)/2\) under treatment. Consequently, the shared perturbation cancels from their difference, leaving exactly \(\tau_\theta(x)\). Here the scalar \(b\) is distinct from the estimator's fine polynomial basis vector \(\mathbf b_J(x)\) in \cref{obj:def:projection}.

In \cref{obj:def:marked-handle}, the product sign-pair weights \leanref{sym:Pi}{\ensuremath{w_\theta(\lambda,\eta)}} encode the hypothesis through the correlation between the two perturbations. The construction itself depends only on \(d\) and \((h,\delta,a,b)\); the parameters \(\alpha,\beta,\gamma,L,\varepsilon\) enter only through the admissibility, occupancy, and information constants of the construction certificate. The deterministic contrast \leanref{sym:tautheta}{\ensuremath{\tau_\theta}} is shared by every primitive law within a hypothesis. Thus mixing over nuisance functions preserves an exact target separation. The compensating terms in the arm means link this separation to the singleton cancellation stated in the construction certificate below.

The information argument uses conditional factorization over spatial components. The calculus in \cref{obj:lem:hellinger-block-calculus}, stated after the construction, records how squared Hellinger distance behaves under products, conditioning on a common marginal, and the addition of a common independent coordinate.

\begin{algorithmv}{def:marked-handle}[Finite-prior construction \(\mathfrak C\)]
The construction \(\mathfrak C\) takes inputs \((h,\delta,a,b)\) satisfying \(0<\delta\le h\le1/2\) and \(a,b>0\), at the public parameters \((d,\alpha,\beta,\gamma,L)\) and fixed interior covariate profile \(x_0\).
\begin{enumerate}
\item Define the real functions
\[
\chi(u)=
\begin{cases}
\exp(-1/u),&u>0,\\
0,&u\le0,
\end{cases}
\qquad
\vartheta(u)=\frac{\pi}{2}\frac{\chi(u)}{\chi(u)+\chi(1-u)},
\qquad u\in\mathbb R,
\]
and, for each \(z\in\mathbb Z\),
\[
\omega_z(u)=
\begin{cases}
\sin\vartheta(u-z+1),&z-1\le u\le z,\\
\cos\vartheta(u-z),&z\le u\le z+1,\\
0,&u\notin[z-1,z+1].
\end{cases}
\]
The two expressions at \(u=z\) agree.

\item For \(w\in\mathbb R^d\) and \(x\in[0,1]^d\), define the localization bump and fine-scale frame by
\[
B(w)=\prod_{v=1}^d\frac{\chi(1-4w_v^2)}{\chi(1)},
\qquad
B_h(x)=B\bigl((x-x_0)/h\bigr),
\]
\[
\psi_{\mathbf z}(x)=
\prod_{v=1}^d
\omega_{z_v}\bigl((x_v-x_{0,v})/\delta\bigr),
\]
\[
\mathcal I=
\left\{\mathbf z\in\mathbb Z^d:
|z_v|\le1+\frac{h}{2\delta}
\text{ for every }v\in\{1,\ldots,d\}\right\},
\qquad
\phi_{\mathbf z}(x)=B_h(x)\psi_{\mathbf z}(x),
\quad \mathbf z\in\mathcal I.
\]

\item For sign arrays \(\lambda,\eta\in\{-1,+1\}^{\mathcal I}\), define
\[
F_\lambda(x)=
\sum_{\mathbf z\in\mathcal I}\lambda_{\mathbf z}\phi_{\mathbf z}(x),
\qquad
G_\eta(x)=
\sum_{\mathbf z\in\mathcal I}\eta_{\mathbf z}\phi_{\mathbf z}(x).
\]
For each hypothesis index \(\theta\in\{-1,+1\}\), define the sign-pair correlation and product weights by
\[
\rho_\theta=\frac{\theta}{2},
\qquad
w_\theta(\lambda,\eta)=
\prod_{\mathbf z\in\mathcal I}
\frac{1+\rho_\theta\lambda_{\mathbf z}\eta_{\mathbf z}}{4}.
\]

\item Define the propensity, contrast, and potential-response means by
\[
e_\lambda(x)=\frac12+aF_\lambda(x),
\qquad
\tau_\theta(x)=-4ab\rho_\theta B_h(x)^2,
\]
\[
\mu_{0,\theta,\eta}(x)=
\frac12+bG_\eta(x)-\frac12\tau_\theta(x),
\qquad
\mu_{1,\theta,\eta}(x)=
\frac12+bG_\eta(x)+\frac12\tau_\theta(x).
\]
The probability-valued domain consists of inputs satisfying
\[
e_\lambda(x),\ \mu_{0,\theta,\eta}(x),\
\mu_{1,\theta,\eta}(x)\in[0,1]
\quad
\text{for every }\theta\in\{-1,+1\},\
\lambda,\eta\in\{-1,+1\}^{\mathcal I},\
x\in[0,1]^d.
\]
On this domain, define the primitive Borel law \(P_{\theta,\lambda,\eta}\) by
\[
X\sim\operatorname{Unif}([0,1]^d),
\qquad
\mathcal L\bigl(A,Y(0),Y(1)\mid X=x\bigr)
=
\operatorname{Bernoulli}\bigl(e_\lambda(x)\bigr)
\otimes
\operatorname{Bernoulli}\bigl(\mu_{0,\theta,\eta}(x)\bigr)
\otimes
\operatorname{Bernoulli}\bigl(\mu_{1,\theta,\eta}(x)\bigr).
\]
Here \(X\) is the covariate vector, \(A\) is the binary treatment indicator, and \(Y(0),Y(1)\) are the binary potential responses under control and treatment. The prescribed functions \(e_\lambda,\mu_{0,\theta,\eta},\mu_{1,\theta,\eta}\) are the designated conditional-margin functions of this law.

\item Output the contrasts and finite pushforward measures
\[
\mathfrak C(h,\delta,a,b)
=
\bigl((\tau_\theta,\varpi_\theta)\bigr)_{\theta\in\{-1,+1\}},
\qquad
\varpi_\theta=
\sum_{\lambda,\eta\in\{-1,+1\}^{\mathcal I}}
w_\theta(\lambda,\eta)\,
\operatorname{Dirac}(P_{\theta,\lambda,\eta}),
\]
where \(\operatorname{Dirac}(P)\) is the unit point mass at the primitive law \(P\).
\end{enumerate}
\end{algorithmv}

The product inequality in \cref{obj:lem:hellinger-block-calculus} aggregates component distances. Its conditional identity integrates these distances over the common covariate distribution, subject to the stated measurable-density condition. The final identity accommodates the independent randomizer available to decisions in \cref{obj:def:risk}.

To translate mixture separation into expected absolute error, the analysis uses the fuzzy testing bound of \citet[Section 3, Lemma 1]{KennedyBalakrishnanRobinsWasserman2024}. \Cref{obj:lem:published-fuzzy-testing}, stated after the Hellinger calculus, fixes both its loss and its Hellinger normalization.

\begin{lemmav}{lem:hellinger-block-calculus}[Hellinger calculus for blocks]
For probability laws \(\mathbb B,\mathbb C\) on a common measurable space, use the normalization of \cref{obj:def:hellinger}:
\[
\mathsf h^2(\mathbb B,\mathbb C)
=\int\left(\sqrt{\frac{d\mathbb B}{d\xi}}
           -\sqrt{\frac{d\mathbb C}{d\xi}}\right)^2\,d\xi,
\qquad \xi=\mathbb B+\mathbb C.
\]
The following statements hold on arbitrary measurable spaces.

For any finite family of pairs of probability laws \(\mathbb B_{u_*},\mathbb C_{u_*}\), each pair on its respective common measurable space,
\[
\mathsf h^2\left(\bigotimes_{u_*}\mathbb B_{u_*},
                \bigotimes_{u_*}\mathbb C_{u_*}\right)
\le \sum_{u_*}\mathsf h^2(\mathbb B_{u_*},\mathbb C_{u_*}).
\]

Let \(\nu\) be a probability law on the conditioning space, and let \(\mathbb B_z,\mathbb C_z\) be probability kernels into a common measurable space. Suppose there exists a common s-finite reference kernel such that both kernels are represented, at every \(z\), by nonnegative extended-real densities with respect to that reference kernel, with both densities jointly measurable in the conditioning and outcome coordinates. Then
\[
\mathsf h^2\bigl(\nu(dz)\mathbb B_z(dw),
                \nu(dz)\mathbb C_z(dw)\bigr)
=\int \mathsf h^2(\mathbb B_z,\mathbb C_z)\,d\nu(z).
\]

For any probability laws \(\mathbb B,\mathbb C\) on a common measurable space and any probability law \(\nu\) on another measurable space,
\[
\mathsf h^2(\mathbb B\otimes\nu,\mathbb C\otimes\nu)
=\mathsf h^2(\mathbb B,\mathbb C).
\]
\end{lemmav}

\begin{proof}[Proof of \cref{obj:lem:hellinger-block-calculus}]
We first establish the finite-reference identity used below. Whenever a real Radon--Nikodym density is used, its value is taken to be zero where the extended-real derivative is infinite; for finite measures, this changes the derivative on a null set.

Let \(\rho_{\mathrm{ref}}\) be a finite measure on the common space \(\Omega\), and suppose measurable nonnegative real densities satisfy
\[
f_{\mathbb B,\mathrm{ref}},f_{\mathbb C,\mathrm{ref}}
   :\Omega\longrightarrow[0,\infty),
\qquad
\mathbb B(dw)=f_{\mathbb B,\mathrm{ref}}(w)\rho_{\mathrm{ref}}(dw),
\qquad
\mathbb C(dw)=f_{\mathbb C,\mathrm{ref}}(w)\rho_{\mathrm{ref}}(dw).
\]
Define
\[
\xi_{\Sigma}:=\mathbb B+\mathbb C,
\qquad
f_{\mathbb B,\Sigma}:=\frac{d\mathbb B}{d\xi_{\Sigma}},
\qquad
f_{\mathbb C,\Sigma}:=\frac{d\mathbb C}{d\xi_{\Sigma}},
\qquad
r_{\Sigma,\mathrm{ref}}:=\frac{d\xi_{\Sigma}}{d\rho_{\mathrm{ref}}},
\]
using the real-density convention above. Radon--Nikodym reconstruction and the chain rule give
\[
\mathbb B(dw)=f_{\mathbb B,\Sigma}(w)\xi_{\Sigma}(dw),
\qquad
\mathbb C(dw)=f_{\mathbb C,\Sigma}(w)\xi_{\Sigma}(dw),
\]
and, \(\rho_{\mathrm{ref}}\)-almost everywhere,
\[
f_{\mathbb B,\mathrm{ref}}
   =f_{\mathbb B,\Sigma}r_{\Sigma,\mathrm{ref}},
\qquad
f_{\mathbb C,\mathrm{ref}}
   =f_{\mathbb C,\Sigma}r_{\Sigma,\mathrm{ref}}.
\]
For these nonnegative densities,
\[
r_{\Sigma,\mathrm{ref}}
 \bigl(\sqrt{f_{\mathbb B,\Sigma}}
             -\sqrt{f_{\mathbb C,\Sigma}}\bigr)^2
=
\bigl(\sqrt{f_{\mathbb B,\Sigma}r_{\Sigma,\mathrm{ref}}}
             -\sqrt{f_{\mathbb C,\Sigma}r_{\Sigma,\mathrm{ref}}}\bigr)^2.
\]
Thus changing measure in the normalization of \cref{obj:def:hellinger} yields
\[
\begin{aligned}
\mathsf h^2(\mathbb B,\mathbb C)
&=\int
 r_{\Sigma,\mathrm{ref}}
 \bigl(\sqrt{f_{\mathbb B,\Sigma}}
             -\sqrt{f_{\mathbb C,\Sigma}}\bigr)^2
 \,d\rho_{\mathrm{ref}}\\
&=\int
 \bigl(\sqrt{f_{\mathbb B,\mathrm{ref}}}
             -\sqrt{f_{\mathbb C,\mathrm{ref}}}\bigr)^2
 \,d\rho_{\mathrm{ref}}.
\end{aligned}
\]
% lean: CausalSmith.Stat.TwosamplePointcateAnnotationFrontier.hellingerSq_real_density

\begin{enumerate}
\item \emph{Finite products.}
Write \(\mathcal U_*\) for the finite indexing set and \(\Omega_{u_*}\) for the measurable space of the pair indexed by \(u_*\). Define
\[
\xi_{u_*}:=\mathbb B_{u_*}+\mathbb C_{u_*},
\qquad
f_{\mathbb B,u_*}:=\frac{d\mathbb B_{u_*}}{d\xi_{u_*}},
\qquad
f_{\mathbb C,u_*}:=\frac{d\mathbb C_{u_*}}{d\xi_{u_*}},
\qquad u_*\in\mathcal U_*,
\]
again as real densities. They are measurable, nonnegative, and integrable, and reconstruction gives
\[
\mathbb B_{u_*}(dw)
   =f_{\mathbb B,u_*}(w)\xi_{u_*}(dw),
\qquad
\mathbb C_{u_*}(dw)
   =f_{\mathbb C,u_*}(w)\xi_{u_*}(dw),
\]
with
\[
\int f_{\mathbb B,u_*}\,d\xi_{u_*}
=
\int f_{\mathbb C,u_*}\,d\xi_{u_*}
=1.
\]
Define the coordinate affinity by
\[
\mathsf a_{\mathrm{aff},u_*}
:=
\int
\sqrt{f_{\mathbb B,u_*}(w)f_{\mathbb C,u_*}(w)}
\,d\xi_{u_*}(w).
\]
The square roots are square-integrable, so their product is integrable. Expanding the square in the finite-reference identity gives
\[
\begin{aligned}
\mathsf h^2(\mathbb B_{u_*},\mathbb C_{u_*})
&=\int
 \left(
 f_{\mathbb B,u_*}+f_{\mathbb C,u_*}
 -2\sqrt{f_{\mathbb B,u_*}f_{\mathbb C,u_*}}
 \right)\,d\xi_{u_*}\\
&=2\bigl(1-\mathsf a_{\mathrm{aff},u_*}\bigr).
\end{aligned}
\]

On the product space, define
\[
\xi_{\otimes}:=\bigotimes_{u_*\in\mathcal U_*}\xi_{u_*},
\qquad
f_{\mathbb B,\otimes}(w)
   :=\prod_{u_*\in\mathcal U_*}f_{\mathbb B,u_*}(w_{u_*}),
\qquad
f_{\mathbb C,\otimes}(w)
   :=\prod_{u_*\in\mathcal U_*}f_{\mathbb C,u_*}(w_{u_*}),
\quad
w\in\prod_{u_*\in\mathcal U_*}\Omega_{u_*}.
\]
To verify the product-density representation, take measurable sets
\[
W_{u_*}\subseteq\Omega_{u_*},
\qquad u_*\in\mathcal U_*.
\]
Product integration gives
\[
\int_{\prod_{u_*}W_{u_*}}
 f_{\mathbb B,\otimes}\,d\xi_{\otimes}
=
\prod_{u_*}
 \int_{W_{u_*}}f_{\mathbb B,u_*}\,d\xi_{u_*}
=
\prod_{u_*}\mathbb B_{u_*}(W_{u_*}).
\]
The same calculation applies to \(\mathbb C\). Uniqueness of finite product measures therefore yields
\[
\left(\bigotimes_{u_*}\mathbb B_{u_*}\right)(dw)
   =f_{\mathbb B,\otimes}(w)\xi_{\otimes}(dw),
\qquad
\left(\bigotimes_{u_*}\mathbb C_{u_*}\right)(dw)
   =f_{\mathbb C,\otimes}(w)\xi_{\otimes}(dw).
\]
In particular,
\[
\int f_{\mathbb B,\otimes}\,d\xi_{\otimes}
=
\int f_{\mathbb C,\otimes}\,d\xi_{\otimes}
=1.
\]
Applying the finite-reference identity and expanding the square as above gives
\[
\mathsf h^2\left(\bigotimes_{u_*}\mathbb B_{u_*},
                \bigotimes_{u_*}\mathbb C_{u_*}\right)
=
2\left(
1-\int\sqrt{f_{\mathbb B,\otimes}f_{\mathbb C,\otimes}}
\,d\xi_{\otimes}
\right).
\]
Nonnegativity of the coordinate densities implies the pointwise identity
\[
\sqrt{f_{\mathbb B,\otimes}(w)f_{\mathbb C,\otimes}(w)}
=
\prod_{u_*}
\sqrt{f_{\mathbb B,u_*}(w_{u_*})
      f_{\mathbb C,u_*}(w_{u_*})}.
\]
Product integration consequently gives
\[
\int\sqrt{f_{\mathbb B,\otimes}f_{\mathbb C,\otimes}}
\,d\xi_{\otimes}
=
\prod_{u_*}\mathsf a_{\mathrm{aff},u_*}.
\]

The affinities satisfy
\[
\mathsf a_{\mathrm{aff},u_*}\ge0,
\qquad
0\le\mathsf h^2(\mathbb B_{u_*},\mathbb C_{u_*})
   =2\bigl(1-\mathsf a_{\mathrm{aff},u_*}\bigr),
\]
and hence
\[
0\le\mathsf a_{\mathrm{aff},u_*}\le1.
\]
For every subset \(\mathcal V_*\subseteq\mathcal U_*\), induction on its cardinality gives
\[
0\le\prod_{u_*\in\mathcal V_*}\mathsf a_{\mathrm{aff},u_*}\le1.
\]
A second induction gives the required affinity-defect bound. The empty product is one and the empty sum is zero. For \(u_*\in\mathcal U_*\setminus\mathcal V_*\), the induction step is
\[
\begin{aligned}
1-\prod_{u'_*\in\mathcal V_*\cup\{u_*\}}
       \mathsf a_{\mathrm{aff},u'_*}
&=
\left(1-\prod_{u'_*\in\mathcal V_*}
       \mathsf a_{\mathrm{aff},u'_*}\right)
+
\bigl(1-\mathsf a_{\mathrm{aff},u_*}\bigr)
 \prod_{u'_*\in\mathcal V_*}\mathsf a_{\mathrm{aff},u'_*}\\
&\le
\sum_{u'_*\in\mathcal V_*}
 \bigl(1-\mathsf a_{\mathrm{aff},u'_*}\bigr)
+
\bigl(1-\mathsf a_{\mathrm{aff},u_*}\bigr).
\end{aligned}
\]
Taking \(\mathcal V_*=\mathcal U_*\) and multiplying by two proves
\[
\begin{aligned}
\mathsf h^2\left(\bigotimes_{u_*}\mathbb B_{u_*},
                \bigotimes_{u_*}\mathbb C_{u_*}\right)
&=2\left(1-\prod_{u_*}\mathsf a_{\mathrm{aff},u_*}\right)\\
&\le2\sum_{u_*}\bigl(1-\mathsf a_{\mathrm{aff},u_*}\bigr)
=\sum_{u_*}\mathsf h^2(\mathbb B_{u_*},\mathbb C_{u_*}).
\end{aligned}
\]
The empty family is included by the stated empty-product and empty-sum conventions.
% lean: CausalSmith.Stat.TwosamplePointcateAnnotationFrontier.hellingerSq_pi_le_sum

\item \emph{A common conditioning marginal.}
Let \(Z\) and \(\Omega\) be the conditioning and outcome spaces. Choose the s-finite reference kernel and jointly measurable densities supplied by the hypothesis, writing them as
\[
\Xi_{\mathrm{ref}}:Z\longrightarrow\{\text{s-finite measures on }\Omega\},
\qquad
f_{\mathbb B,\mathrm{ref}},f_{\mathbb C,\mathrm{ref}}
   :Z\times\Omega\longrightarrow[0,\infty],
\]
so that, for every \(z\in Z\),
\[
\mathbb B_z(dw)
   =f_{\mathbb B,\mathrm{ref}}(z,w)\Xi_{\mathrm{ref},z}(dw),
\qquad
\mathbb C_z(dw)
   =f_{\mathbb C,\mathrm{ref}}(z,w)\Xi_{\mathrm{ref},z}(dw).
\]
Define the sum kernel and normalized densities by
\[
\mathbb K_{\Sigma,z}:=\mathbb B_z+\mathbb C_z
\]
and
\[
\bigl(r_{\mathbb B}(z,w),r_{\mathbb C}(z,w)\bigr)
:=
\begin{cases}
\displaystyle
\left(
\frac{f_{\mathbb B,\mathrm{ref}}(z,w)}
     {f_{\mathbb B,\mathrm{ref}}(z,w)+f_{\mathbb C,\mathrm{ref}}(z,w)},
\frac{f_{\mathbb C,\mathrm{ref}}(z,w)}
     {f_{\mathbb B,\mathrm{ref}}(z,w)+f_{\mathbb C,\mathrm{ref}}(z,w)}
\right),
&
\displaystyle
0<f_{\mathbb B,\mathrm{ref}}(z,w)+f_{\mathbb C,\mathrm{ref}}(z,w)<\infty,
\\[6pt]
(0,0),&\text{otherwise}.
\end{cases}
\]
These are jointly measurable real functions on all of \(Z\times\Omega\), and
\[
0\le r_{\mathbb B}(z,w)\le1,
\qquad
0\le r_{\mathbb C}(z,w)\le1,
\]
because each numerator lies between zero and its positive finite denominator; the other cases have value zero.

For every \(z\),
\[
\mathbb K_{\Sigma,z}(dw)
=
\bigl(f_{\mathbb B,\mathrm{ref}}(z,w)
     +f_{\mathbb C,\mathrm{ref}}(z,w)\bigr)\Xi_{\mathrm{ref},z}(dw),
\]
and the nonnegative extended integral satisfies
\[
\int
\bigl(f_{\mathbb B,\mathrm{ref}}(z,w)
     +f_{\mathbb C,\mathrm{ref}}(z,w)\bigr)
\,d\Xi_{\mathrm{ref},z}(w)
=2.
\]
Thus the sum of the densities is finite \(\Xi_{\mathrm{ref},z}\)-almost everywhere. Where that sum is zero, both densities are zero. Cancellation therefore gives, \(\Xi_{\mathrm{ref},z}\)-almost everywhere,
\[
\begin{aligned}
\bigl(f_{\mathbb B,\mathrm{ref}}+f_{\mathbb C,\mathrm{ref}}\bigr)(z,w)
\,r_{\mathbb B}(z,w)
&=f_{\mathbb B,\mathrm{ref}}(z,w),\\
\bigl(f_{\mathbb B,\mathrm{ref}}+f_{\mathbb C,\mathrm{ref}}\bigr)(z,w)
\,r_{\mathbb C}(z,w)
&=f_{\mathbb C,\mathrm{ref}}(z,w).
\end{aligned}
\]
Consequently, for every \(z\),
\[
\mathbb B_z(dw)=r_{\mathbb B}(z,w)\mathbb K_{\Sigma,z}(dw),
\qquad
\mathbb C_z(dw)=r_{\mathbb C}(z,w)\mathbb K_{\Sigma,z}(dw).
\]

Define the finite joint reference measure by
\[
\Xi_{\Sigma}(dz,dw):=\nu(dz)\mathbb K_{\Sigma,z}(dw).
\]
The kernel density representations and joint measurability imply
\[
\nu(dz)\mathbb B_z(dw)
   =r_{\mathbb B}(z,w)\Xi_{\Sigma}(dz,dw),
\qquad
\nu(dz)\mathbb C_z(dw)
   =r_{\mathbb C}(z,w)\Xi_{\Sigma}(dz,dw).
\]
Moreover,
\[
\Xi_{\Sigma}(Z\times\Omega)
=\int\mathbb K_{\Sigma,z}(\Omega)\,d\nu(z)=2.
\]
The normalized-density bounds give
\[
0\le\sqrt{r_{\mathbb B}(z,w)}\le1,
\qquad
0\le\sqrt{r_{\mathbb C}(z,w)}\le1,
\]
and hence
\[
0\le
\bigl(\sqrt{r_{\mathbb B}(z,w)}
      -\sqrt{r_{\mathbb C}(z,w)}\bigr)^2
\le4.
\]
The measurable square-root difference is therefore integrable against \(\Xi_{\Sigma}\). Applying the finite-reference identity and then integration over the kernel gives
\[
\begin{aligned}
&\mathsf h^2\bigl(\nu(dz)\mathbb B_z(dw),
                 \nu(dz)\mathbb C_z(dw)\bigr)\\
&\quad=
\int
\bigl(\sqrt{r_{\mathbb B}(z,w)}
      -\sqrt{r_{\mathbb C}(z,w)}\bigr)^2
\,d\Xi_{\Sigma}(z,w)\\
&\quad=
\int_Z
\left[
\int_\Omega
\bigl(\sqrt{r_{\mathbb B}(z,w)}
      -\sqrt{r_{\mathbb C}(z,w)}\bigr)^2
\,d\mathbb K_{\Sigma,z}(w)
\right]\,d\nu(z).
\end{aligned}
\]
For each \(z\), the same finite-reference identity and the conditional density representations yield
\[
\int_\Omega
\bigl(\sqrt{r_{\mathbb B}(z,w)}
      -\sqrt{r_{\mathbb C}(z,w)}\bigr)^2
\,d\mathbb K_{\Sigma,z}(w)
=
\mathsf h^2(\mathbb B_z,\mathbb C_z).
\]
Substitution proves
\[
\mathsf h^2\bigl(\nu(dz)\mathbb B_z(dw),
                \nu(dz)\mathbb C_z(dw)\bigr)
=
\int\mathsf h^2(\mathbb B_z,\mathbb C_z)\,d\nu(z).
\]
% lean: CausalSmith.Stat.TwosamplePointcateAnnotationFrontier.hellingerSq_compProd_common

\item \emph{A common independent probability factor.}
For the given \(\mathbb B,\mathbb C\), use
\[
\xi_{\Sigma}:=\mathbb B+\mathbb C,
\qquad
f_{\mathbb B,\Sigma}:=\frac{d\mathbb B}{d\xi_{\Sigma}},
\qquad
f_{\mathbb C,\Sigma}:=\frac{d\mathbb C}{d\xi_{\Sigma}},
\]
with the real-density convention above. The sum of the product laws is
\[
(\mathbb B\otimes\nu)+(\mathbb C\otimes\nu)
=\xi_{\Sigma}\otimes\nu.
\]
The product Radon--Nikodym identities give, \((\xi_{\Sigma}\otimes\nu)\)-almost everywhere,
\[
\frac{d(\mathbb B\otimes\nu)}{d(\xi_{\Sigma}\otimes\nu)}(w,z)
=f_{\mathbb B,\Sigma}(w),
\qquad
\frac{d(\mathbb C\otimes\nu)}{d(\xi_{\Sigma}\otimes\nu)}(w,z)
=f_{\mathbb C,\Sigma}(w).
\]
Substituting into the defining square-root formula and integrating a function of the first coordinate gives
\[
\begin{aligned}
\mathsf h^2(\mathbb B\otimes\nu,\mathbb C\otimes\nu)
&=
\int_{\Omega\times Z}
\bigl(\sqrt{f_{\mathbb B,\Sigma}(w)}
      -\sqrt{f_{\mathbb C,\Sigma}(w)}\bigr)^2
\,d(\xi_{\Sigma}\otimes\nu)(w,z)\\
&=
\nu(Z)
\int_\Omega
\bigl(\sqrt{f_{\mathbb B,\Sigma}(w)}
      -\sqrt{f_{\mathbb C,\Sigma}(w)}\bigr)^2
\,d\xi_{\Sigma}(w)\\
&=\mathsf h^2(\mathbb B,\mathbb C),
\end{aligned}
\]
where the last equality uses \(\nu(Z)=1\).
% lean: CausalSmith.Stat.TwosamplePointcateAnnotationFrontier.hellingerSq_prod_probability
\end{enumerate}
\end{proof}

For the original-record lower bound, the testing comparison concerns the law of the entire fixed-size dataset. At this level, the observation space contains both record blocks, and the testing statement can be used with product count one. This formulation preserves the unequal record counts in the experiment. The certificate also records an exact equality of the auxiliary mixture laws; the relevant marginal is the treatment-record marginal of \cref{obj:synth_4}.

\begin{lemmav}{lem:published-fuzzy-testing}[Absolute-loss fuzzy testing bound]
Let \((\mathbb B_z)\) and \((\mathbb C_z)\) be probability kernels from a common measurable index space to a common measurable observation space, and let \(\omega\) be a probability prior on the index space. Let \(v_0\ge1\) be an integer, let \(\mathcal M\) be a model of measures containing \(\mathbb B_z\) and \(\mathbb C_z\) for every \(z\), and let \(\Psi\) be a real-valued functional on measures. Use the squared Hellinger distance of \cref{obj:def:hellinger}, with normalization
\[
\mathsf h^2(\mathbb B,\mathbb C)
=
\int
\left(
\sqrt{\frac{d\mathbb B}{d\xi}}
-
\sqrt{\frac{d\mathbb C}{d\xi}}
\right)^2\,d\xi,
\qquad
\xi=\mathbb B+\mathbb C.
\]
Suppose the following conditions hold:
\begin{itemize}
\item \textbf{(Constants.)} The separation bound \(s_0\) and distance bound \(u_0\) satisfy \(s_0>0\) and \(0\le u_0<2\).
\item \textbf{(Separation.)} For every pair of indices \(z,z'\),
\[
\Psi(\mathbb B_z)-\Psi(\mathbb C_{z'})\ge s_0.
\]
\item \textbf{(Mixture distance.)} The product-law mixtures satisfy
\[
\mathsf h^2\left(
\int \mathbb B_z^{\otimes v_0}\,d\omega(z),
\int \mathbb C_z^{\otimes v_0}\,d\omega(z)
\right)\le u_0.
\]
\end{itemize}
Then every measurable real-valued decision \(\widehat\Psi\) on the \(v_0\)-observation sample space satisfies
\[
\sup_{P\in\mathcal M}
\int
\left|\widehat\Psi(w)-\Psi(P)\right|
\,dP^{\otimes v_0}(w)
\ge
\frac{s_0}{4}
\left\{1-\sqrt{u_0(1-u_0/4)}\right\},
\]
where the nonnegative loss integral is interpreted in the extended sense. This is the absolute-loss specialization of \citep[Section 3, Lemma 1]{KennedyBalakrishnanRobinsWasserman2024}.
\end{lemmav}

\begin{proof}[Proof of \cref{obj:lem:published-fuzzy-testing}]
Fix a measurable decision \(\widehat\Psi\). All loss integrals and their suprema below take values in \([0,\infty]\).

\begin{enumerate}
\item Define the testing submodel, its component risks, and its worst-case risk by
\[
\begin{aligned}
\mathcal M_{\mathrm{test}}
&:=\{\mathbb B_z:z\text{ is an index}\}
  \cup\{\mathbb C_z:z\text{ is an index}\},\\
\mathcal R_{\mathrm{test}}(P)
&:=\int
  |\widehat\Psi(w)-\Psi(P)|\,dP^{\otimes v_0}(w),
  \qquad P\text{ a measure on the observation space},\\
\mathcal S_{\mathrm{test}}
&:=\sup_{P\in\mathcal M_{\mathrm{test}}}
  \mathcal R_{\mathrm{test}}(P).
\end{aligned}
\]
Every member of \(\mathcal M_{\mathrm{test}}\) is a probability measure. The finite-product maps
\(z\mapsto\mathbb B_z^{\otimes v_0}\) and
\(z\mapsto\mathbb C_z^{\otimes v_0}\) are probability kernels. Define their mixture laws by
\[
\mathbb B_{\mathrm{mix}}
:=\int\mathbb B_z^{\otimes v_0}\,d\omega(z),
\qquad
\mathbb C_{\mathrm{mix}}
:=\int\mathbb C_z^{\otimes v_0}\,d\omega(z).
\]
Both mixtures are probability measures, since their component laws and the prior are probability measures. The distance hypothesis therefore gives
\[
\mathsf h^2(\mathbb B_{\mathrm{mix}},\mathbb C_{\mathrm{mix}})
\le u_0.
\]
% lean: Causalean.Stat.Minimax.absolute_fuzzy_testing_lower_bound_iid

\item The probability prior ensures that the index space is nonempty. Choose an index \(z_0\). Separation gives
\[
\Psi(\mathbb C_z)\le \Psi(\mathbb B_{z_0})-s_0
\qquad\text{for every }z.
\]
Thus the nonempty set of values \(\Psi(\mathbb C_z)\) is bounded above. Define the finite real numbers
\[
\psi_{\mathbb C}^{\mathrm{sup}}
:=\sup_z\Psi(\mathbb C_z),
\qquad
a_{\mathrm{test}}
:=\psi_{\mathbb C}^{\mathrm{sup}}+\frac{s_0}{2}.
\]
For each index \(z\), separation with every \(z'\) gives
\[
\Psi(\mathbb C_{z'})
\le\Psi(\mathbb B_z)-s_0,
\qquad
\psi_{\mathbb C}^{\mathrm{sup}}
\le\Psi(\mathbb B_z)-s_0.
\]
Consequently,
\[
a_{\mathrm{test}}+\frac{s_0}{2}
\le\Psi(\mathbb B_z),
\qquad
\Psi(\mathbb C_z)
\le a_{\mathrm{test}}-\frac{s_0}{2}.
\]
Define the measurable threshold event on the sample space by
\[
\mathcal A_{\mathrm{test}}
:=\{w:\widehat\Psi(w)<a_{\mathrm{test}}\}.
\]
% lean: Causalean.Stat.Minimax.exists_separating_threshold

\item We derive the Hellinger testing comparison for these mixtures. Define
\[
\begin{aligned}
\xi_{\mathrm{test}}&:=\mathbb B_{\mathrm{mix}}+\mathbb C_{\mathrm{mix}},\\
f_{\mathbb B}(w)&:=\frac{d\mathbb B_{\mathrm{mix}}}{d\xi_{\mathrm{test}}}(w),
&
f_{\mathbb C}(w)&:=\frac{d\mathbb C_{\mathrm{mix}}}{d\xi_{\mathrm{test}}}(w),\\
\mathsf h_{\mathrm{test}}^2
&:=\mathsf h^2(\mathbb B_{\mathrm{mix}},\mathbb C_{\mathrm{mix}}).
\end{aligned}
\]
Choose nonnegative versions of the densities. They are integrable, and
\[
\int f_{\mathbb B}\,d\xi_{\mathrm{test}}
=\int f_{\mathbb C}\,d\xi_{\mathrm{test}}=1,
\qquad
\int(f_{\mathbb B}-f_{\mathbb C})\,d\xi_{\mathrm{test}}=0.
\]
For every measurable sample event \(\mathcal A\),
\[
\begin{aligned}
\mathbb B_{\mathrm{mix}}(\mathcal A)
-\mathbb C_{\mathrm{mix}}(\mathcal A)
&=\int_{\mathcal A}(f_{\mathbb B}-f_{\mathbb C})\,d\xi_{\mathrm{test}},\\
\int_{\mathcal A^c}(f_{\mathbb B}-f_{\mathbb C})\,d\xi_{\mathrm{test}}
&=-\int_{\mathcal A}(f_{\mathbb B}-f_{\mathbb C})\,d\xi_{\mathrm{test}}.
\end{aligned}
\]
Applying the integral triangle inequality on each of \(\mathcal A\) and
\(\mathcal A^c\), then adding, yields
\[
2\left|
\int_{\mathcal A}(f_{\mathbb B}-f_{\mathbb C})\,d\xi_{\mathrm{test}}
\right|
\le
\int|f_{\mathbb B}-f_{\mathbb C}|\,d\xi_{\mathrm{test}}.
\]

The normalization in \cref{obj:def:hellinger} gives
\[
\mathsf h_{\mathrm{test}}^2
=\int(\sqrt{f_{\mathbb B}}-\sqrt{f_{\mathbb C}})^2\,d\xi_{\mathrm{test}}.
\]
The square roots of the densities are square integrable. The pointwise identities
\[
\begin{aligned}
(\sqrt{f_{\mathbb B}}+\sqrt{f_{\mathbb C}})^2
+(\sqrt{f_{\mathbb B}}-\sqrt{f_{\mathbb C}})^2
&=2f_{\mathbb B}+2f_{\mathbb C},\\
|f_{\mathbb B}-f_{\mathbb C}|
&=|\sqrt{f_{\mathbb B}}-\sqrt{f_{\mathbb C}}|\,
  |\sqrt{f_{\mathbb B}}+\sqrt{f_{\mathbb C}}|
\end{aligned}
\]
therefore imply
\[
\int(\sqrt{f_{\mathbb B}}+\sqrt{f_{\mathbb C}})^2\,d\xi_{\mathrm{test}}
=4-\mathsf h_{\mathrm{test}}^2.
\]
Also, by nonnegativity and the pointwise bound
\((\sqrt{f_{\mathbb B}}-\sqrt{f_{\mathbb C}})^2
\le f_{\mathbb B}+f_{\mathbb C}\),
\[
0\le\mathsf h_{\mathrm{test}}^2
\le\int(f_{\mathbb B}+f_{\mathbb C})\,d\xi_{\mathrm{test}}
=2.
\]
Cauchy--Schwarz now gives
\[
\begin{aligned}
\int|f_{\mathbb B}-f_{\mathbb C}|\,d\xi_{\mathrm{test}}
&\le
\sqrt{\mathsf h_{\mathrm{test}}^2}\,
\sqrt{4-\mathsf h_{\mathrm{test}}^2}\\
&=\sqrt{\mathsf h_{\mathrm{test}}^2
                 (4-\mathsf h_{\mathrm{test}}^2)}.
\end{aligned}
\]
Combining the preceding bounds proves, for every measurable \(\mathcal A\),
\[
\begin{aligned}
|\mathbb B_{\mathrm{mix}}(\mathcal A)
-\mathbb C_{\mathrm{mix}}(\mathcal A)|
&\le\frac12
 \sqrt{\mathsf h_{\mathrm{test}}^2(4-\mathsf h_{\mathrm{test}}^2)}\\
&=\sqrt{\mathsf h_{\mathrm{test}}^2
                   (1-\mathsf h_{\mathrm{test}}^2/4)}.
\end{aligned}
\]
% lean: Causalean.Stat.Minimax.measure_event_gap_le_sharp_hellinger

\item Define the total variation distance of the mixtures by
\[
\mathsf V_{\mathrm{test}}
:=\sup_{\mathcal A\text{ measurable}}
|\mathbb B_{\mathrm{mix}}(\mathcal A)
-\mathbb C_{\mathrm{mix}}(\mathcal A)|.
\]
Taking the supremum in the preceding event bound gives
\[
\mathsf V_{\mathrm{test}}
\le
\sqrt{\mathsf h_{\mathrm{test}}^2
                   (1-\mathsf h_{\mathrm{test}}^2/4)}.
\]
The complementary-event identity implies
\[
\begin{aligned}
\mathbb B_{\mathrm{mix}}(\mathcal A_{\mathrm{test}})
+\mathbb C_{\mathrm{mix}}(\mathcal A_{\mathrm{test}}^c)
&=1+\mathbb B_{\mathrm{mix}}(\mathcal A_{\mathrm{test}})
       -\mathbb C_{\mathrm{mix}}(\mathcal A_{\mathrm{test}})\\
&\ge 1-\mathsf V_{\mathrm{test}}.
\end{aligned}
\]
Since \(0\le\mathsf h_{\mathrm{test}}^2\le u_0<2\),
\[
u_0(1-u_0/4)
-\mathsf h_{\mathrm{test}}^2(1-\mathsf h_{\mathrm{test}}^2/4)
=
\frac{(u_0-\mathsf h_{\mathrm{test}}^2)
      (4-u_0-\mathsf h_{\mathrm{test}}^2)}{4}
\ge0.
\]
Monotonicity of the square root consequently yields
\[
1-\sqrt{u_0(1-u_0/4)}
\le
\mathbb B_{\mathrm{mix}}(\mathcal A_{\mathrm{test}})
+\mathbb C_{\mathrm{mix}}(\mathcal A_{\mathrm{test}}^c).
\]
% lean: Causalean.Stat.Minimax.binary_test_error_ge_of_hellinger_budget

\item The threshold separation gives the pointwise loss bounds, for every index \(z\),
\[
\begin{aligned}
\frac{s_0}{2}\mathbf 1_{\mathcal A_{\mathrm{test}}}(w)
&\le|\widehat\Psi(w)-\Psi(\mathbb B_z)|,\\
\frac{s_0}{2}\mathbf 1_{\mathcal A_{\mathrm{test}}^c}(w)
&\le|\widehat\Psi(w)-\Psi(\mathbb C_z)|.
\end{aligned}
\]
Indeed, on \(\mathcal A_{\mathrm{test}}\) the decision lies below
\(a_{\mathrm{test}}\), while \(\Psi(\mathbb B_z)\) lies at least \(s_0/2\)
above it. On its complement the decision lies at or above the threshold,
while \(\Psi(\mathbb C_z)\) lies at least \(s_0/2\) below it. Outside the
respective events, the displayed left-hand sides vanish.

Integrating these inequalities under the component product laws gives
\[
\begin{aligned}
\frac{s_0}{2}\mathbb B_z^{\otimes v_0}(\mathcal A_{\mathrm{test}})
&\le\mathcal R_{\mathrm{test}}(\mathbb B_z)
\le\mathcal S_{\mathrm{test}},\\
\frac{s_0}{2}\mathbb C_z^{\otimes v_0}(\mathcal A_{\mathrm{test}}^c)
&\le\mathcal R_{\mathrm{test}}(\mathbb C_z)
\le\mathcal S_{\mathrm{test}}.
\end{aligned}
\]
Kernel event probabilities are measurable functions of the index. Integrating
their bounds against the probability prior therefore gives
\[
\begin{aligned}
\frac{s_0}{2}\mathbb B_{\mathrm{mix}}(\mathcal A_{\mathrm{test}})
&=\int\frac{s_0}{2}
  \mathbb B_z^{\otimes v_0}(\mathcal A_{\mathrm{test}})\,d\omega(z)
\le\int\mathcal S_{\mathrm{test}}\,d\omega
=\mathcal S_{\mathrm{test}},\\
\frac{s_0}{2}\mathbb C_{\mathrm{mix}}(\mathcal A_{\mathrm{test}}^c)
&=\int\frac{s_0}{2}
  \mathbb C_z^{\otimes v_0}(\mathcal A_{\mathrm{test}}^c)\,d\omega(z)
\le\int\mathcal S_{\mathrm{test}}\,d\omega
=\mathcal S_{\mathrm{test}}.
\end{aligned}
\]
Adding yields the extended nonnegative inequality
\[
\frac{s_0}{2}
\left\{
\mathbb B_{\mathrm{mix}}(\mathcal A_{\mathrm{test}})
+\mathbb C_{\mathrm{mix}}(\mathcal A_{\mathrm{test}}^c)
\right\}
\le2\mathcal S_{\mathrm{test}}.
\]
% lean: Causalean.Stat.Minimax.separated_mixture_errors_le_two_mul_sup

\item Combining the testing and risk inequalities gives
\[
\frac{s_0}{2}\left\{1-\sqrt{u_0(1-u_0/4)}\right\}
\le2\mathcal S_{\mathrm{test}}.
\]
The bracket is nonnegative under \(0\le u_0<2\). Using
\[
\frac{s_0}{2}\left\{1-\sqrt{u_0(1-u_0/4)}\right\}
=
2\left[
\frac{s_0}{4}\left\{1-\sqrt{u_0(1-u_0/4)}\right\}
\right]
\]
and cancelling the finite positive factor \(2\) in \([0,\infty]\), we obtain
\[
\frac{s_0}{4}\left\{1-\sqrt{u_0(1-u_0/4)}\right\}
\le\mathcal S_{\mathrm{test}}.
\]
Finally, every member of \(\mathcal M_{\mathrm{test}}\) equals
\(\mathbb B_z\) or \(\mathbb C_z\) for some index \(z\); either membership
hypothesis places that measure in \(\mathcal M\). Thus
\[
\mathcal M_{\mathrm{test}}\subseteq\mathcal M,
\qquad
\mathcal S_{\mathrm{test}}
\le
\sup_{P\in\mathcal M}
\int|\widehat\Psi(w)-\Psi(P)|\,dP^{\otimes v_0}(w).
\]
Transitivity proves the asserted bound, including when the supremum is infinite.
% lean: CausalSmith.Stat.TwosamplePointcateAnnotationFrontier.PublishedFuzzyTesting
\end{enumerate}
\end{proof}

The marginal in \cref{obj:synth_4} is the sampling law of each outcome-free record. The next result establishes model membership, target separation, and the information bound for the prescribed smooth-frame construction. Its occupancy condition restricts the expected number of treatment records at the fine spatial scale.

\begin{theoremv}{thm:marked-component-certificate}[Marked component mixture certificate]
Fix the primitive class \(\mathcal M\) of \cref{obj:def:primitive-class}, with public parameters satisfying
\[
d\in\mathbb N,\qquad d\ge1,\qquad
\alpha,\beta\in(0,1],\qquad
\gamma\in[1,3],\qquad
L>\tfrac12,\qquad
\varepsilon\in(0,\tfrac12).
\]
There exist constants \(c\in(0,1]\), \(c_0>0\), and \(C>0\), depending only on these public parameters, such that the following holds for every choice satisfying:
\begin{itemize}
\item \textbf{(Record counts.)} \(n,m\in\mathbb N\), \(n\ge2\), with total record count \(N=n+m\).
\item \textbf{(Scales.)} \(0<\delta\le h\le\tfrac12\).
\item \textbf{(Propensity amplitude.)} \(0<a\le c\delta^\alpha\).
\item \textbf{(Response amplitude.)} \(0<b\le c\delta^\beta\).
\item \textbf{(Effect calibration.)} \(ab\le ch^\gamma\).
\item \textbf{(Occupancy.)} \(N\delta^d\le c_0\).
\end{itemize}
The prescribed construction \(\mathfrak C\) of \cref{obj:def:marked-handle}, evaluated at \(d,h,\delta,a,b\), has the following properties for both hypothesis indices \(\theta\in\{-1,+1\}\).

Its weights on the finite sign arrays are nonnegative and normalized:
\[
w_\theta(\lambda,\eta)\ge0
\quad\text{for every }\lambda,\eta\in\{-1,+1\}^{\mathcal I},
\qquad
\sum_{\lambda,\eta\in\{-1,+1\}^{\mathcal I}}
w_\theta(\lambda,\eta)=1.
\]
Every constructed law \(P_{\theta,\lambda,\eta}\) belongs to \(\mathcal M\), and its canonical contrast agrees with the deterministic constructed contrast everywhere on the unit cube:
\[
\tau_{P_{\theta,\lambda,\eta}}(x)=\tau_\theta(x)
\qquad
\bigl(x\in[0,1]^d\bigr).
\]
Thus the finite prior \(\varpi_\theta\) is supported on \(\mathcal M\). The designated propensity functions, restricted to the unit cube, have the same distribution under \(\varpi_{+1}\) and \(\varpi_{-1}\).

At the interior center \(x_0=(1/2,\ldots,1/2)\), the contrasts satisfy
\[
\tau_\theta(x_0)=-2\theta ab,
\qquad
\tau_{-1}(x_0)-\tau_{+1}(x_0)=4ab,
\qquad
cab\le4ab.
\]
For the fixed-size original-record mixture laws \(M_{+1}\) and \(M_{-1}\), the unhalved squared Hellinger distance of \cref{obj:def:hellinger} satisfies
\[
H^2(M_{+1},M_{-1})
\le CnNh^d\delta^da^2b^2.
\]

For every integer \(k_{\mathrm{aux}}\ge0\), the auxiliary-record mixtures agree:
\[
\int P_{XA}^{\otimes k_{\mathrm{aux}}}\,d\varpi_{+1}(P)
=
\int P_{XA}^{\otimes k_{\mathrm{aux}}}\,d\varpi_{-1}(P).
\]
Their explicit conditional auxiliary kernels also agree at every collection of \(k_{\mathrm{aux}}\) covariates in \([0,1]^d\). For either hypothesis and every \(x\in[0,1]^d\), the singleton labeled mixture conditional on \(X=x\) is uniform on the four values of \((A,Y)\in\{0,1\}^2\).

Conditional on every collection of all \(N\) covariates in the unit cube, form the graph whose vertices are all \(N\) records. Join two distinct vertices when both covariates lie in the localization cube \(C_h\) and their maximum-norm distance is at most \(2\delta\); records exterior to \(C_h\) are isolated vertices. The conditional mixed density equals the product of the component densities over all connected components of this graph, including the exterior singleton components. Every auxiliary-only component and every component consisting of a single labeled record has the same conditional density, and hence the same conditional law, under the two hypotheses.
\end{theoremv}

The information bound in \cref{obj:thm:marked-component-certificate} reflects where the two hypotheses can be distinguished. Their common propensity distribution gives identical auxiliary mixtures for every auxiliary sample size, including equality conditional on the auxiliary covariates. Compensation in the response means also makes a singleton labeled mixture uniform under either hypothesis. Conditional component factorization therefore locates distinguishing information in components containing a labeled record and at least one additional nearby record.

The factor \(nN\) in the distance bound corresponds to record pairs involving a labeled member, while \(h^d\delta^d\) records the localization and proximity scales. The squared perturbation factor is \(a^2b^2\), and the target separation is \(4ab\). Under the stated sparse-occupancy gate \(N\delta^d\le c_0\), with occupancy constant \leanref{sym:c0}{\ensuremath{c_0}}, the certificate bounds the aggregate information by this pair-scale expression. These clauses give the component structure needed by \cref{obj:lem:hellinger-block-calculus} and the separation needed by the testing bound of \cref{obj:lem:published-fuzzy-testing}.

The next result packages the smoothness, contrast, and occupancy calibrations into a minimax lower bound. Its sample-size interval is the range in which the unequal-count rate is at least the supplied-propensity order.

\begin{lemmav}{lem:nuisance-frontier-converse}[Nuisance frontier lower bound]
Fix \(d\in\mathbb N\) and \(\alpha,\beta,\gamma,L,\varepsilon\in\mathbb R\), where \(\varepsilon\) is the overlap level. Define the nuisance smoothness sum \(S\), the smoothness threshold \(S_{\mathrm{crit}}\), the record-growth exponent \(q_*\), and the rate denominator \(\Delta\) by
\[
S=\alpha+\beta,\qquad
S_{\mathrm{crit}}=\frac{\gamma d}{2\gamma+d},\qquad
q_*=\frac{\gamma d}{(\alpha+\beta)(2\gamma+d)},\qquad
\Delta=2\gamma+d+\frac{\gamma d}{\alpha+\beta}.
\]
Suppose that:
\begin{itemize}
\item \textbf{(Dimension.)} \(d\ge1\).
\item \textbf{(Nuisance smoothness.)} \(0<\alpha\le1\) and \(0<\beta\le1\).
\item \textbf{(Effect smoothness.)} \(1\le\gamma\le3\).
\item \textbf{(Radius.)} \(L>1/2\).
\item \textbf{(Overlap level.)} \(0<\varepsilon<1/2\).
\item \textbf{(Frontier regime.)} \(S<S_{\mathrm{crit}}\).
\end{itemize}
There exists a constant \(c>0\), depending on the fixed public parameters \(d,\alpha,\beta,\gamma,L,\varepsilon\), such that for every \(n,m\in\mathbb N\) with \(n\ge2\), the total treatment-record count
\[
N=n+m
\]
satisfying \(N\le n^{q_*}\) obeys the following bound for the minimax absolute-error risk \(R(n,m)\) of \cref{obj:def:minimax}, evaluated at these fixed parameters:
\[
R(n,m)\ge c(nN)^{-\gamma/\Delta}.
\]
\end{lemmav}

\begin{proof}[Proof of \cref{obj:lem:nuisance-frontier-converse}]
\begin{enumerate}
\item \textit{Public constants.}
Take the constants supplied by \cref{obj:thm:marked-component-certificate}, writing its amplitude constant as \(c_{\mathrm{cert}}\):
\[
0<c_{\mathrm{cert}}\le1,\qquad c_0>0,\qquad C>0.
\]
They depend only on the fixed public parameters. The parameter conditions give
\[
\gamma>0,\qquad S>0,\qquad \Delta>0.
\]
Define
\[
c_h=\frac14,\qquad
c_\delta=\min\{c_h,c_0\},\qquad
\rho_0=\min\left\{
c_{\mathrm{cert}},1,c_{\mathrm{cert}}c_h^\gamma,
\frac1{16C}
\right\}.
\]
Thus
\[
0<c_\delta\le c_h\le\frac12,\qquad
c_\delta\le c_0,\qquad
0<\rho_0\le1,
\]
and
\[
\rho_0\le c_{\mathrm{cert}},\qquad
\rho_0\le c_{\mathrm{cert}}c_h^\gamma,\qquad
C\rho_0\le\frac1{16}.
\]
Since \(d\ge1\), \(S>0\), and \(c_\delta,c_h\le1\),
\[
c_\delta^d\le c_\delta\le c_0,\qquad
c_\delta^S\le1,\qquad
c_h^d\le1,\qquad
c_\delta^{d+2S}\le1.
\]
Also,
\[
\rho_0^2\le\rho_0,\qquad \rho_0^4\le\rho_0,
\]
so the information constant satisfies
\[
C\rho_0^4c_h^dc_\delta^{d+2S}
\le C\rho_0\le\frac1{16}.
\]
Finally, define the separation constant and the claimed lower-bound constant by
\[
c_{\mathrm{sep}}=\rho_0^2c_\delta^S>0,
\qquad
c=\frac34c_{\mathrm{sep}}>0.
\]
% lean: CausalSmith.Stat.TwosamplePointcateAnnotationFrontier.nuisance_scaled_marked_priors

\item \textit{Admissible scales.}
Fix arbitrary \(n,m\in\mathbb N\) in the stated range, and write
\[
N=n+m,\qquad n\ge2,\qquad n\le N\le n^{q_*}.
\]
In particular, \(n,N>0\) and \(nN\ge1\). Choose
\[
h=c_h(nN)^{-1/\Delta},\qquad
\delta=c_\delta(nN)^{-\gamma/(S\Delta)},\qquad
a=\rho_0\delta^\alpha,\qquad
b=\rho_0\delta^\beta.
\]
The frontier condition implies
\[
S<S_{\mathrm{crit}}
=\frac{\gamma d}{2\gamma+d}<\gamma,
\]
because \(2\gamma+d>d>0\). Consequently,
\[
-\frac{\gamma}{S\Delta}\le-\frac1\Delta\le0.
\]
Monotonicity of positive-base powers for \(nN\ge1\) now gives
\[
0<\delta
\le c_\delta(nN)^{-1/\Delta}
\le c_h(nN)^{-1/\Delta}
=h
\le c_h\le\frac12.
\]
The amplitude definitions then give \(a,b>0\).
% lean: CausalSmith.Stat.TwosamplePointcateAnnotationFrontier.nuisance_scales_admissible

\item \textit{Separation and contrast calibration.}
Using \(\delta>0\) and \(S=\alpha+\beta\), multiplication of positive-base powers yields
\[
\begin{aligned}
ab
&=\rho_0^2\delta^{\alpha+\beta}\\
&=\rho_0^2c_\delta^S
  (nN)^{-\gamma S/(S\Delta)}\\
&=c_{\mathrm{sep}}(nN)^{-\gamma/\Delta}.
\end{aligned}
\]
Likewise,
\[
h^\gamma=c_h^\gamma(nN)^{-\gamma/\Delta}.
\]
The bounds \(\rho_0^2\le\rho_0\), \(c_\delta^S\le1\), and \(\rho_0\le c_{\mathrm{cert}}c_h^\gamma\) therefore imply
\[
\begin{aligned}
ab
&=\rho_0^2c_\delta^S(nN)^{-\gamma/\Delta}\\
&\le\rho_0(nN)^{-\gamma/\Delta}\\
&\le c_{\mathrm{cert}}c_h^\gamma(nN)^{-\gamma/\Delta}
=c_{\mathrm{cert}}h^\gamma.
\end{aligned}
\]
% lean: CausalSmith.Stat.TwosamplePointcateAnnotationFrontier.nuisance_scaled_marked_priors

\item \textit{Occupancy, class membership, and mixture distance.}
The definitions of the rate parameters give
\[
q_*=\frac{S_{\mathrm{crit}}}{S}\ge0,\qquad
\Delta=(2\gamma+d)(1+q_*),\qquad
\frac{\gamma d}{S\Delta}=\frac{q_*}{1+q_*}.
\]
Substituting the scale \(\delta\) and collecting powers therefore gives
\[
\begin{aligned}
N\delta^d
&=c_\delta^d
  n^{-q_*/(1+q_*)}N^{1/(1+q_*)}\\
&=c_\delta^d
  \left(\frac{N}{n^{q_*}}\right)^{1/(1+q_*)}
\le c_\delta^d\le c_0.
\end{aligned}
\]
Here the boundary condition makes the positive ratio \(N/n^{q_*}\) at most one, and \(1/(1+q_*)>0\). The remaining amplitude hypotheses follow from \(\rho_0\le c_{\mathrm{cert}}\):
\[
0<a\le c_{\mathrm{cert}}\delta^\alpha,\qquad
0<b\le c_{\mathrm{cert}}\delta^\beta.
\]

All hypotheses of \cref{obj:thm:marked-component-certificate} are thus satisfied for the construction of \cref{obj:def:marked-handle}. For its finite lattice index set \(\mathcal I\), write
\[
\mathcal S_{\mathrm{sign}}
=\{-1,+1\}^{\mathcal I}\times\{-1,+1\}^{\mathcal I}.
\]
The certificate supplies probability priors \(\varpi_\theta\) and, for every
\(\theta\in\{-1,+1\}\), \((\lambda,\eta)\in\mathcal S_{\mathrm{sign}}\), and \(x\in[0,1]^d\),
\[
P_{\theta,\lambda,\eta}\in\mathcal M,\qquad
\tau_{P_{\theta,\lambda,\eta}}(x)=\tau_\theta(x).
\]
At the evaluation profile this gives
\[
\tau_{P_{-1,\lambda,\eta}}(x_0)=2ab,\qquad
\tau_{P_{+1,\lambda,\eta}}(x_0)=-2ab.
\]

Define the mixtures of the original data blocks by
\[
M_\theta
=\int \mathcal L_P(\mathcal D_{n,m})\,d\varpi_\theta(P),
\qquad \theta\in\{-1,+1\}.
\]
The certificate also supplies
\[
H^2(M_{+1},M_{-1})\le CnNh^d\delta^da^2b^2.
\]
To evaluate its right-hand side, first note that
\[
\delta^da^2b^2=\rho_0^4\delta^{d+2S},
\qquad
d+\frac{\gamma(d+2S)}{S}
=d+\frac{\gamma d}{S}+2\gamma=\Delta.
\]
Hence
\[
\begin{aligned}
nNh^d\delta^da^2b^2
&=\rho_0^4c_h^dc_\delta^{d+2S}
  (nN)^{\,1-d/\Delta-\gamma(d+2S)/(S\Delta)}\\
&=\rho_0^4c_h^dc_\delta^{d+2S}.
\end{aligned}
\]
The inequality \(C\rho_0^4c_h^dc_\delta^{d+2S}\le1/16\) proves
\[
H^2(M_{+1},M_{-1})
\le C\rho_0^4c_h^dc_\delta^{d+2S}
\le\frac1{16}.
\]
% lean: CausalSmith.Stat.TwosamplePointcateAnnotationFrontier.nuisance_scaled_marked_priors

\item \textit{The target as a functional of the experiment law.}
The randomized observation space and the marked experiment model are
\[
\Omega_{\mathrm{obs}}
=
\bigl([0,1]^d\times\{0,1\}^2\bigr)^n
\times
\bigl([0,1]^d\times\{0,1\}\bigr)^m
\times[0,1],
\]
\[
\mathcal P_{\mathrm{exp}}
=
\left\{
\mathsf E_{n,m}(P_{\theta,\lambda,\eta}):
\theta\in\{-1,+1\},\
(\lambda,\eta)\in\mathcal S_{\mathrm{sign}}
\right\}.
\]
When \(m=0\), the auxiliary-record factor is the one-point empty product.

We show that equal marked experiment laws have equal point targets. Take two constructed laws
\[
P=P_{\theta,\lambda,\eta},\qquad
P'=P_{\theta',\lambda',\eta'},
\]
where both hypothesis indices belong to \(\{-1,+1\}\) and both sign pairs belong to \(\mathcal S_{\mathrm{sign}}\), and suppose
\[
\mathsf E_{n,m}(P)=\mathsf E_{n,m}(P').
\]
Since \(n>0\), projection onto the first labeled record gives
\[
\mathcal L_P(X,A,Y)=\mathcal L_{P'}(X,A,Y).
\]
Both covariate laws are uniform. Uniqueness of conditional probability kernels therefore gives
\[
\mathcal L_P(A,Y\mid X=x)
=
\mathcal L_{P'}(A,Y\mid X=x)
\quad\text{for Lebesgue-almost every }x\in[0,1]^d.
\]
For the independent Bernoulli construction, with the response selection of \cref{obj:synth_5}, the four conditional cells are
\[
\begin{aligned}
\Pr_P(A=1,Y=1\mid X=x)
 &=e_P(x)\mu_{1,P}(x),\\
\Pr_P(A=1,Y=0\mid X=x)
 &=e_P(x)\bigl(1-\mu_{1,P}(x)\bigr),\\
\Pr_P(A=0,Y=1\mid X=x)
 &=\bigl(1-e_P(x)\bigr)\mu_{0,P}(x),\\
\Pr_P(A=0,Y=0\mid X=x)
 &=\bigl(1-e_P(x)\bigr)\bigl(1-\mu_{0,P}(x)\bigr).
\end{aligned}
\]
The same formulas hold for \(P'\). Summing the first two equal cells identifies the propensity. The equal cells with \(Y=1\), together with overlap, then identify the arm means:
\[
e_P(x)=e_{P'}(x),\qquad
\mu_{0,P}(x)=\mu_{0,P'}(x),\qquad
\mu_{1,P}(x)=\mu_{1,P'}(x)
\quad\text{almost everywhere}.
\]
Indeed, the factors used for division satisfy
\[
e_P(x)\ge\varepsilon>0,\qquad 1-e_P(x)\ge\varepsilon>0.
\]
Conditional independence in the prescribed construction now gives
\[
\begin{aligned}
\mathcal L_P(A,Y(0),Y(1)\mid X=x)
&=\operatorname{Bernoulli}(e_P(x))
  \otimes\operatorname{Bernoulli}(\mu_{0,P}(x))
  \otimes\operatorname{Bernoulli}(\mu_{1,P}(x))\\
&=\mathcal L_{P'}(A,Y(0),Y(1)\mid X=x)
\quad\text{almost everywhere}.
\end{aligned}
\]
Integration against their common uniform covariate law proves equality of the full primitive laws.

Their common primitive law thus gives identical Bernoulli margin laws for \(A,Y(0),Y(1)\), so admissible conditional-margin versions agree almost everywhere. The positive Hölder orders give continuity of \(e_P,\mu_{0,P},\tau_P\), and
\[
\mu_{1,P}=\mu_{0,P}+\tau_P
\quad\text{on }[0,1]^d
\]
gives continuity of the treated mean; the same holds for \(P'\). Since the uniform law has full support on the cube, these continuous versions agree everywhere. In particular,
\[
\tau_P(x_0)=\tau_{P'}(x_0).
\]

We may consequently define a functional on all nonnegative Borel measures on \(\Omega_{\mathrm{obs}}\) by
\[
\Psi(\mathbb P)=
\begin{cases}
\tau_{P_{\theta,\lambda,\eta}}(x_0),
&
\mathbb P=\mathsf E_{n,m}(P_{\theta,\lambda,\eta})
\text{ for some }
\theta\in\{-1,+1\},\
(\lambda,\eta)\in\mathcal S_{\mathrm{sign}},\\
0,&\mathbb P\notin\mathcal P_{\mathrm{exp}}.
\end{cases}
\]
The equality just proved makes the first case independent of the representative. Thus
\[
\Psi\bigl(\mathsf E_{n,m}(P_{\theta,\lambda,\eta})\bigr)
=\tau_{P_{\theta,\lambda,\eta}}(x_0)
\]
throughout the marked family.
% lean: CausalSmith.Stat.TwosamplePointcateAnnotationFrontier.marked_experiment_target_exists

\item \textit{A common prior for testing.}
Use the finite index space
\[
\mathcal Z_{\mathrm{test}}
=\mathcal S_{\mathrm{sign}}\times\mathcal S_{\mathrm{sign}},
\qquad
z_{\mathrm{test}}
=((\lambda_-,\eta_-),(\lambda_+,\eta_+)).
\]
Define the common prior by its point masses:
\[
\omega(\{z_{\mathrm{test}}\})
=
w_{-1}(\lambda_-,\eta_-)\,
w_{+1}(\lambda_+,\eta_+).
\]
The nonnegative weights and their unit sums from \cref{obj:thm:marked-component-certificate} make \(\omega\) a probability measure. Define the two probability kernels by
\[
\mathbb B_{z_{\mathrm{test}}}
=\mathsf E_{n,m}(P_{-1,\lambda_-,\eta_-}),
\qquad
\mathbb C_{z_{\mathrm{test}}}
=\mathsf E_{n,m}(P_{+1,\lambda_+,\eta_+}).
\]
Their values belong to \(\mathcal P_{\mathrm{exp}}\). For every
\(z_{\mathrm{test}},z'_{\mathrm{test}}\in\mathcal Z_{\mathrm{test}}\), the target identities give
\[
\Psi(\mathbb B_{z_{\mathrm{test}}})
-\Psi(\mathbb C_{z'_{\mathrm{test}}})
=2ab-(-2ab)=4ab>0.
\]

Summing out the unused sign pair in each kernel gives the mixture identities
\[
\int\mathbb B_{z_{\mathrm{test}}}\,d\omega(z_{\mathrm{test}})
=M_{-1}\otimes\operatorname{Unif}([0,1]),
\]
\[
\int\mathbb C_{z_{\mathrm{test}}}\,d\omega(z_{\mathrm{test}})
=M_{+1}\otimes\operatorname{Unif}([0,1]).
\]
The uniform factor is the independent randomizer. By the common-factor identity in \cref{obj:lem:hellinger-block-calculus} and symmetry of the squared difference in \cref{obj:def:hellinger},
\[
\begin{aligned}
\mathsf h^2\left(
\int\mathbb B_{z_{\mathrm{test}}}\,d\omega,
\int\mathbb C_{z_{\mathrm{test}}}\,d\omega
\right)
&=H^2(M_{-1},M_{+1})\\
&=H^2(M_{+1},M_{-1})
\le\frac1{16}.
\end{aligned}
\]
% lean: CausalSmith.Stat.TwosamplePointcateAnnotationFrontier.marked_common_prior_hellinger

\item \textit{Testing implies a minimax bound.}
Apply \cref{obj:lem:published-fuzzy-testing} to the model \(\mathcal P_{\mathrm{exp}}\), the functional \(\Psi\), and the preceding kernels and prior, with
\[
v_0=1,\qquad s_0=4ab>0,\qquad
u_0=\frac1{16}\in[0,2).
\]
Here one observation is the entire randomized data block; the one-fold product space is identified with that block by its sole coordinate. For every Borel real-valued decision \(T\) on \(\Omega_{\mathrm{obs}}\), the testing bound gives
\[
\sup_{\mathbb P\in\mathcal P_{\mathrm{exp}}}
\int_{\Omega_{\mathrm{obs}}}
\left|T(w_{\mathrm{obs}})-\Psi(\mathbb P)\right|
\,d\mathbb P(w_{\mathrm{obs}})
\ge
ab\left(1-\frac{\sqrt{63}}{32}\right).
\]
The numerical overlap factor satisfies
\[
63\le64,\qquad
\frac{\sqrt{63}}{32}\le\frac14,\qquad
1-\frac{\sqrt{63}}{32}\ge\frac34.
\]
Moreover, the target representation and \cref{obj:def:risk} give
\[
\begin{aligned}
&\sup_{\mathbb P\in\mathcal P_{\mathrm{exp}}}
\int_{\Omega_{\mathrm{obs}}}
\left|T(w_{\mathrm{obs}})-\Psi(\mathbb P)\right|
\,d\mathbb P(w_{\mathrm{obs}})\\
&\qquad=
\sup_{\substack{\theta\in\{-1,+1\}\\
(\lambda,\eta)\in\mathcal S_{\mathrm{sign}}}}
\mathcal R_{n,m}(T,P_{\theta,\lambda,\eta})
\le
\sup_{P\in\mathcal M}\mathcal R_{n,m}(T,P),
\end{aligned}
\]
where the last inequality uses the certified class membership. Thus every such decision satisfies
\[
\sup_{P\in\mathcal M}\mathcal R_{n,m}(T,P)\ge\frac34ab.
\]
Taking the infimum over the decisions in \cref{obj:def:minimax} proves
\[
R(n,m)\ge\frac34ab.
\]
All loss integrals and suprema here retain their extended nonnegative interpretation.
% lean: CausalSmith.Stat.TwosamplePointcateAnnotationFrontier.marked_fuzzy_testing_minimax

\item \textit{Substitution of the separation scale.}
The identity \(ab=c_{\mathrm{sep}}(nN)^{-\gamma/\Delta}\) now yields
\[
R(n,m)
\ge\frac34ab
=\frac34c_{\mathrm{sep}}(nN)^{-\gamma/\Delta}
=c(nN)^{-\gamma/\Delta}.
\]
The positive constant \(c\) depends only on the fixed public parameters, and the sample sizes were arbitrary in the stated range.
% lean: CausalSmith.Stat.TwosamplePointcateAnnotationFrontier.nuisance_frontier_converse
\end{enumerate}
\end{proof}

Under the low-nuisance-smoothness condition in \cref{obj:lem:nuisance-frontier-converse}, the exponent \(q_*\) exceeds one. The stated interval consequently covers total record counts from the fully labeled experiment through the crossover to supplied-propensity precision. Within this interval, the lower bound captures the joint contribution of labeled and outcome-free records through their count product.

The complete minimax characterization in \cref{obj:thm:sharp-annotation-frontier} organizes the lower bounds into complementary ranges. When \(S<S_{\mathrm{crit}}\) and \(n\le N\le n^{q_*}\), \cref{obj:lem:nuisance-frontier-converse} supplies the unequal-count term. Supplying the propensity enlarges the information available to a decision, so the supplied-propensity risk in \cref{obj:def:oracle-risk} is a lower bound for the original-record minimax risk in \cref{obj:def:minimax}. For \(N\ge n^{q_*}\), the supplied-propensity lower bound in \cref{obj:thm:supplied-propensity} supplies the order \(n^{-\gamma/(2\gamma+d)}\), which is the larger term in the sharp rate. When \(S\ge S_{\mathrm{crit}}\), that same order is the larger term for every \(N\ge n\). Thus the supplied-propensity floor resolves the ranges complementary to the sparse-occupancy construction.

Together with \cref{obj:thm:rectangular-upper}, these lower bounds yield the matching characterization in \cref{obj:thm:sharp-annotation-frontier}. For the formal crosswalk, the rate and attaining decision are collected in the ordered pair of \cref{obj:def:frontier-handle} in \cref{sec:verification-note}.



The statistical interpretation and count consequences of this pair are stated once in \cref{obj:thm:sharp-annotation-frontier,obj:prop:annotation-threshold}.


\section{Acquisition consequences and supervised comparison}\label{sec:acquisition-consequences}

The sample requirements in \cref{obj:prop:annotation-threshold} follow from the two terms of the sharp rate. Throughout the acquisition analysis, fix \(d\ge1\), \(\alpha,\beta\in(0,1]\), \(\gamma\in[1,3]\), \(L>1/2\), and an overlap level \(\varepsilon\in(0,1/2)\), as in \cref{obj:thm:sharp-annotation-frontier,obj:thm:supplied-propensity}. Retain the total treatment-record count \(N=n+m\), summed nuisance smoothness \(S=\alpha+\beta\), rate denominator \(\Delta\), smoothness threshold \(S_{\mathrm{crit}}\), and total-record exponent \(q_*\) defined there. The identities below express how labeled responses and treatment records jointly determine precision.

\subsection{Complete acquisition statement}

\begin{propositionv}{prop:annotation-threshold}[Annotation thresholds and acquisition requirements]
Fix public parameters satisfying
\begin{itemize}
\item \textbf{(Dimension.)} \(d\in\mathbb N\) and \(d\ge1\).
\item \textbf{(Nuisance smoothness.)} \(0<\alpha\le1\) and \(0<\beta\le1\).
\item \textbf{(Effect smoothness.)} \(1\le\gamma\le3\).
\item \textbf{(Radius.)} \(L>1/2\).
\item \textbf{(Overlap.)} \(0<\varepsilon<1/2\).
\end{itemize}
Let \(R(n,m)\) and \(R^e(n,m)\) be the minimax and supplied-propensity risks of \cref{obj:def:minimax,obj:def:oracle-risk}, and let \(r_*(n,m)\) be the sharp rate of \cref{obj:def:sharp-rate}, with the fixed parameter arguments suppressed. Define the total treatment-record count \(N\), the summed nuisance smoothness \(S\), the tuning denominator \(\Delta\), the smoothness cutoff \(S_{\mathrm{crit}}\), the total-record exponent \(q_*\), and the oracle rate \(r_o\) by
\[
N=n+m,\qquad S=\alpha+\beta,\qquad
\Delta=2\gamma+d+\frac{\gamma d}{S},
\qquad
S_{\mathrm{crit}}=\frac{\gamma d}{2\gamma+d},
\qquad
q_*=\frac{\gamma d}{S(2\gamma+d)},
\qquad
r_o(n)=n^{-\gamma/(2\gamma+d)}.
\]
All constants below may depend on the fixed public parameters.

For every sequence \(m_n\in\mathbb N\), the following equivalences hold:
\[
\begin{aligned}
\frac{R(n,m_n)}{R^e(n,m_n)}=O(1)
&\quad\Longleftrightarrow\quad
r_*(n,m_n)=O(r_o(n))\\
&\quad\Longleftrightarrow\quad
\exists\,\zeta\in[1,\infty)\\ 
\forall\,n\ \text{ sufficiently large},\\
(n,m_n)\in\mathcal B_*(\zeta)\\
&\quad\Longleftrightarrow\quad
\liminf_{n\to\infty}\frac{n+m_n}{n^{q_*}}>0,
\end{aligned}
\]
where \(\mathcal B_*(\zeta)\) is the region of \cref{obj:def:annotation-region}; the lower limit is understood in the extended nonnegative reals. Moreover,
\[
\begin{aligned}
\frac{R(n,m_n)}{R(n,0)}\longrightarrow0
&\quad\Longleftrightarrow\quad
\frac{r_*(n,m_n)}{r_*(n,0)}\longrightarrow0\\
&\quad\Longleftrightarrow\quad
S<S_{\mathrm{crit}}
\ \text{ and }\\ 
\frac{m_n}{n}\longrightarrow\infty.
\end{aligned}
\]
All sequence limits and order relations are as \(n\to\infty\). If \(S<S_{\mathrm{crit}}\), then \(q_*>1\). For every sequence with \(m_n/n=O(1)\),
\[
R(n,m_n)=O(R(n,0)),
\qquad
R(n,0)=O(R(n,m_n)).
\]
If \(S\ge S_{\mathrm{crit}}\), then every auxiliary-size sequence satisfies
\[
\frac{R(n,m_n)}{R^e(n,m_n)}=O(1).
\]

There exist real constants \(c,C\) with \(0<c\le C\) such that, for every \(n,m\in\mathbb N\) with \(n\ge2\),
\[
c\,r_*(n,m)\le R(n,m)\le C\,r_*(n,m).
\]
For these sample sizes, the exact rate branches are
\[
r_*(n,m)=
\begin{cases}
r_o(n),&S\ge S_{\mathrm{crit}},\\
(nN)^{-\gamma/\Delta},
&S<S_{\mathrm{crit}},\ N\le n^{q_*},\\
r_o(n),
&S<S_{\mathrm{crit}},\ N\ge n^{q_*}.
\end{cases}
\]
At every boundary pair with \(N=n^{q_*}\),
\[
(nN)^{-\gamma/\Delta}=r_o(n).
\]
In particular, for every \(n\ge2\),
\[
r_*(n,0)=
\begin{cases}
n^{-2\gamma/\Delta},&S<S_{\mathrm{crit}},\\
r_o(n),&S\ge S_{\mathrm{crit}}.
\end{cases}
\]
Writing \(s_{\mathrm{pub}}=S/2\) for the average nuisance smoothness, the numerical supervised benchmark coincides with \(r_*(n,0)\): for every \(n\in\mathbb N\),
\[
\begin{cases}
n^{-1/(1+d/(2\gamma)+d/(4s_{\mathrm{pub}}))},
&s_{\mathrm{pub}}<\dfrac{d/4}{1+d/(2\gamma)},\\
n^{-1/(2+d/\gamma)},
&s_{\mathrm{pub}}\ge\dfrac{d/4}{1+d/(2\gamma)}
\end{cases}
=r_*(n,0).
\]
In this identity at \(n=0\), negative powers of zero are assigned the value zero.

Finally, there exist real constants \(c,C>0\), separate from the preceding pair, such that the following acquisition guarantees hold for every target absolute-error scale \(0<\eta<1\) and every \(n,m\in\mathbb N\) with \(n\ge2\). Let \(\mathcal M\) be the primitive class at the fixed parameters from \cref{obj:def:primitive-class}, and use the risk of \cref{obj:def:risk}.

If there exists a Borel decision \(T\) satisfying
\[
\mathcal R_{n,m}(T,P)\le\eta
\qquad\text{for every }P\in\mathcal M,
\]
then
\[
c\,\eta^{-(2+d/\gamma)}\le n,
\qquad
c\,\eta^{-(2+d/\gamma+d/S)}\le n(n+m).
\]
Conversely, if
\[
C\,\eta^{-(2+d/\gamma)}\le n,
\qquad
C\,\eta^{-(2+d/\gamma+d/S)}\le n(n+m),
\]
then the sharp decision \(T_*\) of \cref{obj:def:sharp-decision} is Borel and satisfies
\[
\mathcal R_{n,m}(T_*,P)\le\eta
\qquad\text{for every }P\in\mathcal M.
\]
The exact rate-level acquisition condition is
\[
r_*(n,m)\le\eta
\quad\Longleftrightarrow\quad
\eta^{-(2+d/\gamma)}\le n
\ \text{ and }\
\eta^{-(2+d/\gamma+d/S)}\le n(n+m).
\]
The total treatment-record count also satisfies \(n+m\ge n\).
\end{propositionv}

\subsection{Oracle attainment and strict benefit}

Since \(\Delta=(2\gamma+d)(1+q_*)\), the sharp rate in \cref{obj:def:sharp-rate} satisfies
\[
\frac{r_*(n,m)}{r_o(n)}
=
\max\left\{
1,\,
\left(\frac{n^{q_*}}{N}\right)^{\gamma/\Delta}
\right\}.
\]
Thus, for any finite tolerance \(\zeta\ge1\), the annotation region in \cref{obj:def:annotation-region} has the exact characterization
\[
(n,m)\in\mathcal B_*(\zeta)
\quad\Longleftrightarrow\quad
N\ge \zeta^{-\Delta/\gamma}n^{q_*},
\qquad n\ge2,\quad m\ge0.
\]
The tolerance measures the remaining nuisance contribution relative to the supplied-propensity order. At \(\zeta=1\), the count threshold is the point at which the two terms coincide.

For an auxiliary-size sequence \((m_n)\), write \(N_n=n+m_n\). The uniform risk comparisons in \cref{obj:thm:sharp-annotation-frontier,obj:thm:supplied-propensity} give
\[
\frac{R(n,m_n)}{R^e(n,m_n)}
\asymp
\frac{r_*(n,m_n)}{r_o(n)}.
\]
Consequently, bounded relative risk is equivalent to an eventual positive lower bound on \(N_n/n^{q_*}\), or, equivalently,
\[
\frac{R(n,m_n)}{R^e(n,m_n)}=O(1)
\quad\Longleftrightarrow\quad
\liminf_{n\to\infty}\frac{N_n}{n^{q_*}}>0.
\]
This is the oracle-attainment equivalence in \cref{obj:prop:annotation-threshold}. When \(S\ge S_{\mathrm{crit}}\), the inequality \(q_*\le1\), together with \(N_n\ge n\), places every auxiliary-size sequence in this regime.

Strict benefit compares precision with labeled records alone. In the regime \(S<S_{\mathrm{crit}}\), one has \(q_*>1\) and \(r_*(n,0)=n^{-2\gamma/\Delta}\). Hence
\[
\frac{r_*(n,m_n)}{r_*(n,0)}
=
\max\left\{
n^{-\gamma(q_*-1)/\Delta},\,
\left(\frac{n}{N_n}\right)^{\gamma/\Delta}
\right\}.
\]
The first term tends to zero, and the second tends to zero exactly when \(m_n/n\to+\infty\). For \(S\ge S_{\mathrm{crit}}\), the rate ratio equals one. The uniform minimax comparisons therefore yield the criterion in \cref{obj:prop:annotation-threshold}:
\[
\frac{R(n,m_n)}{R(n,0)}\longrightarrow0
\quad\Longleftrightarrow\quad
S<S_{\mathrm{crit}}
\quad\text{and}\quad
\frac{m_n}{n}\longrightarrow+\infty.
\]
Strict benefit and oracle attainment thus describe distinct acquisition requirements. In the low-smoothness regime, a diverging auxiliary-to-labeled ratio produces a strict improvement, while attainment of the supplied-propensity order requires total treatment information of order at least \(n^{q_*}\). A bounded auxiliary-to-labeled ratio preserves the supervised minimax order.

\subsection{Target-error acquisition}

For a target absolute-error scale \(0<\eta<1\), inversion of the two rate terms gives the exact identity
\[
r_*(n,m)\le\eta
\quad\Longleftrightarrow\quad
\begin{cases}
n\ge \eta^{-(2+d/\gamma)},\\
nN\ge \eta^{-(2+d/\gamma+d/S)}.
\end{cases}
\]
The first inequality specifies the response-label requirement. The second specifies the joint requirement on labeled and treatment-record counts.

To connect these numerical conditions to risk, let \(c_R,C_R>0\) be comparison constants from \cref{obj:thm:sharp-annotation-frontier}, so that
\[
c_Rr_*(n,m)
\le R(n,m)
\le \sup_{P\in\mathcal M}\mathcal R_{n,m}(T_*,P)
\le C_Rr_*(n,m).
\]
A Borel decision whose risk is at most \(\eta\) uniformly over \(\mathcal M\) necessarily satisfies
\[
n\ge (c_R/\eta)^{2+d/\gamma},
\qquad
nN\ge (c_R/\eta)^{2+d/\gamma+d/S}.
\]
Conversely, the sufficient conditions
\[
n\ge (C_R/\eta)^{2+d/\gamma},
\qquad
nN\ge (C_R/\eta)^{2+d/\gamma+d/S}
\]
ensure that the decision \(T_*\) in \cref{obj:def:sharp-decision} has uniform risk at most \(\eta\). Absorbing the fixed powers of these comparison constants gives the necessity and sufficiency constants in \cref{obj:prop:annotation-threshold}.

At the rate level, a fixed labeled count satisfying the first inequality requires
\[
N\ge
\max\left\{
n,\,
\frac{\eta^{-(2+d/\gamma+d/S)}}{n}
\right\}.
\]
In particular, choosing \(n\asymp\eta^{-(2+d/\gamma)}\) gives the total-record requirement
\[
N\gtrsim
\max\left\{
\eta^{-(2+d/\gamma)},\,
\eta^{-d/S}
\right\}.
\]
The smoothness threshold compares these two powers. Below \(S_{\mathrm{crit}}\), the treatment-record requirement grows faster than the labeled requirement; at and above the threshold, the labeled count already supplies the required treatment information.

\subsection{The supervised numerical benchmark}

At \(m=0\), the sharp rate becomes
\[
r_*(n,0)
=
\max\left\{
n^{-\gamma/(2\gamma+d)},\,
n^{-2\gamma/\Delta}
\right\}.
\]
Its exponents and smoothness cutoff coincide with the numerical benchmark in \citet[Section 3, Theorem 1]{KennedyBalakrishnanRobinsWasserman2024}.

The cited theorem concerns a supervised model with a covariate density bounded above by a constant, an \(\alpha\)-smooth propensity, a \(\beta\)-smooth control regression, and a \(\gamma\)-smooth contrast. It gives a minimax mean-absolute-error lower bound for sample sizes larger than a constant depending on \((\alpha,\beta,\gamma,d)\). Its implicit multiplicative constant is positive and independent of sample size, with the density bound and H\"older constants fixed; these are source-model constants. The density restriction accommodates the construction-specific designs used in the source testing argument \citep[Section 3.1]{KennedyBalakrishnanRobinsWasserman2024}. The following identities express the numerical comparison in the notation of this paper.

\begin{lemmav}{lem:published-cate-benchmark}[Published CATE benchmark identities]
Let the dimension \(d\in\mathbb N\) and smoothness orders \(\alpha,\beta,\gamma\in\mathbb R\) satisfy:
\begin{itemize}
\item \textbf{(Dimension.)} \(1\le d\).
\item \textbf{(First nuisance smoothness.)} \(0<\alpha\).
\item \textbf{(Second nuisance smoothness.)} \(0<\beta\).
\item \textbf{(Effect smoothness.)} \(1\le\gamma\).
\end{itemize}
Put \(s_{\mathrm{pub}}=(\alpha+\beta)/2\), the average nuisance smoothness. The numerical exponents and smoothness cutoff displayed in \citet[Section 3, Theorem 1]{KennedyBalakrishnanRobinsWasserman2024} satisfy
\[
\frac{1}{1+d/(2\gamma)+d/(4s_{\mathrm{pub}})}
=\frac{2\gamma}{2\gamma+d+\gamma d/(\alpha+\beta)},
\qquad
\frac{1}{2+d/\gamma}=\frac{\gamma}{2\gamma+d},
\]
and
\[
s_{\mathrm{pub}}<\frac{d/4}{1+d/(2\gamma)}
\quad\Longleftrightarrow\quad
\alpha+\beta<\frac{\gamma d}{2\gamma+d}.
\]
For every \(n\in\mathbb N\) with \(n\ge2\), the numerical piecewise benchmark satisfies
\[
\begin{cases}
n^{-1/(1+d/(2\gamma)+d/(4s_{\mathrm{pub}}))},
& s_{\mathrm{pub}}<(d/4)/(1+d/(2\gamma)),\\
n^{-1/(2+d/\gamma)},
& s_{\mathrm{pub}}\ge(d/4)/(1+d/(2\gamma))
\end{cases}
=
\max\left\{
n^{-\gamma/(2\gamma+d)},
n^{-2\gamma/(2\gamma+d+\gamma d/(\alpha+\beta))}
\right\}.
\]
\end{lemmav}

\begin{proof}[Proof of \cref{obj:lem:published-cate-benchmark}]
\begin{enumerate}
\item The hypotheses give
\[
d>0,\qquad \gamma>0,\qquad \alpha+\beta>0.
\]
Define the real denominators
\[
D_{\mathrm{or}}:=2\gamma+d,
\qquad
D_{\mathrm{pair}}:=2\gamma+d+\frac{\gamma d}{\alpha+\beta}.
\]
Both are strictly positive:
\[
D_{\mathrm{or}}>0,\qquad D_{\mathrm{pair}}>0.
\]
% lean: CausalSmith.Stat.TwosamplePointcateAnnotationFrontier.published_cate_benchmark

\item Clearing the positive denominators and simplifying gives
\[
\frac{1}{1+d/(2\gamma)+d/(4s_{\mathrm{pub}})}
=\frac{2\gamma}{D_{\mathrm{pair}}},
\qquad
\frac{1}{2+d/\gamma}
=\frac{\gamma}{D_{\mathrm{or}}}.
\]
Likewise,
\[
\frac{d/4}{1+d/(2\gamma)}
=\frac{1}{2}\frac{\gamma d}{D_{\mathrm{or}}}.
\]
Since \(s_{\mathrm{pub}}=(\alpha+\beta)/2\), multiplication by \(2>0\) therefore yields
\[
s_{\mathrm{pub}}<\frac{d/4}{1+d/(2\gamma)}
\quad\Longleftrightarrow\quad
\alpha+\beta<\frac{\gamma d}{D_{\mathrm{or}}}.
\]
These are the asserted exponent identities and cutoff equivalence.
% lean: CausalSmith.Stat.TwosamplePointcateAnnotationFrontier.published_cate_benchmark

\item Fix \(n\in\mathbb N\) with \(n\ge2\), so that
\[
n\ge1.
\]
After substituting the exponent identities, the numerical piecewise benchmark becomes
\[
\begin{cases}
n^{-2\gamma/D_{\mathrm{pair}}},
& s_{\mathrm{pub}}<(d/4)/(1+d/(2\gamma)),\\
n^{-\gamma/D_{\mathrm{or}}},
& s_{\mathrm{pub}}\ge(d/4)/(1+d/(2\gamma)).
\end{cases}
\]
We compare these two powers by splitting at the equivalent cutoff in the preceding step.
% lean: CausalSmith.Stat.TwosamplePointcateAnnotationFrontier.published_cate_benchmark

\item Suppose first that
\[
\alpha+\beta<\frac{\gamma d}{D_{\mathrm{or}}}.
\]
Multiplying by \(D_{\mathrm{or}}>0\), then dividing by \(\alpha+\beta>0\), gives
\[
(\alpha+\beta)D_{\mathrm{or}}<\gamma d,
\qquad
D_{\mathrm{or}}<\frac{\gamma d}{\alpha+\beta}.
\]
Consequently,
\[
2D_{\mathrm{or}}
\le D_{\mathrm{or}}+\frac{\gamma d}{\alpha+\beta}
=D_{\mathrm{pair}}.
\]
Multiplying by \(\gamma>0\) and dividing by the positive product
\(D_{\mathrm{pair}}D_{\mathrm{or}}\) yields
\[
\frac{2\gamma}{D_{\mathrm{pair}}}
\le\frac{\gamma}{D_{\mathrm{or}}},
\qquad
-\frac{\gamma}{D_{\mathrm{or}}}
\le-\frac{2\gamma}{D_{\mathrm{pair}}}.
\]
For a base \(n\ge1\), real powers are nondecreasing in their exponent. Hence
\[
n^{-\gamma/D_{\mathrm{or}}}
\le n^{-2\gamma/D_{\mathrm{pair}}},
\]
and thus
\[
\max\left\{
n^{-\gamma/D_{\mathrm{or}}},
n^{-2\gamma/D_{\mathrm{pair}}}
\right\}
=n^{-2\gamma/D_{\mathrm{pair}}}.
\]
The cutoff equivalence selects precisely this first branch of the piecewise benchmark.
% lean: CausalSmith.Stat.TwosamplePointcateAnnotationFrontier.published_cate_benchmark

\item In the complementary case,
\[
\frac{\gamma d}{D_{\mathrm{or}}}\le\alpha+\beta.
\]
Multiplication by \(D_{\mathrm{or}}>0\), followed by division by
\(\alpha+\beta>0\), gives
\[
\gamma d\le(\alpha+\beta)D_{\mathrm{or}},
\qquad
\frac{\gamma d}{\alpha+\beta}\le D_{\mathrm{or}}.
\]
It follows that
\[
D_{\mathrm{pair}}
=D_{\mathrm{or}}+\frac{\gamma d}{\alpha+\beta}
\le2D_{\mathrm{or}}.
\]
Multiplying by \(\gamma>0\) and dividing by
\(D_{\mathrm{or}}D_{\mathrm{pair}}>0\) now gives
\[
\frac{\gamma}{D_{\mathrm{or}}}
\le\frac{2\gamma}{D_{\mathrm{pair}}},
\qquad
-\frac{2\gamma}{D_{\mathrm{pair}}}
\le-\frac{\gamma}{D_{\mathrm{or}}}.
\]
Monotonicity of real powers for \(n\ge1\) therefore implies
\[
n^{-2\gamma/D_{\mathrm{pair}}}
\le n^{-\gamma/D_{\mathrm{or}}},
\]
so
\[
\max\left\{
n^{-\gamma/D_{\mathrm{or}}},
n^{-2\gamma/D_{\mathrm{pair}}}
\right\}
=n^{-\gamma/D_{\mathrm{or}}}.
\]
The cutoff equivalence selects the second branch, including equality at the cutoff. Substituting the definitions of \(D_{\mathrm{or}}\) and \(D_{\mathrm{pair}}\) proves the asserted maximum identity in both cases.
% lean: CausalSmith.Stat.TwosamplePointcateAnnotationFrontier.published_cate_benchmark
\end{enumerate}
\end{proof}

\Cref{obj:lem:published-cate-benchmark} identifies both the exponents and the boundary between their regimes. Equality at the cutoff places the supervised benchmark at the supplied-propensity order. Under the present public-parameter restrictions, its maximum is exactly \(r_*(n,0)\).

For the known-uniform primitive class in \cref{obj:def:primitive-class}, the minimax interpretation of this expression follows from the original-record lower bounds. The construction in \cref{obj:thm:marked-component-certificate} and the converse in \cref{obj:lem:nuisance-frontier-converse} provide the nuisance contribution in its governing range. The same-class subfamily with \(e_P=\mu_{0,P}=1/2\) in \cref{obj:thm:supplied-propensity} provides the labeled-response floor. Together with \cref{obj:thm:rectangular-upper}, these results establish the sharp risk comparison used throughout the acquisition analysis, while \cref{obj:lem:published-cate-benchmark} supplies the exponent-and-cutoff identity with the published supervised benchmark.


\section{Verification note}\label{sec:verification-note}

The formal results checked in Lean 4 are the estimator guarantee in \cref{obj:thm:rectangular-upper}, the supplied-propensity benchmark in \cref{obj:thm:supplied-propensity}, the smooth-frame construction certificate in \cref{obj:thm:marked-component-certificate}, and the nuisance lower bound in \cref{obj:lem:nuisance-frontier-converse}. Their combination yields the minimax sandwich in \cref{obj:thm:sharp-annotation-frontier} and the acquisition consequences in \cref{obj:prop:annotation-threshold}. Auxiliary checked results cover local polynomial projection, population positivity, rectangular sampling error, and Hellinger calculus in \cref{obj:lem:local-polynomial-projection-facts,obj:lem:population-positivity,obj:lem:rectangular-sampling-bound,obj:lem:hellinger-block-calculus}.

The companion artifact for this paper version, whose working-paper number and version appear in the first-page stamp, is available from \url{https://causalsmith.org/papers/stat_twosample_pointcate_annotation_frontier_v1}. Its interactive formal links provide the theorem-to-file crosswalk and exact source excerpts. In the source archive, the checked modules are rooted at \texttt{CausalSmith/Stat/STAT\_TwosamplePointcateAnnotationFrontier\_Research}; the generated file \texttt{presentation\_crosswalk.json} records the mapping from every displayed result to its Lean module and declaration. The crosswalk is regenerated with each paper version, so it describes the results as printed in this version.

The primitive-model and sampling restrictions are assumed inputs. These comprise the known uniform covariate design, conditional exchangeability, a fixed public overlap level \(\varepsilon\in(0,1/2)\), binary potential responses, H\"older restrictions, and probability-valued arm-mean bounds defining \cref{obj:def:primitive-class}, together with the independent product experiment and randomizer in \cref{obj:ass:sample-law}. Each guarantee retains its stated public-parameter, sample-size, scale, amplitude, and occupancy conditions.

The absolute-loss testing result in \cref{obj:lem:published-fuzzy-testing} is proved in the formal artifact without external formal assumptions; \citet[Section 3, Lemma 1]{KennedyBalakrishnanRobinsWasserman2024} remains its bibliographic source and statistical motivation. For the supervised comparison, \cref{obj:lem:published-cate-benchmark} checks the numerical exponent identities, smoothness cutoff, and piecewise-rate equality associated with \citet[Section 3, Theorem 1]{KennedyBalakrishnanRobinsWasserman2024}. The published statistical lower bound belongs to that source's model and assumptions. The minimax interpretation for the present known-uniform primitive class is established by \cref{obj:thm:sharp-annotation-frontier}, through this paper's estimator guarantee and original-record lower bounds. The formal numerical identity also assigns totalized negative powers the value zero at \(n=0\); every statistical risk and acquisition guarantee in the paper uses \(n\ge2\), so this representation value has no precision interpretation.

For the machine-readable crosswalk, the verified rate and decision are also packaged in the following representation-level handle; no statistical argument depends on this packaging.

\begin{definitionv}{def:frontier-handle}[Rate and decision handle]
The frontier handle is the ordered pair \((r_*,T_*)\), parameterized by \(d,n,m\in\mathbb N\) and \(\alpha,\beta,\gamma,\varepsilon\in\mathbb R\), where \(\varepsilon\) is the overlap parameter. Its components are the sharp rate and decision of \cref{obj:def:sharp-rate,obj:def:sharp-decision}:
\[
r_*=
\max\left\{
n^{-\gamma/(2\gamma+d)},
\bigl(n(n+m)\bigr)^{-\gamma/(2\gamma+d+\gamma d/(\alpha+\beta))}
\right\},
\qquad
T_*(\mathcal D_{n,m},U)
=
\widehat T_{\mathrm{up}}(\mathcal D_{n,m}).
\]
Here \(\mathcal D_{n,m}\) is the original dataset and \(U\in\mathbb R\) is the randomizer coordinate. The decision component uses the publicly tuned estimator of \cref{obj:def:tuning} with parameters \(d,\alpha,\beta,\gamma,\varepsilon,n,m\).
\end{definitionv}


\section{Proofs of the main results}\label{sec:deferred-proofs}

\begin{proof}[Proof of \cref{obj:thm:rectangular-upper}]
Fix the public parameters in the statement. Every constant constructed below depends only on these parameters.

\begin{enumerate}
\item Fix admissible \(n,m,h,J\) and \(P\in\mathcal M\). Define
\[
N=n+m,\qquad S=\alpha+\beta,\qquad
\delta=\frac hJ,\qquad
q=\binom{d+2}{2},\qquad k=qJ^d.
\]
In particular, \(N\ge n\ge2\) and \(\delta>0\).
By \cref{obj:lem:local-polynomial-projection-facts}, there is a constant
\(0<C_{\mathrm{loc}}<\infty\), uniform over these choices, and a coefficient vector
\(\vartheta\in\mathbb R^q\) such that
\[
p_\vartheta(x)=r(x)^\top\vartheta,\qquad
\sup_{x\in C_h}|\tau_P(x)-p_\vartheta(x)|
 \le C_{\mathrm{loc}}h^\gamma,\qquad
\|\vartheta\|_2\le C_{\mathrm{loc}},
\qquad
r_0^\top\vartheta=\tau_P(x_0).
\]
The same result supplies
\[
\|r_0\|_2\le C_{\mathrm{loc}},\qquad
\|e_Pr_u-\Pi_J(e_Pr_u)\|_{L_2(\nu_h)}
 \le C_{\mathrm{loc}}\delta^\alpha,\qquad
\|\mu_{0,P}-\Pi_J\mu_{0,P}\|_{L_2(\nu_h)}
 \le C_{\mathrm{loc}}\delta^\beta.
\]
Here \(x_0\in C_h\), so the bound on \(r_0\) follows from the stated uniform coarse-basis bound.

The finite Borel basis functions and the finite moment sums in
\cref{obj:def:moments} make both empirical moments finite Borel functions of the dataset.
The eigenvalue cutoff is Borel: its matrix superlevel set is the intersection of the closed unit-vector quadratic-form inequalities. Matrix inversion on the invertible branch and clipping are Borel. Thus the decision in
\cref{obj:def:estimator}, with value zero on the failed-cutoff branch, is finite and Borel on every dataset. Composing with the dataset coordinate gives a Borel decision on the original dataset and randomizer.%
% lean: CausalSmith.Stat.TwosamplePointcateAnnotationFrontier.tHat_measurable

\item Define the population moments and the localized Taylor remainder by
\[
\bar R=\mathbb E_{\mathsf E_{n,m}(P)}\widehat R_{h,J},
\qquad
Q_{P,h,J}=\mathbb E_{\mathsf E_{n,m}(P)}\widehat Q_{h,J},
\qquad
g_{\mathrm{rem}}:C_h\longrightarrow\mathbb R,\quad
g_{\mathrm{rem}}(x)=\tau_P(x)-p_\vartheta(x).
\]
All localized integrals below are over \(C_h\).

The observed-response identity and the exchangeability condition in
\cref{obj:def:primitive-class} give, almost surely,
\[
\mathbb E_P[AY\mid X]=e_P(X)\mu_{1,P}(X),
\qquad
\mathbb E_P[Y\mid X]
 =\mu_{0,P}(X)+e_P(X)\tau_P(X).
\]
Indeed, \(AY=AY(1)\), and conditional exchangeability factors its conditional mean. Applying the same argument to \((1-A)Y(0)\) and adding gives the second identity.

The role counts in \cref{obj:def:roles} satisfy
\[
\ell=\lfloor n/2\rfloor\ge1,\qquad
t=N-\ell\ge1.
\]
Their records are independent under the sampling law in
\cref{obj:ass:sample-law}. Consequently, taking expectations in
\cref{obj:def:moments}, using uniform design and factoring each rectangular product, yields
\[
\begin{aligned}
\bar R_u
&=\int r_u e_P\mu_{1,P}\,d\nu_h
 -\int e_Pr_u\,\Pi_J(\mu_{0,P}+e_P\tau_P)\,d\nu_h,\\
(Q_{P,h,J})_{uv}
&=\int r_u e_Pr_v\,d\nu_h
 -\int e_Pr_u\,\Pi_J(e_Pr_v)\,d\nu_h.
\end{aligned}
\]
The response in the rectangular correction is \(Y\), so its conditional mean is the second conditional mean displayed above.

For completeness, writing the orthonormal fine basis from
\cref{obj:def:projection} as
\[
\mathbf b_J=(b_{J,\iota})_{\iota=1}^k,
\]
its finite expansion gives, for arbitrary
\(f_1,f_2\in L_2(\nu_h)\),
\[
\int f_1\Pi_Jf_2\,d\nu_h
=
\sum_{\iota=1}^k
 \left(\int b_{J,\iota} f_1\,d\nu_h\right)
 \left(\int b_{J,\iota} f_2\,d\nu_h\right)
=
\int (\Pi_Jf_1)f_2\,d\nu_h.
\]
Orthonormality also gives
\[
\int f_1f_2\,d\nu_h-\int f_1\Pi_Jf_2\,d\nu_h
=
\int(f_1-\Pi_Jf_1)(f_2-\Pi_Jf_2)\,d\nu_h,
\]
and
\[
\|\Pi_Jf_1\|_{L_2(\nu_h)}^2
+\|f_1-\Pi_Jf_1\|_{L_2(\nu_h)}^2
=\|f_1\|_{L_2(\nu_h)}^2.
\]
These identities show that the population matrix displayed above is symmetric, so it agrees with the expectation of the symmetrized empirical matrix.

Multiplying its entries by \(\vartheta\), substituting
\(\mu_{1,P}=\mu_{0,P}+\tau_P\), and using the residual identity gives
\[
\begin{aligned}
\bar R_u-(Q_{P,h,J}\vartheta)_u
&=\int
 \bigl(e_Pr_u-\Pi_J(e_Pr_u)\bigr)
 \bigl(\mu_{0,P}-\Pi_J\mu_{0,P}\bigr)\,d\nu_h\\
&\quad+\int r_u e_Pg_{\mathrm{rem}}\,d\nu_h
 -\int \Pi_J(e_Pr_u)e_Pg_{\mathrm{rem}}\,d\nu_h.
\end{aligned}
\]
Cauchy--Schwarz and the two residual bounds from the preceding step imply
\[
\left|\int
 \bigl(e_Pr_u-\Pi_J(e_Pr_u)\bigr)
 \bigl(\mu_{0,P}-\Pi_J\mu_{0,P}\bigr)\,d\nu_h\right|
\le C_{\mathrm{loc}}^2\delta^S.
\]
The coarse coordinates are orthonormal, overlap bounds \(|e_P|\) by one, and projection contraction follows from the displayed Pythagorean identity. Hence
\[
\|r_u\|_{L_2(\nu_h)}=1,\qquad
\|e_Pr_u\|_{L_2(\nu_h)}\le1,\qquad
\|\Pi_J(e_Pr_u)\|_{L_2(\nu_h)}\le1,\qquad
\|e_Pg_{\mathrm{rem}}\|_{L_2(\nu_h)}
 \le C_{\mathrm{loc}}h^\gamma.
\]
Applying Cauchy--Schwarz to the remaining two integrals therefore gives
\[
|\bar R_u-(Q_{P,h,J}\vartheta)_u|
\le C_{\mathrm{loc}}^2\delta^S+2C_{\mathrm{loc}}h^\gamma
\le
(C_{\mathrm{loc}}^2+2C_{\mathrm{loc}})
(h^\gamma+\delta^S).
\]
Define
\[
C_{\mathrm{bias}}
=\sqrt q\,(C_{\mathrm{loc}}^2+2C_{\mathrm{loc}}).
\]
Squaring the coordinate bounds, summing over the \(q\) coordinates, and taking square roots proves
\[
\|\bar R-Q_{P,h,J}\vartheta\|_2
\le C_{\mathrm{bias}}(h^\gamma+\delta^S).
\]%
% lean: CausalSmith.Stat.TwosamplePointcateAnnotationFrontier.population_bias_coordinate_bound

\item Define the squared deviations, the bias scale, and the good event by
\[
\begin{aligned}
Z_R(\mathcal D_{n,m},U)
 &=\|\widehat R_{h,J}-\bar R\|_2^2,\\
Z_Q(\mathcal D_{n,m},U)
 &=\|\widehat Q_{h,J}-Q_{P,h,J}\|_F^2,\\
b_{\mathrm{bias}}&=h^\gamma+\delta^S,\\
g_\varepsilon&=\frac{\varepsilon(1-\varepsilon)}2,\\
G_{\mathrm{good}}
 &=\{(\mathcal D_{n,m},U):
       \lambda_{\min}(\widehat Q_{h,J})\ge g_\varepsilon\}.
\end{aligned}
\]
The first two functions ignore \(U\).

On \(G_{\mathrm{good}}\), the Rayleigh bound is
\[
\xi_{\mathrm{inv}}^\top\widehat Q_{h,J}\xi_{\mathrm{inv}}
 \ge g_\varepsilon\|\xi_{\mathrm{inv}}\|_2^2
\qquad(\xi_{\mathrm{inv}}\in\mathbb R^q).
\]
It excludes a nonzero kernel and thus makes the matrix invertible.
For arbitrary \(\xi_{\mathrm{rhs}}\in\mathbb R^q\), set
\[
\xi_{\mathrm{inv}}=\widehat Q_{h,J}^{-1}\xi_{\mathrm{rhs}}.
\]
Cauchy--Schwarz then gives
\[
g_\varepsilon\|\xi_{\mathrm{inv}}\|_2^2
\le \xi_{\mathrm{inv}}^\top\xi_{\mathrm{rhs}}
\le \|\xi_{\mathrm{inv}}\|_2\|\xi_{\mathrm{rhs}}\|_2,
\]
and consequently
\[
\|\widehat Q_{h,J}^{-1}\xi_{\mathrm{rhs}}\|_2
\le g_\varepsilon^{-1}\|\xi_{\mathrm{rhs}}\|_2.
\]
For \(\xi_{\mathrm{inv}}=0\) this is immediate; otherwise division by its positive norm proves it.

The inverse identity and the residual decomposition are
\[
\widehat Q_{h,J}^{-1}\widehat R_{h,J}-\vartheta
=
\widehat Q_{h,J}^{-1}
 (\widehat R_{h,J}-\widehat Q_{h,J}\vartheta),
\]
\[
\widehat R_{h,J}-\widehat Q_{h,J}\vartheta
=
(\widehat R_{h,J}-\bar R)
+(\bar R-Q_{P,h,J}\vartheta)
+(Q_{P,h,J}-\widehat Q_{h,J})\vartheta.
\]
The triangle inequality and rowwise Cauchy--Schwarz imply
\[
\|\widehat R_{h,J}-\widehat Q_{h,J}\vartheta\|_2
\le
\sqrt{Z_R}
+C_{\mathrm{bias}}b_{\mathrm{bias}}
+C_{\mathrm{loc}}\sqrt{Z_Q}.
\]
Because both arm means lie in \([0,1]\),
\[
\tau_P(x_0)\in[-1,1].
\]
Clipping toward this interval decreases absolute error: below its lower endpoint the clipped endpoint lies between the value and the target, above its upper endpoint the same holds, and inside the interval clipping preserves the value. Thus, on \(G_{\mathrm{good}}\),
\[
\begin{aligned}
|\widehat T_{h,J}-\tau_P(x_0)|
&\le g_\varepsilon^{-1}\|r_0\|_2
 \|\widehat R_{h,J}-\widehat Q_{h,J}\vartheta\|_2\\
&\le g_\varepsilon^{-1}C_{\mathrm{loc}}
 \left(\sqrt{Z_R}
 +C_{\mathrm{bias}}b_{\mathrm{bias}}
 +C_{\mathrm{loc}}\sqrt{Z_Q}\right).
\end{aligned}
\]
Define the positive finite constant
\[
C_{\mathrm{pt}}
=g_\varepsilon^{-1}C_{\mathrm{loc}}
 (1+C_{\mathrm{bias}}+C_{\mathrm{loc}}).
\]
All three coefficients in the preceding parentheses are at most
\(1+C_{\mathrm{bias}}+C_{\mathrm{loc}}\). On the complementary event the decision is zero and its error is at most \(1\). Combining these branches gives, on every sample,
\[
|\widehat T_{h,J}-\tau_P(x_0)|
\le
C_{\mathrm{pt}}
 (b_{\mathrm{bias}}+\sqrt{Z_R}+\sqrt{Z_Q})
+\mathbf1_{G_{\mathrm{good}}^c}.
\]%
% lean: CausalSmith.Stat.TwosamplePointcateAnnotationFrontier.rectangular_pointwise_error_bound

\item By \cref{obj:lem:population-positivity},
\[
v_{\mathrm{unit}}^\top Q_{P,h,J}v_{\mathrm{unit}}
\ge\varepsilon(1-\varepsilon)=2g_\varepsilon
\qquad
(v_{\mathrm{unit}}\in\mathbb R^q,\ \|v_{\mathrm{unit}}\|_2=1).
\]
For each such vector, finite Cauchy--Schwarz gives
\[
\left|
v_{\mathrm{unit}}^\top
(\widehat Q_{h,J}-Q_{P,h,J})v_{\mathrm{unit}}
\right|^2
\le Z_Q.
\]
If \(Z_Q<g_\varepsilon^2\), every unit-vector quadratic form of
\(\widehat Q_{h,J}\) is at least \(g_\varepsilon\), by subtracting this deviation from \(2g_\varepsilon\). Its infimum is therefore at least \(g_\varepsilon\). Consequently,
\[
G_{\mathrm{good}}^c
\ \Longrightarrow\
Z_Q\ge g_\varepsilon^2,
\qquad
\mathbf1_{G_{\mathrm{good}}^c}
\le g_\varepsilon^{-1}\sqrt{Z_Q}.
\]
Define
\[
\kappa_{\mathrm{mean}}=C_{\mathrm{pt}}+g_\varepsilon^{-1}.
\]
The preceding pointwise bound now becomes
\[
|\widehat T_{h,J}-\tau_P(x_0)|
\le
\kappa_{\mathrm{mean}}
 (b_{\mathrm{bias}}+\sqrt{Z_R}+\sqrt{Z_Q}).
\]

Define the variance scale by
\[
V_{\mathrm{var}}
=\frac1{nh^d}+\frac{qJ^d}{nNh^{2d}}.
\]
By \cref{obj:lem:rectangular-sampling-bound}, there is a uniform positive finite constant \(C_{\mathrm{samp}}\) such that
\[
\mathbb E_{\mathsf E_{n,m}(P)}Z_R
+\mathbb E_{\mathsf E_{n,m}(P)}Z_Q
\le C_{\mathrm{samp}}V_{\mathrm{var}}.
\]
Both deviations are nonnegative. Cauchy--Schwarz under the probability sampling law therefore gives
\[
\begin{aligned}
\mathbb E_{\mathsf E_{n,m}(P)}\sqrt{Z_R}
&\le
\sqrt{\mathbb E_{\mathsf E_{n,m}(P)}Z_R}
\le\sqrt{C_{\mathrm{samp}}}\sqrt{V_{\mathrm{var}}},\\
\mathbb E_{\mathsf E_{n,m}(P)}\sqrt{Z_Q}
&\le
\sqrt{\mathbb E_{\mathsf E_{n,m}(P)}Z_Q}
\le\sqrt{C_{\mathrm{samp}}}\sqrt{V_{\mathrm{var}}}.
\end{aligned}
\]
The decision and target both belong to \([-1,1]\), so their absolute difference is at most \(2\) and is integrable. Integrating the pointwise inequality and defining
\[
C_{\mathrm{mean}}
=\kappa_{\mathrm{mean}}(1+2\sqrt{C_{\mathrm{samp}}})
\]
proves
\[
\begin{aligned}
\mathcal R_{n,m}(\widehat T_{h,J},P)
&\le\kappa_{\mathrm{mean}}
 \left(b_{\mathrm{bias}}
       +2\sqrt{C_{\mathrm{samp}}}\sqrt{V_{\mathrm{var}}}\right)\\
&\le C_{\mathrm{mean}}
 (b_{\mathrm{bias}}+\sqrt{V_{\mathrm{var}}}).
\end{aligned}
\]%
% lean: CausalSmith.Stat.TwosamplePointcateAnnotationFrontier.rectangular_risk_sqrt_variance

\item Define
\[
\rho_{\mathrm{lab}}=(nh^d)^{-1/2},\qquad
\rho_{\mathrm{rect}}=(nNh^d\delta^d)^{-1/2},\qquad
a_{\mathrm{noise}}=\rho_{\mathrm{lab}}+\rho_{\mathrm{rect}},
\qquad
\kappa_{\mathrm{rank}}=1+\sqrt q.
\]
All these quantities are nonnegative, and
\(\kappa_{\mathrm{rank}}\ge1\). Since \(\delta=h/J\),
\[
\rho_{\mathrm{lab}}^2=\frac1{nh^d},\qquad
\rho_{\mathrm{rect}}^2
=\frac{J^d}{nNh^{2d}},\qquad
V_{\mathrm{var}}
=\rho_{\mathrm{lab}}^2+q\rho_{\mathrm{rect}}^2.
\]
It follows that
\[
V_{\mathrm{var}}
\le
(\rho_{\mathrm{lab}}+\sqrt q\,\rho_{\mathrm{rect}})^2,
\qquad
\rho_{\mathrm{lab}}+\sqrt q\,\rho_{\mathrm{rect}}
\le\kappa_{\mathrm{rank}}a_{\mathrm{noise}},
\]
and hence
\[
\sqrt{V_{\mathrm{var}}}
\le\kappa_{\mathrm{rank}}a_{\mathrm{noise}}.
\]
Define
\[
C_{\mathrm{fixed}}
=C_{\mathrm{mean}}\kappa_{\mathrm{rank}}+2.
\]
Using also
\(b_{\mathrm{bias}}\le\kappa_{\mathrm{rank}}b_{\mathrm{bias}}\), the preceding risk bound gives
\[
\mathcal R_{n,m}(\widehat T_{h,J},P)
\le C_{\mathrm{fixed}}
 (b_{\mathrm{bias}}+a_{\mathrm{noise}}).
\]
If \(b_{\mathrm{bias}}+a_{\mathrm{noise}}\le1\), this is the desired bound with the minimum. If
\(b_{\mathrm{bias}}+a_{\mathrm{noise}}>1\), the pointwise diameter bound gives
\[
\mathcal R_{n,m}(\widehat T_{h,J},P)
\le2\le C_{\mathrm{fixed}}.
\]
Thus, uniformly over \(P\in\mathcal M\),
\[
\mathcal R_{n,m}(\widehat T_{h,J},P)
\le C_{\mathrm{fixed}}
\min\left\{1,\,
h^\gamma+\delta^S+(nh^d)^{-1/2}
 +(nNh^d\delta^d)^{-1/2}\right\}.
\]%
% lean: CausalSmith.Stat.TwosamplePointcateAnnotationFrontier.rectangular_fixed_tuning_risk

\item We next establish the public tuning bound. Define
\[
\Delta_{\mathrm{o}}=2\gamma+d,\qquad
\Delta=2\gamma+d+\frac{\gamma d}{S},
\qquad
c_{\mathrm{lab}}=2^{\Delta_{\mathrm{o}}/2},
\qquad
c_{\mathrm{rect}}
=2^{d/2}2^{\Delta/2}+2^{\Delta_{\mathrm{o}}},
\qquad
C_{\mathrm{bal}}=2+c_{\mathrm{lab}}+c_{\mathrm{rect}}.
\]
The parameter domain gives \(S>0\), \(\gamma>0\),
\(\Delta_{\mathrm{o}}>0\), and \(\Delta>0\), so these constants are positive and finite.

For fixed \(n\ge2\) and \(m\ge0\), define the two bandwidth orders
\[
u_{\mathrm{o}}=n^{-1/\Delta_{\mathrm{o}}},
\qquad
u_{\mathrm{rect}}=(nN)^{-1/\Delta}.
\]
Because \(n\ge1\) and \(nN\ge1\),
\[
0<u_{\mathrm{o}}\le1,\qquad
0<u_{\mathrm{rect}}\le1.
\]
The definition in \cref{obj:def:tuning} consequently gives
\[
h_*=\frac12\max\{u_{\mathrm{o}},u_{\mathrm{rect}}\},
\qquad
0<h_*\le\frac12,\qquad
u_{\mathrm{o}}\le2h_*,
\qquad
u_{\mathrm{rect}}\le2h_*.
\]
Moreover,
\[
\begin{aligned}
h_*^\gamma
&\le \max\{u_{\mathrm{o}}^\gamma,u_{\mathrm{rect}}^\gamma\}\\
&=\max\left\{
 n^{-\gamma/\Delta_{\mathrm{o}}},
 (nN)^{-\gamma/\Delta}\right\}
=r_{\mathrm{up}}(n,m),
\end{aligned}
\]
using \cref{obj:def:sharp-rate}.

Taking logarithms of \(u_{\mathrm{o}}\le2h_*\) gives
\[
-\log n
\le\Delta_{\mathrm{o}}(\log2+\log h_*).
\]
Therefore
\[
\begin{aligned}
\log\bigl((nh_*^d)^{-1/2}\bigr)
&=-\frac12(\log n+d\log h_*)\\
&\le\frac{\Delta_{\mathrm{o}}}{2}\log2
   +\gamma\log h_*,
\end{aligned}
\]
where \(\Delta_{\mathrm{o}}-d=2\gamma\). Exponentiating yields
\[
(nh_*^d)^{-1/2}\le c_{\mathrm{lab}}h_*^\gamma.
\]%
% lean: CausalSmith.Stat.TwosamplePointcateAnnotationFrontier.bandwidth_monomial_balance

\item Suppose first that \(S<\gamma\). Define
\[
z_{\mathrm{grid}}=h_*^{-(\gamma/S-1)}.
\]
Then \(z_{\mathrm{grid}}\ge1\), and \cref{obj:def:tuning} gives
\[
J_*=\lceil z_{\mathrm{grid}}\rceil,\qquad
1\le J_*,\qquad
z_{\mathrm{grid}}\le J_*<z_{\mathrm{grid}}+1
\le2z_{\mathrm{grid}}.
\]
Since \(h_*/z_{\mathrm{grid}}=h_*^{\gamma/S}\), these inequalities imply
\[
0<\frac{h_*}{J_*},\qquad
\frac12h_*^{\gamma/S}
\le\frac{h_*}{J_*}\le h_*^{\gamma/S},
\qquad
\left(\frac{h_*}{J_*}\right)^S\le h_*^\gamma.
\]
The lower cell-side bound gives
\[
\log(h_*/J_*)\ge\frac{\gamma}{S}\log h_*-\log2.
\]
Substituting this inequality into the logarithm of the rectangular sampling term and exponentiating proves
\[
\left(nNh_*^d(h_*/J_*)^d\right)^{-1/2}
\le
2^{d/2}
\left(nNh_*^{d+\gamma d/S}\right)^{-1/2}.
\]
Similarly, \(u_{\mathrm{rect}}\le2h_*\) gives
\[
-\log(nN)\le\Delta(\log2+\log h_*).
\]
As \(\Delta-(d+\gamma d/S)=2\gamma\),
\[
\begin{aligned}
\log\left(\left(nNh_*^{d+\gamma d/S}\right)^{-1/2}\right)
&=-\frac12\left(\log(nN)
 +(d+\gamma d/S)\log h_*\right)\\
&\le\frac{\Delta}{2}\log2+\gamma\log h_*.
\end{aligned}
\]
Consequently,
\[
\left(nNh_*^d(h_*/J_*)^d\right)^{-1/2}
\le2^{d/2}2^{\Delta/2}h_*^\gamma
\le c_{\mathrm{rect}}h_*^\gamma.
\]%
% lean: CausalSmith.Stat.TwosamplePointcateAnnotationFrontier.ceil_grid_side_bounds

\item Suppose instead that \(S\ge\gamma\), including equality. The exponent in
\cref{obj:def:tuning} is then zero, so
\[
J_*=1,\qquad
0<\frac{h_*}{J_*}=h_*,
\qquad
\left(\frac{h_*}{J_*}\right)^S
=h_*^S\le h_*^\gamma,
\]
because \(0<h_*\le1\).

The oracle bandwidth inequality also reads
\[
(n^2)^{-1/(2\Delta_{\mathrm{o}})}
=u_{\mathrm{o}}\le2h_*.
\]
Taking logarithms gives
\[
-\log(n^2)
\le2\Delta_{\mathrm{o}}(\log2+\log h_*).
\]
Using \(2\Delta_{\mathrm{o}}-2d=4\gamma\), we obtain
\[
\begin{aligned}
\log\left((n^2h_*^{2d})^{-1/2}\right)
&=-\frac12\left(\log(n^2)+2d\log h_*\right)\\
&\le\Delta_{\mathrm{o}}\log2+2\gamma\log h_*.
\end{aligned}
\]
Since \(N\ge n\) and \(h_*^{2\gamma}\le h_*^\gamma\),
\[
\begin{aligned}
\left(nNh_*^d(h_*/J_*)^d\right)^{-1/2}
&=(nNh_*^{2d})^{-1/2}\\
&\le(n^2h_*^{2d})^{-1/2}\\
&\le2^{\Delta_{\mathrm{o}}}h_*^{2\gamma}\\
&\le c_{\mathrm{rect}}h_*^\gamma.
\end{aligned}
\]
Thus both cases give admissible tuning and the four-term balance
\[
\begin{gathered}
0<h_*\le\frac12,\qquad 1\le J_*,\\
h_*^\gamma+(h_*/J_*)^S+(nh_*^d)^{-1/2}
+\left(nNh_*^d(h_*/J_*)^d\right)^{-1/2}\\
\le(2+c_{\mathrm{lab}}+c_{\mathrm{rect}})h_*^\gamma
\le C_{\mathrm{bal}}r_{\mathrm{up}}(n,m).
\end{gathered}
\]%
% lean: CausalSmith.Stat.TwosamplePointcateAnnotationFrontier.tuning_balance

\item Finally, define one constant for both assertions:
\[
C=C_{\mathrm{fixed}}(C_{\mathrm{bal}}+1).
\]
It is positive and finite, and
\[
C_{\mathrm{fixed}}\le C,\qquad
C_{\mathrm{fixed}}C_{\mathrm{bal}}\le C.
\]
The fixed-tuning envelope has a nonnegative minimum, so enlarging its constant from \(C_{\mathrm{fixed}}\) to \(C\) and taking the supremum over \(P\in\mathcal M\) proves the first risk assertion.

For the second assertion, the tuning bounds make
\((h_*,J_*)\) admissible for the fixed-tuning result. The identity in
\cref{obj:def:tuning} supplies
\[
\widehat T_{\mathrm{up}}=\widehat T_{h_*,J_*},
\]
and hence also its finiteness and Borel measurability. Using the minimum's upper bound by its four-term argument gives, for every \(P\in\mathcal M\),
\[
\begin{aligned}
\mathcal R_{n,m}(\widehat T_{\mathrm{up}},P)
&\le C_{\mathrm{fixed}}
 \left\{h_*^\gamma+(h_*/J_*)^S+(nh_*^d)^{-1/2}
 +\left(nNh_*^d(h_*/J_*)^d\right)^{-1/2}\right\}\\
&\le C_{\mathrm{fixed}}C_{\mathrm{bal}}r_{\mathrm{up}}(n,m)\\
&\le Cr_{\mathrm{up}}(n,m).
\end{aligned}
\]
The rate is nonnegative as the maximum of two positive powers of positive sample counts. All constants are uniform over the primitive class, so taking the supremum proves the tuned assertion.%
% lean: CausalSmith.Stat.TwosamplePointcateAnnotationFrontier.rectangular_upper
\end{enumerate}
\end{proof}

\begin{proof}[Proof of \cref{obj:thm:sharp-annotation-frontier}]
Fix the public parameters as in the statement. Throughout the proof, use
\[
N=n+m,\qquad S=\alpha+\beta,\qquad
\Delta=2\gamma+d+\frac{\gamma d}{S},\qquad
S_{\mathrm{crit}}=\frac{\gamma d}{2\gamma+d},\qquad
q_*=\frac{\gamma d}{S(2\gamma+d)},
\]
and
\[
r_o(n)=n^{-\gamma/(2\gamma+d)},\qquad
r_*=\max\{r_o(n),(nN)^{-\gamma/\Delta}\},
\]
where the last equality is \cref{obj:def:sharp-rate}. All sample sizes below satisfy \(n,m\in\mathbb N\) and \(n\ge2\).

\begin{enumerate}
\item \emph{Supplied-propensity lower bound.}
By \cref{obj:thm:supplied-propensity}, there is a constant \(c_O>0\), depending on the fixed public parameters, such that
\[
c_O r_o(n)\le R^e(n,m)
\qquad\text{for every }n,m.
\]
Every admissible decision \(T\) induces a supplied-propensity decision by
\[
T^e(\mathcal D_{n,m},U,e)=T(\mathcal D_{n,m},U)
\]
for every function \(e:[0,1]^d\to\mathbb R\). For each supplied function this decision is Borel, and its loss under each \(P\in\mathcal M\) is
\[
\int
\left|T^e(\mathcal D_{n,m},U,e_P)-\tau_P(x_0)\right|
\,d\mathsf E_{n,m}(P)
=
\mathcal R_{n,m}(T,P).
\]
Taking the supremum over the same primitive class and then the infimum over decisions, using \cref{obj:def:oracle-risk,obj:def:minimax}, gives
\[
R^e(n,m)\le R(n,m).
\]
Consequently,
\[
c_O r_o(n)\le R(n,m)
\qquad\text{for every }n,m.
\]
% lean: CausalSmith.Stat.TwosamplePointcateAnnotationFrontier.oracleRisk_le_minimaxRisk

\item \emph{A uniform lower constant.}
First suppose \(S<S_{\mathrm{crit}}\). By
\cref{obj:lem:nuisance-frontier-converse}, there is a constant \(c_B>0\), depending on the fixed public parameters, such that
\[
c_B(nN)^{-\gamma/\Delta}\le R(n,m)
\qquad\text{whenever }n\le N\le n^{q_*}.
\]
Choose
\[
c=\min\{c_O,c_B\}>0.
\]

We establish the rate comparisons used in this case and in the remaining smoothness case. The parameter conditions imply
\[
S>0,\qquad \gamma>0,\qquad 2\gamma+d>0,\qquad
\Delta>0,\qquad q_*>0,\qquad N\ge n\ge2.
\]
Direct substitution gives
\[
q_*=
\frac{\gamma d}{S(2\gamma+d)}
=
\frac{S_{\mathrm{crit}}}{S},
\]
and
\[
\Delta
=
(2\gamma+d)+\frac{\gamma d}{S}
=
(2\gamma+d)
\left(1+\frac{\gamma d}{S(2\gamma+d)}\right)
=
(2\gamma+d)(1+q_*).
\]
% lean: CausalSmith.Stat.TwosamplePointcateAnnotationFrontier.rate_exponent_identities

Since \(n>0\), the laws of real powers therefore give the boundary identity
\[
\begin{aligned}
(n\,n^{q_*})^{-\gamma/\Delta}
&=(n^{1+q_*})^{-\gamma/\Delta}\\
&=n^{-(1+q_*)\gamma/\Delta}\\
&=n^{-\gamma/(2\gamma+d)}
=r_o(n).
\end{aligned}
\]
The exponent \(-\gamma/\Delta\) is negative. Thus, when \(N\ge n^{q_*}\),
\[
0<n\,n^{q_*}\le nN
\quad\Longrightarrow\quad
(nN)^{-\gamma/\Delta}
\le(n\,n^{q_*})^{-\gamma/\Delta}
=r_o(n).
\]
When \(N\le n^{q_*}\), instead,
\[
0<nN\le n\,n^{q_*}
\quad\Longrightarrow\quad
r_o(n)
=(n\,n^{q_*})^{-\gamma/\Delta}
\le(nN)^{-\gamma/\Delta}.
\]
For completeness, if \(S\ge S_{\mathrm{crit}}\), then
\[
q_*=\frac{S_{\mathrm{crit}}}{S}\le1,
\qquad
n^{q_*}\le n^1=n\le N,
\]
where the power comparison uses \(n\ge1\). Taking the maximum in the definition of \(r_*\) consequently gives
\[
\begin{aligned}
S\ge S_{\mathrm{crit}}
&\quad\Longrightarrow\quad r_*=r_o(n),\\
S<S_{\mathrm{crit}},\ N\le n^{q_*}
&\quad\Longrightarrow\quad r_*=(nN)^{-\gamma/\Delta},\\
S<S_{\mathrm{crit}},\ N\ge n^{q_*}
&\quad\Longrightarrow\quad r_*=r_o(n).
\end{aligned}
\]
The boundary identity also gives
\[
N=n^{q_*}
\quad\Longrightarrow\quad
(nN)^{-\gamma/\Delta}=r_o(n).
\]
% lean: CausalSmith.Stat.TwosamplePointcateAnnotationFrontier.rate_branches

Return to the case \(S<S_{\mathrm{crit}}\). If \(N\le n^{q_*}\), then \(N\ge n\), so the nuisance lower bound applies and yields
\[
c r_*
=c(nN)^{-\gamma/\Delta}
\le c_B(nN)^{-\gamma/\Delta}
\le R(n,m).
\]
If \(N>n^{q_*}\), the supplied-propensity lower bound yields
\[
c r_*=c r_o(n)\le c_O r_o(n)\le R(n,m).
\]
Here the comparisons of constants preserve the inequalities because both rate terms are nonnegative.

In the remaining case \(S\ge S_{\mathrm{crit}}\), choose
\[
c=c_O>0.
\]
Then \(r_*=r_o(n)\), and hence
\[
c r_*=c_O r_o(n)\le R(n,m).
\]
This constructs a positive constant \(c\), depending on the fixed public parameters and valid simultaneously for every admissible sample-size pair, such that
\[
c r_*\le R(n,m).
\]
% lean: CausalSmith.Stat.TwosamplePointcateAnnotationFrontier.sharpRate_minimax_lower

\item \emph{Attainment and the upper constant.}
By \cref{obj:thm:rectangular-upper}, there is a real constant \(C_U>0\), depending on the fixed public parameters, such that the publicly tuned decision is Borel measurable and its risk under every \(P\in\mathcal M\) satisfies
\[
\mathcal R_{n,m}(\widehat T_{\mathrm{up}},P)
\le C_U r_{\mathrm{up}}(n,m).
\]
By \cref{obj:def:sharp-decision,obj:def:sharp-rate},
\[
T_*=\widehat T_{\mathrm{up}},
\qquad
r_{\mathrm{up}}(n,m)=r_*.
\]
Set
\[
C=\max\{c,C_U\}.
\]
Then \(0<c\le C<\infty\). For every admissible sample-size pair, \(T_*\) is Borel measurable, and taking the supremum of the pointwise risk bound gives
\[
\sup_{P\in\mathcal M}\mathcal R_{n,m}(T_*,P)
\le C_U r_*
\le C r_*.
\]
The last inequality follows from \(C_U\le C\) and \(r_*\ge0\).
% lean: CausalSmith.Stat.TwosamplePointcateAnnotationFrontier.sharp_annotation_frontier

\item \emph{Minimax comparison and arbitrary decisions.}
Since \(T_*\) belongs to the decision class in \cref{obj:def:minimax},
\[
R(n,m)\le
\sup_{P\in\mathcal M}\mathcal R_{n,m}(T_*,P).
\]
Together with the bounds already established, this proves
\[
c r_*
\le R(n,m)
\le\sup_{P\in\mathcal M}\mathcal R_{n,m}(T_*,P)
\le C r_*.
\]
For any Borel real-valued decision \(T\) using the original dataset and the independent uniform randomizer, the same infimum comparison gives
\[
c r_*\le R(n,m)
\le\sup_{P\in\mathcal M}\mathcal R_{n,m}(T,P).
\]
The rate comparisons and boundary identity established above supply all three stated branches and their equality at \(N=n^{q_*}\).
% lean: CausalSmith.Stat.TwosamplePointcateAnnotationFrontier.sharp_annotation_frontier
\end{enumerate}
\end{proof}

\begin{proof}[Proof of \cref{obj:thm:supplied-propensity}]
Fix the public parameters in the stated domain. Throughout, inverse-propensity quotients are assigned value zero when their denominator is zero:
\[
\frac{t_{\mathrm{num}}}{0}:=0
\qquad(t_{\mathrm{num}}\in\mathbb R).
\]
The designated propensity is extended by zero outside \([0,1]^d\).

\begin{enumerate}
\item Define the constant-nuisance subfamily and its supplied-propensity risk by
\[
\mathcal M_{1/2}
=
\left\{
P\in\mathcal M:
e_P(x)=\mu_{0,P}(x)=\frac12
\text{ for every }x\in[0,1]^d
\right\},
\]
\[
\mathcal R^e_{n,m}(T^e,P)
=
\int
\left|T^e(\mathcal D_{n,m},U,e_P)-\tau_P(x_0)\right|
\,d\mathsf E_{n,m}(P),
\qquad
R^e_{1/2}(n,m)
=
\inf_{T^e}\sup_{P\in\mathcal M_{1/2}}
\mathcal R^e_{n,m}(T^e,P).
\]
These loss integrals may take the value \(+\infty\). Since the subfamily uses the same decision class as \cref{obj:def:oracle-risk}, restriction of the supremum gives
\[
R^e_{1/2}(n,m)\le R^e(n,m).
\]
It therefore suffices to establish a lower bound for \(R^e_{1/2}(n,m)\), a uniform upper bound for the displayed smoother, and its admissibility. We first construct the lower-bound family.
% lean: CausalSmith.Stat.TwosamplePointcateAnnotationFrontier.oracleSubfamilyRisk_le_oracleRisk

\item Use the smooth profile from \cref{obj:def:marked-handle}, with the explicit formulas
\[
\chi(\upsilon)=
\begin{cases}
\exp(-1/\upsilon),&\upsilon>0,\\
0,&\upsilon\le0,
\end{cases}
\qquad \upsilon\in\mathbb R,
\qquad
B(w)=\prod_{\iota_{\mathrm{cov}}=1}^d\frac{\chi(1-4w_{\iota_{\mathrm{cov}}}^2)}{\chi(1)},
\qquad w\in\mathbb R^d.
\]
For this argument, extend the scaled-bump formula to every positive localization width and every covariate vector:
\[
B_h(x)=B\bigl((x-x_0)/h\bigr),
\qquad
C_h=\prod_{j=1}^d
\left[\frac12-\frac h2,\frac12+\frac h2\right],
\qquad h>0,\quad x\in\mathbb R^d.
\]
Smoothness of the exponential glue function and finite products give
\[
B\in C^\infty(\mathbb R^d),
\qquad
\operatorname{supp}(B)\subseteq[-1/2,1/2]^d,
\qquad
0\le B\le1,
\qquad
B(0)=1.
\]
In particular,
\[
0\le B_h\le1,\qquad
B_h(x)=0\quad(x\notin C_h),\qquad
B_h(x_0)=1.
\]

We need a uniform Hölder bound after scaling. Define
\[
p_*=\lceil\gamma\rceil-1,
\qquad
s_o=\gamma-p_*,
\qquad
0\le p_*\le2,\quad 0<s_o\le1.
\]
Write the iterated derivatives as
\[
D^0B=B,\qquad D^{\iota_{\mathrm{der}}+1}B=D(D^{\iota_{\mathrm{der}}}B),
\qquad \iota_{\mathrm{der}}\in\mathbb N,
\]
and define their finite operator-norm bounds by
\[
C_{B,\iota_{\mathrm{der}}}=\sup_{w\in\mathbb R^d}\|D^{\iota_{\mathrm{der}}}B(w)\|,
\qquad \iota_{\mathrm{der}}\in\mathbb N.
\]
Their finiteness follows from continuity and compact support. Set
\[
A_B=\max\{C_{B,p_*},C_{B,p_*+1}\},
\qquad
M_B=2A_B.
\]
The derivative bound and the mean-value inequality give, for every \(w,w'\in\mathbb R^d\),
\[
\|D^{p_*}B(w)-D^{p_*}B(w')\|
\le A_B\|w-w'\|_2,
\qquad
\|D^{p_*}B(w)-D^{p_*}B(w')\|\le2A_B.
\]
If \(\|w-w'\|_2\le1\), use
\(\|w-w'\|_2\le\|w-w'\|_2^{s_o}\); if
\(\|w-w'\|_2>1\), use \(1\le\|w-w'\|_2^{s_o}\).
Thus both cases yield
\[
\|D^{p_*}B(w)-D^{p_*}B(w')\|
\le M_B\|w-w'\|_2^{s_o}.
\]

Translation and dilation satisfy
\[
D^{\iota_{\mathrm{der}}}B_h(x)=h^{-\iota_{\mathrm{der}}}D^{\iota_{\mathrm{der}}}B\bigl((x-x_0)/h\bigr).
\]
Evaluating these multilinear maps in coordinate directions gives the bounds
\[
|\partial^\kappa B_h(x)|
\le d^{\iota_{\mathrm{der}}} C_{B,\iota_{\mathrm{der}}}h^{-\iota_{\mathrm{der}}}
\qquad(|\kappa|=\iota_{\mathrm{der}}\le p_*),
\]
and
\[
|\partial^\kappa B_h(x)-\partial^\kappa B_h(x')|
\le d^{p_*}M_Bh^{-\gamma}\|x-x'\|_2^{s_o}
\qquad(|\kappa|=p_*).
\]
Here the dimension factors are valid because every coordinate direction has norm at most \(d\). Define
\[
S_B=\sum_{\iota_{\mathrm{der}}=0}^{p_*}d^{\iota_{\mathrm{der}}}C_{B,\iota_{\mathrm{der}}},
\qquad
D_B=S_B+d^{p_*}M_B+1.
\]
For \(0<h\le1\) and \(\iota_{\mathrm{der}}\le p_*<\gamma\), one has
\(h^{-\iota_{\mathrm{der}}}\le h^{-\gamma}\). Applying the norm definition in
\cref{obj:def:holder-norm} to the preceding derivative and modulus bounds proves
\[
D_B>0,
\qquad
\|B_h\|_{H^\gamma}\le D_Bh^{-\gamma}
\qquad(0<h\le1).
\]
This also covers the integer smoothness endpoints, where \(s_o=1\).
% lean: CausalSmith.Stat.TwosamplePointcateAnnotationFrontier.oracle_macro_holder

\item Choose the public amplitude constant
\[
\eta_o=\frac{1}{128(D_B+1)}.
\]
It satisfies
\[
\eta_o>0,\qquad
\eta_o(D_B+1)=\frac1{128},\qquad
\eta_o\le\frac18,\qquad
\eta_o\le\frac1{16},\qquad
\eta_oD_B\le L.
\]
For each \(n\ge2\), define
\[
h_n=n^{-1/(2\gamma+d)},
\qquad
b_n=\eta_o r_o(n).
\]
Since \(2\gamma+d>0\) and \(n\ge2\), the power identities give
\[
0<h_n\le1,
\qquad
h_n^\gamma=r_o(n),
\qquad
n\,r_o(n)^2h_n^d=1,
\qquad
0<b_n\le\frac18.
\]
Consequently,
\[
\|b_nB_{h_n}\|_{H^\gamma}
\le b_nD_Bh_n^{-\gamma}
=\eta_oD_B
\le L.
\]

For \(\theta\in\{-1,+1\}\), define the primitive law \(P_{o,\theta}\) by
\[
X\sim\operatorname{Unif}([0,1]^d),
\]
\[
\mathcal L_{P_{o,\theta}}\bigl(A,Y(0),Y(1)\mid X=x\bigr)
=
\operatorname{Bernoulli}(1/2)
\otimes
\operatorname{Bernoulli}(1/2)
\otimes
\operatorname{Bernoulli}\bigl(1/2+\theta b_nB_{h_n}(x)\bigr),
\qquad x\in[0,1]^d.
\]
These conditional margins are Borel. Conditional independence gives exchangeability, and their bounds give
\[
e(x)=\mu_0(x)=\frac12,
\qquad
\frac18\le\frac12+\theta b_nB_{h_n}(x)\le\frac78.
\]
The constant propensity and control functions have Hölder norms at most \(1/2\), while multiplication by either sign preserves the preceding contrast-norm bound. Thus the construction satisfies every condition of \cref{obj:def:primitive-class}.

The designated versions agree with these admissible margins on the cube. Indeed, conditional-margin uniqueness gives almost-sure equality, and continuity together with the full support of the uniform design promotes that equality to pointwise equality. Hence
\[
P_{o,\theta}\in\mathcal M_{1/2},
\qquad
e_{P_{o,\theta}}(x)=\mu_{0,P_{o,\theta}}(x)=\frac12,
\qquad
\tau_{P_{o,\theta}}(x)=\theta b_nB_{h_n}(x)
\quad(x\in[0,1]^d),
\]
so that
\[
\tau_{P_{o,\theta}}(x_0)=\theta b_n.
\]
% lean: CausalSmith.Stat.TwosamplePointcateAnnotationFrontier.oracle_bump_testing_family

\item Define the observed-record and full-experiment laws by
\[
P_{o,\theta}^{\mathrm{obs}}
=\mathcal L_{P_{o,\theta}}(X,A,Y),
\qquad
M_{o,\theta}=\mathsf E_{n,m}(P_{o,\theta}).
\]
Use the common probability reference measure
\[
\nu_o
=
\operatorname{Unif}([0,1]^d)
\otimes\operatorname{Bernoulli}(1/2)
\otimes\operatorname{Bernoulli}(1/2)
\]
on the covariate, treatment, and observed-response coordinates. Relative to this reference, the observed-record density is
\[
\ell_{o,\theta}(x,a_{\mathrm{obs}},y_{\mathrm{obs}})
=
\begin{cases}
1,&a_{\mathrm{obs}}=0,\\
1+2(2y_{\mathrm{obs}}-1)\theta b_nB_{h_n}(x),
&a_{\mathrm{obs}}=1,
\end{cases}
\]
for \(x\in[0,1]^d\) and \(a_{\mathrm{obs}},y_{\mathrm{obs}}\in\{0,1\}\). The amplitude bound implies
\[
\frac12\le\ell_{o,\theta}\le2.
\]
For any two density values at least \(1/2\), the squared sum of their square roots is at least one. The difference-of-squares identity therefore gives, pointwise,
\[
\bigl(\sqrt{\ell_{o,-1}}-\sqrt{\ell_{o,+1}}\bigr)^2
\le
(\ell_{o,-1}-\ell_{o,+1})^2
\le16b_n^2B_{h_n}(x)^2.
\]
The Hellinger normalization in \cref{obj:def:hellinger} can be evaluated under \(\nu_o\): the densities relative to the sum of the two laws are
\(\ell_{o,\theta}/(\ell_{o,-1}+\ell_{o,+1})\), and the common denominator cancels. Thus
\[
\mathsf h^2(P_{o,-1}^{\mathrm{obs}},P_{o,+1}^{\mathrm{obs}})
=
\int
\bigl(\sqrt{\ell_{o,-1}}-\sqrt{\ell_{o,+1}}\bigr)^2\,d\nu_o.
\]
Since \(B_{h_n}^2\le\mathbf1_{C_{h_n}}\), the containment
\(C_{h_n}\subseteq[0,1]^d\) and the product volume of this cube give
\[
\int_{[0,1]^d}B_{h_n}(x)^2\,dx
\le\operatorname{vol}(C_{h_n})
=h_n^d.
\]
Combining the preceding displays proves
\[
\mathsf h^2(P_{o,-1}^{\mathrm{obs}},P_{o,+1}^{\mathrm{obs}})
\le16b_n^2h_n^d.
\]

The auxiliary covariate-treatment marginal is identical under the two laws:
\[
P_{o,-1,XA}=P_{o,+1,XA}
=
\operatorname{Unif}([0,1]^d)\otimes\operatorname{Bernoulli}(1/2).
\]
The product bound and common-factor identity in
\cref{obj:lem:hellinger-block-calculus} therefore yield
\[
\mathsf h^2(M_{o,-1},M_{o,+1})
\le
n\,\mathsf h^2(P_{o,-1}^{\mathrm{obs}},P_{o,+1}^{\mathrm{obs}})
\le16n b_n^2h_n^d
=16\eta_o^2
\le\frac1{16}.
\]
This applies to every \(m\ge0\); at \(m=0\), the auxiliary product is the probability law on the empty record block. The independent uniform randomizer is another common probability factor. Finally, symmetry of the squared square-root difference gives the orientation used below:
\[
\mathsf h^2(M_{o,+1},M_{o,-1})
=
\mathsf h^2(M_{o,-1},M_{o,+1})
\le\frac1{16}.
\]
% lean: CausalSmith.Stat.TwosamplePointcateAnnotationFrontier.oracle_bump_information_at_rate

\item The target is a well-defined functional of these full-experiment laws. To see this, equality of two such laws implies equality of their first labeled-record marginals, since \(n>0\). Under either constructed law, the observed conditional probabilities identify the arm means through
\[
\mu_{1,P_{o,\theta}}(x)
=
2P_{o,\theta}(A=1,Y=1\mid X=x),
\qquad
\mu_{0,P_{o,\theta}}(x)
=
2P_{o,\theta}(A=0,Y=1\mid X=x)
\]
almost everywhere under the uniform design. Thus equality of the observed laws identifies all three conditional margins. Their conditional product construction then identifies the primitive laws, and uniqueness of the canonical continuous contrast gives
\[
M_{o,\theta}=M_{o,\theta'}
\quad\Longrightarrow\quad
\tau_{P_{o,\theta}}(x_0)=\tau_{P_{o,\theta'}}(x_0).
\]
In particular, \(M_{o,+1}\) and \(M_{o,-1}\) are distinct because \(b_n>0\). Define a functional on all measures on the full sample space by
\[
\Psi_o(\mathbb Q_o)=
\begin{cases}
b_n,&\mathbb Q_o=M_{o,+1},\\
-b_n,&\mathbb Q_o=M_{o,-1},\\
0,&\mathbb Q_o\notin\{M_{o,+1},M_{o,-1}\}.
\end{cases}
\]
It satisfies
\[
\Psi_o(M_{o,\theta})=\tau_{P_{o,\theta}}(x_0).
\]

The supplied propensity is the same function for both hypotheses, including its extension outside the cube:
\[
e_o(x)=
\begin{cases}
1/2,&x\in[0,1]^d,\\
0,&x\notin[0,1]^d.
\end{cases}
\]
For any admissible supplied-propensity decision, its section at \(e_o\) is measurable. Apply \cref{obj:lem:published-fuzzy-testing}, the absolute-loss bound of
\citep[Section 3, Lemma 1]{KennedyBalakrishnanRobinsWasserman2024},
with singleton priors, one observation consisting of the entire record block and randomizer, and the choices
\[
\omega_o=\operatorname{Dirac}(*),
\qquad
\mathbb B_*^{\,o}=M_{o,+1},
\qquad
\mathbb C_*^{\,o}=M_{o,-1},
\qquad
\mathcal P_o^{\mathrm{block}}=\{M_{o,+1},M_{o,-1}\},
\]
\[
v_{\mathrm{test}}=1,
\qquad
s_{\mathrm{test}}=2b_n,
\qquad
u_{\mathrm{test}}=\frac1{16}.
\]
The separation is \(2b_n>0\), and the preceding Hellinger bound supplies the distance hypothesis. Consequently,
\[
\sup_{\theta\in\{-1,+1\}}
\mathcal R^e_{n,m}(T^e,P_{o,\theta})
\ge
\frac{b_n}{2}
\left(1-\frac{\sqrt{63}}{32}\right).
\]
Since \(63\le64\),
\[
1-\frac{\sqrt{63}}{32}\ge\frac34,
\qquad
\frac{b_n}{2}
\left(1-\frac{\sqrt{63}}{32}\right)\ge\frac38b_n.
\]
The two primitive laws belong to \(\mathcal M_{1/2}\), so taking the infimum over decisions proves
\[
c=\frac38\eta_o>0,
\qquad
c\,r_o(n)\le R^e_{1/2}(n,m)
\qquad(n\ge2,\ m\ge0).
\]
% lean: CausalSmith.Stat.TwosamplePointcateAnnotationFrontier.oracle_subfamily_lower_rate

\item We next establish the upper bound. For every \(h>0\), define the localization measure and weight by
\[
d\nu_h(x)=h^{-d}\mathbf1_{C_h}(x)\,dx,
\qquad
w_h(x)=h^{-d}\mathbf1_{C_h}(x).
\]
The measure \(\nu_h\) is a probability measure because
\(\operatorname{vol}(C_h)=h^d\). For \(0<h\le1\), its support lies in the design cube.

To make the coarse basis explicit, define
\[
\mathcal I_2=
\{\kappa\in\mathbb N^d:|\kappa|\le2\},
\qquad
q=|\mathcal I_2|,
\]
and the one-dimensional polynomials
\[
\mathsf P_0(\upsilon)=1,\qquad
\mathsf P_1(\upsilon)=\sqrt{12}\,\upsilon,\qquad
\mathsf P_2(\upsilon)=\sqrt5\,(6\upsilon^2-1/2),
\qquad \upsilon\in\mathbb R.
\]
At the current localization width \(h\), the basis and evaluation vector are
\[
r_\kappa(x)=
\prod_{\iota_{\mathrm{cov}}=1}^d
\mathsf P_{\kappa_{\iota_{\mathrm{cov}}}}\bigl((x_{\iota_{\mathrm{cov}}}-1/2)/h\bigr),
\qquad
r(x)=(r_\kappa(x))_{\kappa\in\mathcal I_2},
\qquad
r_0=r(x_0),
\]
for every \(x\in\mathbb R^d\). Define the evaluation kernel and its dimension constant by
\[
\mathcal K_h(x)=r_0^\top r(x),
\qquad
K_{\mathrm{ev}}=q(4^d)^2.
\]
The constant coordinate belongs to \(\mathcal I_2\), so \(q\ge1\) and \(K_{\mathrm{ev}}>0\).

Direct integration of the three one-dimensional formulas gives
\[
\int_{-1/2}^{1/2}
\mathsf P_{j_{\mathrm{pol}}}(\upsilon)
\mathsf P_{j'_{\mathrm{pol}}}(\upsilon)\,d\upsilon
=
\mathbf1_{\{j_{\mathrm{pol}}=j'_{\mathrm{pol}}\}}
\qquad
(j_{\mathrm{pol}},j'_{\mathrm{pol}}\in\{0,1,2\}).
\]
Scaling each coordinate and multiplying these identities yields
\[
\int r_\kappa(x)r_{\kappa'}(x)\,d\nu_h(x)
=
\mathbf1_{\{\kappa=\kappa'\}}.
\]
Expanding the finite sums consequently gives polynomial reproduction and kernel energy:
\[
\int\mathcal K_h(x)\,r(x)^\top v\,d\nu_h(x)
=r_0^\top v
\qquad(v\in\mathbb R^q),
\]
\[
\int\mathcal K_h(x)^2\,d\nu_h(x)=\|r_0\|_2^2.
\]
The displayed polynomial formulas also give
\(|\mathsf P_{j_{\mathrm{pol}}}(\upsilon)|\le4\) on \([-1/2,1/2]\). Thus
\[
|r_\kappa(x)|\le4^d\quad(x\in C_h),
\qquad
|\mathcal K_h(x)|\le K_{\mathrm{ev}}\quad(x\in C_h),
\qquad
\|r_0\|_2^2\le K_{\mathrm{ev}}.
\]
All these identities hold for every \(h>0\).
% lean: CausalSmith.Stat.TwosamplePointcateAnnotationFrontier.oracle_polynomial_reproduction

\item Fix \(P\in\mathcal M\). Define the supplied-propensity response, its localized record contribution, the raw average, and its mean by
\[
Z_P(x,a_{\mathrm{obs}},y_{\mathrm{obs}})
=
\frac{a_{\mathrm{obs}}y_{\mathrm{obs}}}{e_P(x)}
-
\frac{(1-a_{\mathrm{obs}})y_{\mathrm{obs}}}{1-e_P(x)},
\]
\[
F_{P,h}(x,a_{\mathrm{obs}},y_{\mathrm{obs}})
=
w_h(x)\mathcal K_h(x)
Z_P(x,a_{\mathrm{obs}},y_{\mathrm{obs}}),
\]
\[
\overline T_{P,h}(\mathcal D_{n,m},U)
=
\frac1n\sum_{i=1}^nF_{P,h}(X_i,A_i,Y_i),
\qquad
\mu_{P,h}
=
\int\overline T_{P,h}\,d\mathsf E_{n,m}(P).
\]
These definitions apply to all \(x\in\mathbb R^d\), binary record coordinates, and every dataset and randomizer, using the quotient convention stated above.

On the design cube, overlap and the binary record coordinates give
\[
|Z_P(x,a_{\mathrm{obs}},y_{\mathrm{obs}})|\le\varepsilon^{-1}.
\]
Outside the cube, the zero extension of the propensity gives
\[
Z_P(x,a_{\mathrm{obs}},y_{\mathrm{obs}})
=-(1-a_{\mathrm{obs}})y_{\mathrm{obs}},
\qquad
|Z_P(x,a_{\mathrm{obs}},y_{\mathrm{obs}})|\le1\le\varepsilon^{-1}.
\]
Hence the response bound is global, and
\[
|F_{P,h}(x,a_{\mathrm{obs}},y_{\mathrm{obs}})|
\le\varepsilon^{-1}h^{-d}K_{\mathrm{ev}}.
\]
Borel measurability and this bound supply all integrability and finite second moments used below.

For \(0<h\le1\), exchangeability and the conditional-margin identities in
\cref{obj:def:primitive-class}, followed by cancellation of the overlap denominators, give
\[
\mathbb E_P\left[
w_h(X)\mathcal K_h(X)\frac{AY}{e_P(X)}
\right]
=
\int_{[0,1]^d}w_h(x)\mathcal K_h(x)\mu_{1,P}(x)\,dx,
\]
\[
\mathbb E_P\left[
w_h(X)\mathcal K_h(X)\frac{(1-A)Y}{1-e_P(X)}
\right]
=
\int_{[0,1]^d}w_h(x)\mathcal K_h(x)\mu_{0,P}(x)\,dx.
\]
Subtracting and using the canonical contrast identity gives
\[
\mathbb E_P[F_{P,h}(X,A,Y)]
=
\int\mathcal K_h(x)\tau_P(x)\,d\nu_h(x).
\]
Every labeled coordinate has this expectation, while the auxiliary block and randomizer integrate out as probability factors. Since \(n>0\),
\[
\mu_{P,h}
=
\int\mathcal K_h(x)\tau_P(x)\,d\nu_h(x).
\]
% lean: CausalSmith.Stat.TwosamplePointcateAnnotationFrontier.oracle_average_mean

\item We now bound this mean's bias over the full range \(0<h\le1\). Define the Taylor polynomial
\[
p(x)=
\sum_{|\kappa|\le p_*}
\frac{\partial^\kappa\tau_P(x_0)}{\kappa!}
(x-x_0)^\kappa,
\qquad
\kappa!=\prod_{\iota_{\mathrm{cov}}=1}^d\kappa_{\iota_{\mathrm{cov}}}!.
\]
Because \(p_*\le2\), the coarse polynomial basis spans \(p\). Thus there exists a coefficient vector \(\vartheta_{\mathrm{Tay}}\in\mathbb R^q\) with
\[
p_{\vartheta_{\mathrm{Tay}}}(x)
=
r(x)^\top\vartheta_{\mathrm{Tay}}
=
p(x)
\qquad(x\in\mathbb R^d),
\qquad
p_{\vartheta_{\mathrm{Tay}}}(x_0)=\tau_P(x_0).
\]

If \(\gamma=1\), then \(p_*=0\) and \(p(x)=\tau_P(x_0)\). The Hölder norm condition gives the Lipschitz bound, and the geometry of \(C_h\) gives
\[
|\tau_P(x)-p(x)|
\le L\|x-x_0\|_2
\le L\sqrt d\,h
\le L(\sqrt d+1)h
\qquad(x\in C_h).
\]

If \(1<\gamma\le3\), then \(1\le p_*\le2\). For each \(x\in C_h\), define the segment function
\[
g_{P,x}(\sigma_{\mathrm{seg}})=\tau_P\bigl(x_0+\sigma_{\mathrm{seg}}(x-x_0)\bigr),
\qquad 0\le\sigma_{\mathrm{seg}}\le1.
\]
The segment lies in the design cube. Expanding its \(p_*\)-th derivative in coordinate partials and using their \(s_o\)-Hölder modulus gives
\[
\left|g_{P,x}^{(p_*)}(\sigma_{\mathrm{seg}})-g_{P,x}^{(p_*)}(0)\right|
\le
(d+1)^{p_*}L
\|x-x_0\|_2^{p_*+s_o}\sigma_{\mathrm{seg}}^{s_o}.
\]
Indeed, the multinomial expansion contributes at most
\((d+1)^{p_*}\|x-x_0\|_2^{p_*}\), and each top-partial difference is bounded by
\(L\sigma_{\mathrm{seg}}^{s_o}\|x-x_0\|_2^{s_o}\).
The one-dimensional Lagrange remainder at degree \(p_*-1\), followed by subtraction of the degree-\(p_*\) Taylor term, supplies some
\(\xi_{P,x}\in(0,1)\) such that
\[
\tau_P(x)-p(x)
=
\frac{
g_{P,x}^{(p_*)}(\xi_{P,x})
-g_{P,x}^{(p_*)}(0)
}{p_*!}.
\]
Since \(\xi_{P,x}^{s_o}\le1\) and \(p_*!\ge1\),
\[
|\tau_P(x)-p(x)|
\le(d+1)^{p_*}L\|x-x_0\|_2^\gamma.
\]
Using \(\|x-x_0\|_2\le(d+1)h\) and
\(\gamma=p_*+s_o\le p_*+1\) yields
\[
|\tau_P(x)-p(x)|
\le L(d+1)^{2p_*+1}h^\gamma
\qquad(x\in C_h).
\]

Define the positive public constants
\[
C_{\mathrm{Tay}}=
\begin{cases}
L(\sqrt d+1),&\gamma=1,\\
L(d+1)^{2p_*+1},&1<\gamma\le3,
\end{cases}
\qquad
C_{\mathrm{bias}}=K_{\mathrm{ev}}C_{\mathrm{Tay}}.
\]
Both cases establish
\[
\sup_{x\in C_h}|\tau_P(x)-p_{\vartheta_{\mathrm{Tay}}}(x)|
\le C_{\mathrm{Tay}}h^\gamma.
\]
Polynomial reproduction and the mean identity therefore imply
\[
\begin{aligned}
|\mu_{P,h}-\tau_P(x_0)|
&=
\left|
\int\mathcal K_h(x)
\bigl(\tau_P(x)-p_{\vartheta_{\mathrm{Tay}}}(x)\bigr)
\,d\nu_h(x)
\right|\\
&\le
\int|\mathcal K_h(x)|
\left|\tau_P(x)-p_{\vartheta_{\mathrm{Tay}}}(x)\right|
\,d\nu_h(x)\\
&\le C_{\mathrm{bias}}h^\gamma.
\end{aligned}
\]
At \(h=h_n\), this becomes
\[
|\mu_{P,h_n}-\tau_P(x_0)|
\le C_{\mathrm{bias}}r_o(n).
\]
% lean: CausalSmith.Stat.TwosamplePointcateAnnotationFrontier.oracle_tuned_average_bias

\item The global response bound gives the recordwise squared bound
\[
F_{P,h}(X,A,Y)^2
\le\varepsilon^{-2}\bigl(w_h(X)\mathcal K_h(X)\bigr)^2.
\]
Under the uniform design, localization and the kernel-energy identity give
\[
\begin{aligned}
\mathbb E_P[F_{P,h}(X,A,Y)^2]
&\le
\varepsilon^{-2}\int_{[0,1]^d}
\bigl(w_h(x)\mathcal K_h(x)\bigr)^2\,dx\\
&=
\varepsilon^{-2}h^{-d}\int\mathcal K_h(x)^2\,d\nu_h(x)\\
&=
\varepsilon^{-2}h^{-d}\|r_0\|_2^2.
\end{aligned}
\]
Independence of the \(n\) labeled records yields
\[
\begin{aligned}
\int(\overline T_{P,h}-\mu_{P,h})^2\,d\mathsf E_{n,m}(P)
&=
\frac1n\operatorname{Var}_P(F_{P,h}(X,A,Y))\\
&\le\frac1n\mathbb E_P[F_{P,h}(X,A,Y)^2]\\
&\le\frac{\varepsilon^{-2}\|r_0\|_2^2}{nh^d}.
\end{aligned}
\]
The first inequality uses
\[
\operatorname{Var}_P(F_{P,h})
=
\mathbb E_P[F_{P,h}^2]
-\bigl(\mathbb E_P[F_{P,h}]\bigr)^2
\le\mathbb E_P[F_{P,h}^2].
\]
Cauchy--Schwarz on the probability experiment then gives
\[
\int|\overline T_{P,h}-\mu_{P,h}|\,d\mathsf E_{n,m}(P)
\le
\left(
\int(\overline T_{P,h}-\mu_{P,h})^2\,d\mathsf E_{n,m}(P)
\right)^{1/2}
\le
\sqrt{\frac{\varepsilon^{-2}\|r_0\|_2^2}{nh^d}}.
\]
At \(h=h_n\), use the target-coordinate bound and
\[
(nh_n^d)^{-1/2}=r_o(n)
\]
to obtain
\[
\int|\overline T_{P,h_n}-\mu_{P,h_n}|\,d\mathsf E_{n,m}(P)
\le\sqrt{\varepsilon^{-2}K_{\mathrm{ev}}}\,r_o(n)
\le C_{\mathrm{dev}}r_o(n),
\qquad
C_{\mathrm{dev}}=\sqrt{\varepsilon^{-2}K_{\mathrm{ev}}}+1>0.
\]
The estimate is uniform in \(P\), \(n\ge2\), and \(m\ge0\).
% lean: CausalSmith.Stat.TwosamplePointcateAnnotationFrontier.oracle_tuned_absolute_deviation

\item The probability-valued arm-mean conditions in \cref{obj:def:primitive-class} imply
\[
-1
\le
\mu_{1,P}(x_0)-\mu_{0,P}(x_0)
=
\tau_P(x_0)
\le1.
\]
The displayed smoother equals
\[
T^e(\mathcal D_{n,m},U,e_P)
=
\operatorname{clip}_{[-1,1]}
\bigl(\overline T_{P,h_n}(\mathcal D_{n,m},U)\bigr).
\]
Projection onto this interval contracts distance to any target inside it. Thus, for every dataset and randomizer,
\[
\begin{aligned}
|T^e(\mathcal D_{n,m},U,e_P)-\tau_P(x_0)|
&\le|\overline T_{P,h_n}(\mathcal D_{n,m},U)-\tau_P(x_0)|\\
&\le
|\overline T_{P,h_n}(\mathcal D_{n,m},U)-\mu_{P,h_n}|
+
|\mu_{P,h_n}-\tau_P(x_0)|.
\end{aligned}
\]
Integrating and applying the preceding two bounds gives
\[
\mathcal R^e_{n,m}(T^e,P)
\le
(C_{\mathrm{dev}}+C_{\mathrm{bias}})r_o(n).
\]
Define
\[
C_{\mathrm{up}}=C_{\mathrm{dev}}+C_{\mathrm{bias}}>0.
\]
The bounded measurable raw average and continuous clipping map give measurability of the supplied-propensity section. The clipping formula itself gives the pointwise guarantee
\[
\left|T^e(\mathcal D_{n,m},U,e_P)\right|\le1
\]
for every dataset and randomizer.
% lean: CausalSmith.Stat.TwosamplePointcateAnnotationFrontier.oracle_smoother_upper_rate

\item For completeness, the smoother has a total admissible extension to arbitrary supplied functions. Define
\[
\widetilde T^e(\mathcal D_{n,m},U,e)=
\begin{cases}
T^e(\mathcal D_{n,m},U,e),&e\text{ is Borel},\\
0,&e\text{ is not Borel}.
\end{cases}
\]
For every fixed Borel \(e\), the totalized reciprocal, record-coordinate maps, finite sum, and clipping are Borel; the other section is constant. Since each designated \(e_P\) is Borel, this extension agrees with the displayed smoother at every admitted law. It is therefore a candidate in \cref{obj:def:oracle-risk}, and the uniform bound proves
\[
R^e(n,m)
\le
\sup_{P\in\mathcal M}
\mathcal R^e_{n,m}(\widetilde T^e,P)
\le C_{\mathrm{up}}r_o(n).
\]

Finally, retain the lower constant already constructed and set
\[
C=\max\{c,C_{\mathrm{up}}\}.
\]
These are finite real constants depending on the fixed public parameters, with
\[
0<c\le C.
\]
Because \(r_o(n)\ge0\), enlarging the upper constant preserves both the minimax and individual-decision bounds. Combining the lower bound with subfamily monotonicity gives, for every \(n\ge2\) and \(m\ge0\),
\[
c\,r_o(n)
\le R^e_{1/2}(n,m)
\le R^e(n,m)
\le C\,r_o(n),
\]
and, for every \(P\in\mathcal M\),
\[
\mathcal R^e_{n,m}(T^e,P)\le C\,r_o(n).
\]
Together with the measurability and pointwise clipping bound above, these are all the asserted conclusions.
% lean: CausalSmith.Stat.TwosamplePointcateAnnotationFrontier.supplied_propensity
\end{enumerate}
\end{proof}

\begin{proof}[Proof of \cref{obj:thm:marked-component-certificate}]
Fix the public parameters. All constants constructed below depend only on these parameters. We use the functions and sign arrays prescribed by \cref{obj:def:marked-handle}, extending the displayed function formulas to \(\mathbb R^d\) when taking analytic bounds.

\begin{enumerate}
\item \textit{Regularity and membership.}
Set
\[
D_d:=2^d.
\]
The function \(\chi\) is smooth and flat at zero: each of its derivatives on the positive half-line is a polynomial in \(1/u\) multiplied by \(\exp(-1/u)\), and this expression tends to zero as \(u\downarrow0\). The denominator defining the scalar angular transition in \cref{obj:def:marked-handle} is positive everywhere. Consequently the prescribed sine and cosine pieces join smoothly, including at their support endpoints. The bump \(B\) is also smooth and compactly supported, and
\[
0\le B_h(x)\le1,\qquad B_h(x_0)=1.
\]
Define the finite derivative bounds
\[
K_\omega:=1+\sup_{u\in\mathbb R}|\omega_0'(u)|,
\qquad
K_B:=1+\sup_{w\in\mathbb R^d}\|\nabla B(w)\|_2,
\]
and put
\[
K_\phi:=K_B+dK_\omega,
\qquad
K_F:=2D_dK_\phi.
\]
Translation invariance gives the same derivative bound for every \(\omega_z\). Telescoping a product of factors bounded by one, and then using \(\delta\le h\), yields
\[
|\phi_{\mathbf z}(x)-\phi_{\mathbf z}(x')|
\le
\left(\frac{K_B}{h}+\frac{dK_\omega}{\delta}\right)
\|x-x'\|_2
\le
\frac{K_\phi}{\delta}\|x-x'\|_2.
\]

At each point, a nonzero frame entry has, in each coordinate, one of the two lattice indices adjacent to the scaled coordinate. Thus at most \(D_d\) entries are nonzero. Each entry has absolute value at most one, so
\[
|F_\lambda(x)|\le D_d,
\qquad
|G_\eta(x)|\le D_d.
\]
For a difference between two points, entries outside the union of their active index sets contribute zero. That union has at most \(2D_d\) elements. Therefore
\[
|F_\lambda(x)-F_\lambda(x')|
\le \frac{K_F}{\delta}\|x-x'\|_2,
\qquad
|G_\eta(x)-G_\eta(x')|
\le \frac{K_F}{\delta}\|x-x'\|_2.
\]
For a Hölder order \(s_{\mathrm{ord}}\in(0,1]\), define
\[
C_{F,s_{\mathrm{ord}}}:=2D_d+K_F+1.
\]
When \(\|x-x'\|_2\le\delta\), the Lipschitz bound and
\[
\frac{\|x-x'\|_2}{\delta}
\le
\left(\frac{\|x-x'\|_2}{\delta}\right)^{s_{\mathrm{ord}}}
\]
give the required Hölder difference bound. When \(\|x-x'\|_2>\delta\), use the bound \(2D_d\) instead. Since \(\delta\le1\), the supremum term is covered as well. With the norm of \cref{obj:def:holder-norm}, this proves
\[
\|F_\lambda\|_{H^{s_{\mathrm{ord}}}}
\le C_{F,s_{\mathrm{ord}}}\delta^{-s_{\mathrm{ord}}},
\qquad
\|G_\eta\|_{H^{s_{\mathrm{ord}}}}
\le C_{F,s_{\mathrm{ord}}}\delta^{-s_{\mathrm{ord}}}.
\]

For the macro bump, use its fixed compact profile \(B^2\). For every integer \(j\ge0\), define
\[
D_{B,j}:=
1+\sup_{w\in\mathbb R^d}\|\mathrm D^j(B^2)(w)\|_{\mathrm{op}},
\]
where \(\mathrm D^j\) is the total derivative of order \(j\), with the order-zero norm interpreted as absolute value. These constants are finite by smoothness and compact support. For \(s_{\mathrm{ord}}>0\), define
\[
j_{\mathrm{der}}:=\lceil s_{\mathrm{ord}}\rceil-1,
\qquad
\xi_{\mathrm{der}}:=s_{\mathrm{ord}}-j_{\mathrm{der}}\in(0,1],
\]
\[
M_{B,s_{\mathrm{ord}}}
:=
2\max\{D_{B,j_{\mathrm{der}}},D_{B,j_{\mathrm{der}}+1}\},
\]
\[
C_{B,s_{\mathrm{ord}}}
:=
\sum_{j=0}^{j_{\mathrm{der}}}D_{B,j}d^j
+
M_{B,s_{\mathrm{ord}}}d^{j_{\mathrm{der}}}+1.
\]
The derivative of order \(j_{\mathrm{der}}\) has both a bounded supremum and a Lipschitz bound supplied by the next derivative. Splitting at distance one gives
\[
\|\mathrm D^{j_{\mathrm{der}}}(B^2)(w)
-\mathrm D^{j_{\mathrm{der}}}(B^2)(w')\|_{\mathrm{op}}
\le
M_{B,s_{\mathrm{ord}}}\|w-w'\|_2^{\xi_{\mathrm{der}}}.
\]
Dilation by \(h\) multiplies a derivative of order \(j\) by \(h^{-j}\). Evaluating a derivative on coordinate directions bounds a partial derivative by its operator norm times \(d^j\). Hence, for \(0<h\le1\),
\[
|\partial^\kappa(B_h^2)(x)|
\le D_{B,|\kappa|}d^{|\kappa|}h^{-|\kappa|}
\le C_{B,s_{\mathrm{ord}}}h^{-s_{\mathrm{ord}}}
\qquad (|\kappa|\le j_{\mathrm{der}}),
\]
and
\[
|\partial^\kappa(B_h^2)(x)-\partial^\kappa(B_h^2)(x')|
\le
M_{B,s_{\mathrm{ord}}}d^{j_{\mathrm{der}}}
h^{-s_{\mathrm{ord}}}\|x-x'\|_2^{\xi_{\mathrm{der}}}
\qquad (|\kappa|=j_{\mathrm{der}}).
\]
Thus, in particular for the orders \(\beta\) and \(\gamma\),
\[
\|B_h^2\|_{H^{s_{\mathrm{ord}}}}
\le C_{B,s_{\mathrm{ord}}}h^{-s_{\mathrm{ord}}}.
\]

Define
\[
K_M:=
D_d+C_{F,\alpha}+C_{F,\beta}
+C_{B,\gamma}+C_{B,\beta}+1,
\qquad
\rho:=\min\left\{1,L-1/2,\frac12-\varepsilon\right\},
\qquad
c_M:=\frac{\rho}{32K_M}.
\]
The public domain gives \(0<\rho\le1\). Hence \(0<c_M\le\rho/32\le1/32\), and multiplying any of the five summands preceding the final \(1\) in \(K_M\) by \(c_M\) gives at most \(\rho/32\).

Suppose the amplitude and contrast inequalities hold with \(c_M\). Since \(0<\delta\le h\le1/2\),
\[
a\le c_M,\qquad b\le c_M,\qquad ab\le c_M,
\]
\[
a\delta^{-\alpha}\le c_M,\qquad
b\delta^{-\beta}\le c_M,\qquad
ab h^{-\gamma}\le c_M.
\]
Also \(\beta\le\gamma\) and \(h\le1\), so
\[
ab\le c_Mh^\gamma\le c_Mh^\beta,
\qquad
ab h^{-\beta}\le c_M.
\]
The prescribed probabilities therefore satisfy
\[
|e_\lambda(x)-1/2|
\le aD_d\le\frac{\rho}{32}\le\frac14,
\qquad
aD_d\le\frac{\rho}{32}\le\frac12-\varepsilon,
\]
\[
|\mu_{0,\theta,\eta}(x)-1/2|,
\ |\mu_{1,\theta,\eta}(x)-1/2|
\le bD_d+ab
\le\frac{\rho}{32}+c_M\le\frac1{16}\le\frac38.
\]
These bounds imply the overlap and both probability-valued arm-mean conditions.

Subadditivity and homogeneity of the Hölder norm give
\[
\|aF_\lambda\|_{H^\alpha}\le C_{F,\alpha}c_M,
\qquad
\|bG_\eta\|_{H^\beta}\le C_{F,\beta}c_M,
\]
\[
\|\tau_\theta\|_{H^\gamma}
\le2C_{B,\gamma}c_M,
\qquad
\|\theta abB_h^2\|_{H^\beta}
\le C_{B,\beta}c_M.
\]
Consequently
\[
\|e_\lambda\|_{H^\alpha}
\le\frac12+\frac{\rho}{32}\le L,
\]
\[
\|\mu_{0,\theta,\eta}\|_{H^\beta}
\le\frac12+\frac{\rho}{16}\le L,
\qquad
\|\tau_\theta\|_{H^\gamma}\le\frac{\rho}{16}\le L.
\]
All the prescribed functions are Borel. The independent conditional Bernoulli construction gives uniform covariates and conditional exchangeability, and its conditional margins are exactly the prescribed functions. Together with the preceding bounds, these facts establish every condition in \cref{obj:def:primitive-class}.

To identify the canonical contrast, the prescribed control mean is continuous, and the prescribed treated mean is continuous because it is the sum of the control mean and \(\tau_\theta\). Their admissible designated versions agree with them almost surely by uniqueness of conditional margins. Two continuous versions that agree almost everywhere under the uniform design agree everywhere on the cube: a nonzero continuous difference would persist on a relative neighborhood of positive uniform measure, including at a boundary point. Thus
\[
P_{\theta,\lambda,\eta}\in\mathcal M,
\qquad
\tau_{P_{\theta,\lambda,\eta}}(x)
=
\mu_{1,\theta,\eta}(x)-\mu_{0,\theta,\eta}(x)
=
\tau_\theta(x)
\quad(x\in[0,1]^d).
\]
% lean: CausalSmith.Stat.TwosamplePointcateAnnotationFrontier.marked_membership

\item \textit{Normalization and singleton cancellation.}
Define expectation over the finite sign arrays by
\[
\mathbb E_\theta Q
:=
\sum_{\lambda,\eta}
w_\theta(\lambda,\eta)Q(\lambda,\eta).
\]
At each index, direct summation over the four sign pairs gives
\[
\frac{1+\theta l s_\eta/2}{4}\ge0
\quad(l,s_\eta\in\{-1,+1\}),
\qquad
\sum_{l,s_\eta\in\{-1,+1\}}\frac{1+\theta l s_\eta/2}{4}=1.
\]
Multiplying these identities over the finite index set proves
\[
w_\theta(\lambda,\eta)\ge0,
\qquad
\sum_{\lambda,\eta}w_\theta(\lambda,\eta)=1.
\]
The same finite-product summation gives
\[
\mathbb E_\theta\lambda_{\mathbf z}
=
\mathbb E_\theta\eta_{\mathbf z}=0,
\qquad
\mathbb E_\theta[\lambda_{\mathbf z}\eta_{\mathbf w}]
=
\begin{cases}
\theta/2,&\mathbf z=\mathbf w,\\
0,&\mathbf z\ne\mathbf w.
\end{cases}
\]
Indeed, matching indices use the scalar sign-pair correlation, while distinct indices include a zero-mean scalar factor.

The sine-cosine identity on each lattice interval gives a sum of squares equal to one. Tensoring it and multiplying by \(B_h^2\) yields
\[
\sum_{\mathbf z\in\mathcal I}\phi_{\mathbf z}(x)^2=B_h(x)^2.
\]
For \(x\in C_h\), all nonzero tensor entries lie in the prescribed finite index set; for \(x\notin C_h\), both sides vanish. Expanding the finite signed sums therefore gives
\[
\mathbb E_\theta F_\lambda(x)
=
\mathbb E_\theta G_\eta(x)=0,
\qquad
\mathbb E_\theta[F_\lambda(x)G_\eta(x)]
=\frac{\theta}{2}B_h(x)^2.
\]

Set
\[
c_S:=\frac1{64D_d}.
\]
If the amplitude inequalities hold with \(c_S\), then
\[
a,b\le c_S,\qquad
aD_d,bD_d\le\frac1{64},
\qquad
ab\le\frac1{4096}.
\]
Thus the raw propensity and both raw arm probabilities lie in \([0,1]\), at every covariate, so their Bernoulli atom formulas apply directly.

For a labeled record, define the treatment and observed-response spins by
\[
s^A:=2A-1,\qquad
Y:=AY(1)+(1-A)Y(0),\qquad
s^Y:=2Y-1.
\]
Relative to the uniform probability measure on the four spin pairs, its likelihood is
\[
L_{\theta,\lambda,\eta}(x;s^A,s^Y)
:=
(1+2s^AaF_\lambda(x))
(1+2s^YbG_\eta(x)+s^As^Y\tau_\theta(x)).
\]
This follows by multiplying the two Bernoulli atom probabilities, each divided by its fair atom probability. Expanding the product gives
\[
\begin{aligned}
\mathbb E_\theta L_{\theta,\lambda,\eta}(x;s^A,s^Y)
={}&1+s^As^Y\tau_\theta(x)\\
&+\bigl(2s^Aa+2(s^A)^2s^Ya\tau_\theta(x)\bigr)
  \mathbb E_\theta F_\lambda(x)\\
&+2s^Yb\,\mathbb E_\theta G_\eta(x)
+4s^As^Yab\,\mathbb E_\theta[F_\lambda(x)G_\eta(x)]\\
={}&1+s^As^Y
\bigl(-2\theta abB_h(x)^2+2\theta abB_h(x)^2\bigr)
=1.
\end{aligned}
\]
Hence every atom of the singleton labeled mixture has probability \(1/4\). This proves its uniform conditional law for every hypothesis and every cube covariate.
% lean: CausalSmith.Stat.TwosamplePointcateAnnotationFrontier.singleton_cancellation

\item \textit{Equality on uninformative components.}
For integers \(n\ge2\) and \(m\ge0\), write
\[
N:=n+m,\qquad
\mathbf x:=(x_1,\ldots,x_N)\in([0,1]^d)^N.
\]
The first \(n\) indices are labeled. A spin assignment has domain
\[
\mathbf s\in
\bigl(\{-1,+1\}^2\bigr)^n
\times\{-1,+1\}^m.
\]
Define the record likelihoods by
\[
R_{\theta,\lambda,\eta,i}(\mathbf x,\mathbf s)
:=
\begin{cases}
L_{\theta,\lambda,\eta}(x_i;s_i^A,s_i^Y),&i\le n,\\
1+2s_i^AaF_\lambda(x_i),&i>n,
\end{cases}
\]
and, for every subset \(\mathcal V\subseteq\{1,\ldots,N\}\), define
\[
g_{\theta,\mathcal V}(\mathbf x,\mathbf s)
:=
\mathbb E_\theta
\prod_{i\in\mathcal V}
R_{\theta,\lambda,\eta,i}(\mathbf x,\mathbf s).
\]
Empty products equal one, so \(g_{\theta,\varnothing}=1\).

The sign flip
\[
\mathcal F_\eta(\lambda,\eta):=(\lambda,-\eta)
\]
is a bijection, and its weights satisfy
\[
w_{+1}(\lambda,\eta)=w_{-1}(\lambda,-\eta).
\]
Every auxiliary likelihood is unchanged by this flip, because it depends on \(\lambda\) alone. Reindexing the finite sum consequently gives
\[
g_{+1,\mathcal V}(\mathbf x,\mathbf s)
=
g_{-1,\mathcal V}(\mathbf x,\mathbf s)
\qquad
\text{if }\mathcal V\subseteq\{n+1,\ldots,N\}.
\]
If \(\mathcal V=\{i\}\) with \(i\le n\), singleton cancellation instead gives
\[
g_{+1,\{i\}}(\mathbf x,\mathbf s)
=
g_{-1,\{i\}}(\mathbf x,\mathbf s)=1.
\]
Thus we may take
\[
c_U:=c_S.
\]
Under its amplitude conditions, every auxiliary-only component and every singleton labeled component has equal densities under the two hypotheses. Equality for all spin assignments gives equality of their full conditional laws.
% lean: CausalSmith.Stat.TwosamplePointcateAnnotationFrontier.uninformative_component_equality

\item \textit{The component density difference.}
Define
\[
c_\Delta:=\frac1{64D_d},
\qquad
C_\Delta:=8D_d^2+6.
\]
Assume the amplitude inequalities with \(c_\Delta\). Fix \(\mathbf x,\mathbf s\) and a subset \(\mathcal V\), and define its size by
\[
p_{\mathrm{cmp}}:=|\mathcal V|.
\]
For every record, introduce the contrast-deleted factors
\[
f_{\lambda,i}:=1+2s_i^AaF_\lambda(x_i),
\qquad
g_{\eta,i}:=
\begin{cases}
1+2s_i^YbG_\eta(x_i),&i\le n,\\
1,&i>n.
\end{cases}
\]
The bounds from the previous steps give
\[
|f_{\lambda,i}-1|\le2aD_d,\qquad
|g_{\eta,i}-1|\le2bD_d,
\qquad
|f_{\lambda,i}|,|g_{\eta,i}|\le\frac98.
\]
For labeled records,
\[
R_{\theta,\lambda,\eta,i}
=f_{\lambda,i}
\bigl(g_{\eta,i}+s_i^As_i^Y\tau_\theta(x_i)\bigr),
\qquad
|\tau_\theta(x_i)|\le2ab.
\]
For auxiliary records, \(R_{\theta,\lambda,\eta,i}=f_{\lambda,i}\). Since \(ab\le1/4096\), both cases satisfy
\[
|R_{\theta,\lambda,\eta,i}|
\le\frac32,
\qquad
|R_{\theta,\lambda,\eta,i}-f_{\lambda,i}g_{\eta,i}|
\le3ab,
\qquad
|f_{\lambda,i}g_{\eta,i}|\le\frac32.
\]

Define
\[
\mathcal A_{\mathcal V}(\lambda)
:=\prod_{i\in\mathcal V}f_{\lambda,i},
\qquad
\mathcal G_{\mathcal V}(\eta)
:=\prod_{i\in\mathcal V}g_{\eta,i},
\]
\[
U_{\mathcal V}
:=p_{\mathrm{cmp}}(2aD_d)(9/8)^{p_{\mathrm{cmp}}},
\qquad
W_{\mathcal V}
:=p_{\mathrm{cmp}}(2bD_d)(9/8)^{p_{\mathrm{cmp}}},
\]
\[
D_{\mathcal V}
:=p_{\mathrm{cmp}}(3ab)(3/2)^{p_{\mathrm{cmp}}}.
\]
Telescoping a product, using an upper bound \(D_{\mathrm{prod}}\ge1\) on both sets of factors, gives
\[
\left|
\prod_{i\in\mathcal V}f_{\mathrm{tel},i}
-\prod_{i\in\mathcal V}g_{\mathrm{tel},i}
\right|
\le
D_{\mathrm{prod}}^{|\mathcal V|}
\sum_{i\in\mathcal V}
|f_{\mathrm{tel},i}-g_{\mathrm{tel},i}|.
\]
Here the telescoping sum replaces one factor at a time; each remaining product is bounded by \(D_{\mathrm{prod}}^{|\mathcal V|-1}\), which is at most the displayed power. For the empty set both sides are zero. Applying this inequality gives
\[
|\mathcal A_{\mathcal V}-1|\le U_{\mathcal V},
\qquad
|\mathcal G_{\mathcal V}-1|\le W_{\mathcal V},
\]
\[
\left|
\prod_{i\in\mathcal V}R_{\theta,\lambda,\eta,i}
-\mathcal A_{\mathcal V}\mathcal G_{\mathcal V}
\right|
\le D_{\mathcal V}.
\]

Flipping \(\eta\) preserves every function of \(\lambda\) and changes the positive weights to the negative weights. Flipping \(\lambda\) similarly preserves every function of \(\eta\). Consequently
\[
\mathbb E_{+1}\mathcal A_{\mathcal V}
=\mathbb E_{-1}\mathcal A_{\mathcal V},
\qquad
\mathbb E_{+1}\mathcal G_{\mathcal V}
=\mathbb E_{-1}\mathcal G_{\mathcal V}.
\]
Expanding the centered product and using normalization therefore yields
\[
\begin{aligned}
&\mathbb E_{+1}[\mathcal A_{\mathcal V}\mathcal G_{\mathcal V}]
-\mathbb E_{-1}[\mathcal A_{\mathcal V}\mathcal G_{\mathcal V}]\\
&\qquad=
\mathbb E_{+1}[(\mathcal A_{\mathcal V}-1)(\mathcal G_{\mathcal V}-1)]
-\mathbb E_{-1}[(\mathcal A_{\mathcal V}-1)(\mathcal G_{\mathcal V}-1)].
\end{aligned}
\]
Each expectation on the right has absolute value at most
\(U_{\mathcal V}W_{\mathcal V}\). Adding the two contrast-deletion errors gives
\[
|g_{+1,\mathcal V}-g_{-1,\mathcal V}|
\le D_{\mathcal V}+2U_{\mathcal V}W_{\mathcal V}+D_{\mathcal V}.
\]
Finally,
\[
(9/8)^{p_{\mathrm{cmp}}}(9/8)^{p_{\mathrm{cmp}}}
\le(3/2)^{p_{\mathrm{cmp}}},
\qquad
p_{\mathrm{cmp}}\le p_{\mathrm{cmp}}^2
\quad(p_{\mathrm{cmp}}\in\mathbb N).
\]
Substitution proves
\[
|g_{+1,\mathcal V}(\mathbf x,\mathbf s)
-g_{-1,\mathcal V}(\mathbf x,\mathbf s)|
\le
C_\Delta(3/2)^{p_{\mathrm{cmp}}}
p_{\mathrm{cmp}}^2ab.
\]
For \(p_{\mathrm{cmp}}=0\), both densities equal one and the bound is zero.
% lean: CausalSmith.Stat.TwosamplePointcateAnnotationFrontier.component_information_bound

\item \textit{A positive floor and component Hellinger distance.}
Define
\[
c_H:=\min\left\{c_\Delta,\frac1{64D_d}\right\},
\qquad
K_{\mathrm{info}}:=C_\Delta^2.
\]
Under the amplitude inequalities with \(c_H\),
\[
aD_d\le\frac1{16},
\qquad
bD_d+ab\le\frac1{16}.
\]
Thus the treatment factor of a labeled likelihood is at least \(7/8\), and its outcome factor is also at least \(7/8\). Their product is at least \(1/2\). An auxiliary likelihood is at least \(7/8\), hence at least \(1/2\). Averaging the resulting product floor with nonnegative normalized weights gives
\[
g_{\theta,\mathcal V}(\mathbf x,\mathbf s)
\ge(1/2)^{p_{\mathrm{cmp}}}.
\]

Define the common probability reference on the record spins by
\[
\mu_{\mathrm{spin}}
:=
\bigotimes_{i=1}^{n}\operatorname{Unif}(\{-1,+1\}^2)
\otimes
\bigotimes_{i=n+1}^{N}\operatorname{Unif}(\{-1,+1\}),
\]
and define
\[
H_{\mathcal V}^2(\mathbf x)
:=
\int
\left(
\sqrt{g_{+1,\mathcal V}(\mathbf x,\mathbf s)}
-\sqrt{g_{-1,\mathcal V}(\mathbf x,\mathbf s)}
\right)^2
\,d\mu_{\mathrm{spin}}(\mathbf s).
\]
For nonnegative numbers bounded below by \(t_{\mathrm{floor}}>0\), the identity
\[
(f_{\mathrm{den}}-g_{\mathrm{den}})^2
=
(\sqrt{f_{\mathrm{den}}}-\sqrt{g_{\mathrm{den}}})^2
(\sqrt{f_{\mathrm{den}}}+\sqrt{g_{\mathrm{den}}})^2
\]
and the denominator bound
\[
(\sqrt{f_{\mathrm{den}}}+\sqrt{g_{\mathrm{den}}})^2
\ge4t_{\mathrm{floor}}
\]
give
\[
(\sqrt{f_{\mathrm{den}}}-\sqrt{g_{\mathrm{den}}})^2
\le
\frac{(f_{\mathrm{den}}-g_{\mathrm{den}})^2}{4t_{\mathrm{floor}}}.
\]
Apply this with the preceding component floor and density difference, and integrate against the probability reference. We obtain
\[
\begin{aligned}
H_{\mathcal V}^2(\mathbf x)
&\le
\frac{
\bigl(C_\Delta(3/2)^{p_{\mathrm{cmp}}}
p_{\mathrm{cmp}}^2ab\bigr)^2
}{
4(1/2)^{p_{\mathrm{cmp}}}
}\\
&=
\frac{C_\Delta^2}{4}(9/2)^{p_{\mathrm{cmp}}}
p_{\mathrm{cmp}}^4a^2b^2\\
&\le
K_{\mathrm{info}}(9/2)^{p_{\mathrm{cmp}}}
p_{\mathrm{cmp}}^4a^2b^2.
\end{aligned}
\]
% lean: CausalSmith.Stat.TwosamplePointcateAnnotationFrontier.marked_component_hellinger_bound

\item \textit{Conditional factorization and the count bound.}
For fixed covariates, define the augmented graph
\[
\mathscr G_{\mathbf x}
:=
\left(
\{1,\ldots,N\},
\left\{\{i,j\}:
i\ne j,\ x_i,x_j\in C_h,\
\|x_i-x_j\|_\infty\le2\delta
\right\}
\right),
\]
and let
\[
\mathscr K_{\mathbf x}
:=
\{\mathcal V:\mathcal V
\text{ is the vertex set of a connected component of }
\mathscr G_{\mathbf x}\}.
\]
The cube \(C_h\) is that of \cref{obj:synth_3}. Records outside \(C_h\) are isolated in this augmented graph and have likelihood one.

If \(\phi_{\mathbf z}(x_i)\ne0\), then
\[
x_i\in C_h,
\qquad
\left|\frac{x_{i,j_{\mathrm{coord}}}-x_{0,j_{\mathrm{coord}}}}{\delta}-z_{j_{\mathrm{coord}}}\right|<1
\quad(j_{\mathrm{coord}}\in\{1,\ldots,d\}).
\]
Two records sharing a nonzero frame entry therefore satisfy
\[
\|x_i-x_j\|_\infty\le2\delta,
\]
and belong to the same component. Assign each frame index active at any record to this unique component; inactive indices contribute a factor one. The product prior is independent across indices, and each component likelihood depends solely on its assigned indices. Regrouping the finite sign sums by these disjoint groups consequently proves
\[
g_\theta(\mathbf x,\mathbf s)
:=
\mathbb E_\theta
\prod_{i=1}^{N}R_{\theta,\lambda,\eta,i}(\mathbf x,\mathbf s)
=
\prod_{\mathcal V\in\mathscr K_{\mathbf x}}
g_{\theta,\mathcal V}(\mathbf x,\mathbf s).
\]
Each component density depends solely on its component spins. Integrating the fixed-sign record product over those fair spins gives one, because every record likelihood is normalized. Averaging signs therefore gives
\[
g_{\theta,\mathcal V}\ge0,
\qquad
\int g_{\theta,\mathcal V}\,d\mu_{\mathrm{spin}}=1.
\]

Define the component affinities by
\[
\mathcal A_{\mathrm{aff},\mathcal V}(\mathbf x)
:=
\int\sqrt{g_{+1,\mathcal V}g_{-1,\mathcal V}}\,
d\mu_{\mathrm{spin}}.
\]
Normalization and \(2\sqrt{f_{\mathrm{den}}g_{\mathrm{den}}}
\le f_{\mathrm{den}}+g_{\mathrm{den}}\) imply
\[
0\le\mathcal A_{\mathrm{aff},\mathcal V}\le1,
\qquad
H_{\mathcal V}^2
=2(1-\mathcal A_{\mathrm{aff},\mathcal V}).
\]
Integration over disjoint spin blocks multiplies the affinities. The elementary finite-product inequality
\[
1-\prod_{\mathcal V\in\mathscr K_{\mathbf x}}
\mathcal A_{\mathrm{aff},\mathcal V}
\le
\sum_{\mathcal V\in\mathscr K_{\mathbf x}}
(1-\mathcal A_{\mathrm{aff},\mathcal V})
\]
follows by induction from the fact that all the affinities lie in \([0,1]\). Hence
\[
\int(\sqrt{g_{+1}}-\sqrt{g_{-1}})^2\,d\mu_{\mathrm{spin}}
\le
\sum_{\mathcal V\in\mathscr K_{\mathbf x}}
H_{\mathcal V}^2(\mathbf x).
\]

Set
\[
c_I:=\min\{c_H,c_U\}.
\]
Its amplitude conditions permit both the component bound and the uninformative-component equalities. A component containing a labeled record and having size less than two has size one, so its distance is zero by singleton cancellation. A component containing no labeled record is auxiliary-only, so its distance is also zero. Define the informative component counts, for every integer \(p_{\mathrm{cmp}}\ge0\), by
\[
C_{p_{\mathrm{cmp}}}^{\mathrm L}(\mathbf x)
:=
\#\left\{
\mathcal V\in\mathscr K_{\mathbf x}:
|\mathcal V|=p_{\mathrm{cmp}},\
\mathcal V\cap\{1,\ldots,n\}\ne\varnothing
\right\}.
\]
Regrouping the remaining component sum by size gives
\[
\int(\sqrt{g_{+1}}-\sqrt{g_{-1}})^2\,d\mu_{\mathrm{spin}}
\le
K_{\mathrm{info}}a^2b^2
\sum_{p_{\mathrm{cmp}}=2}^{N}
(9/2)^{p_{\mathrm{cmp}}}
p_{\mathrm{cmp}}^4C_{p_{\mathrm{cmp}}}^{\mathrm L}(\mathbf x).
\]
The graph \(\mathscr G_{\mathbf x}\) is exactly the graph in the statement: it retains all \(N\) vertices, including the exterior singleton components, whose density factors equal one.
% lean: CausalSmith.Stat.TwosamplePointcateAnnotationFrontier.marked_conditional_hellinger_count_bound

\item \textit{The labeled-root tree series.}
Define the common covariate probability measure by
\[
\nu_{\mathrm{cov}}
:=
\operatorname{Unif}([0,1]^d)^{\otimes N}.
\]
For \(p_{\mathrm{cmp}}\ge2\), a component counted by
\(C_{p_{\mathrm{cmp}}}^{\mathrm L}\) has a labeled root and a spanning tree on its vertex set. For a fixed root, there are
\(\binom{N-1}{p_{\mathrm{cmp}}-1}\) choices of the other vertices. Orient a spanning tree toward the root. Each nonroot vertex chooses a parent among the \(p_{\mathrm{cmp}}\) selected vertices, so the number of valid rooted trees is at most
\(p_{\mathrm{cmp}}^{p_{\mathrm{cmp}}-1}\).

For a fixed valid tree, integrate its leaf covariates successively. Each child is constrained to a maximum-norm box of side \(4\delta\) around its parent; intersection with the unit cube has uniform mass at most \((4\delta)^d\). The root lies in \(C_h\), whose uniform mass is \(h^d\). Thus the probability of this tree event is at most
\[
h^d\bigl((4\delta)^d\bigr)^{p_{\mathrm{cmp}}-1}.
\]
Dropping the remaining component conditions enlarges the event. Summing these tree witnesses over the \(n\) possible labeled roots gives
\[
\int C_{p_{\mathrm{cmp}}}^{\mathrm L}(\mathbf x)\,
d\nu_{\mathrm{cov}}(\mathbf x)
\le
nh^d
\binom{N-1}{p_{\mathrm{cmp}}-1}
p_{\mathrm{cmp}}^{p_{\mathrm{cmp}}-1}
\bigl((4\delta)^d\bigr)^{p_{\mathrm{cmp}}-1}.
\]
For \(p_{\mathrm{cmp}}>N\), the count and the binomial coefficient are both zero.

The binomial bound and the exponential series imply
\[
\binom{N-1}{p_{\mathrm{cmp}}-1}
\le\frac{N^{p_{\mathrm{cmp}}-1}}{(p_{\mathrm{cmp}}-1)!},
\qquad
\frac{p_{\mathrm{cmp}}^{p_{\mathrm{cmp}}-1}}
{(p_{\mathrm{cmp}}-1)!}
\le \exp(p_{\mathrm{cmp}}).
\]
For the second inequality, its left side is one nonnegative summand of the exponential series at \(p_{\mathrm{cmp}}\). Consequently
\[
\binom{N-1}{p_{\mathrm{cmp}}-1}
p_{\mathrm{cmp}}^{p_{\mathrm{cmp}}-1}
\le N^{p_{\mathrm{cmp}}-1}\exp(p_{\mathrm{cmp}}).
\]

Define
\[
\kappa_{\mathrm{tree}}:=\frac92,
\qquad
S_{\mathrm{tree}}
:=
\sum_{j=0}^{\infty}(j+2)^4\,2^{-j},
\qquad
E_{\mathrm{occ}}
:=
\kappa_{\mathrm{tree}}\exp(1)4^d,
\]
\[
c_0:=\frac1{2E_{\mathrm{occ}}},
\qquad
C_{\mathrm{tree}}
:=
(S_{\mathrm{tree}}+1)
\bigl(\kappa_{\mathrm{tree}}\exp(1)\bigr)^2\,4^d.
\]
The series defining \(S_{\mathrm{tree}}\) converges because it is a shifted polynomial-geometric series. These definitions give \(c_0>0\) and \(C_{\mathrm{tree}}>0\).

Under \(N\delta^d\le c_0\), define
\[
z_{\mathrm{occ}}:=E_{\mathrm{occ}}N\delta^d,
\qquad
0\le z_{\mathrm{occ}}\le\frac12.
\]
Applying the binomial-tree bound with \(p_{\mathrm{cmp}}=j+2\) gives
\[
\begin{aligned}
&\kappa_{\mathrm{tree}}^{j+2}(j+2)^4
\binom{N-1}{j+1}(j+2)^{j+1}
\bigl((4\delta)^d\bigr)^{j+1}\\
&\qquad\le
\kappa_{\mathrm{tree}}\exp(1)(j+2)^4
z_{\mathrm{occ}}^{j+1}.
\end{aligned}
\]
For every finite truncation,
\[
\sum_{j=0}^{N-1}(j+2)^4z_{\mathrm{occ}}^{j+1}
\le
z_{\mathrm{occ}}\sum_{j=0}^{N-1}(j+2)^4\,2^{-j}
\le z_{\mathrm{occ}}S_{\mathrm{tree}}.
\]
Multiplying the tree-count inequality by its size weight and summing therefore yields
\[
\begin{aligned}
&\sum_{j=0}^{N-1}
\kappa_{\mathrm{tree}}^{j+2}(j+2)^4
\int C_{j+2}^{\mathrm L}\,d\nu_{\mathrm{cov}}\\
&\qquad\le
nh^d\kappa_{\mathrm{tree}}\exp(1)
z_{\mathrm{occ}}S_{\mathrm{tree}}\\
&\qquad\le C_{\mathrm{tree}}nNh^d\delta^d.
\end{aligned}
\]
The final inequality replaces \(S_{\mathrm{tree}}\) by
\(S_{\mathrm{tree}}+1\) and substitutes the definition of
\(z_{\mathrm{occ}}\).
% lean: CausalSmith.Stat.TwosamplePointcateAnnotationFrontier.labeled_component_series_bound

\item \textit{Integration for the original experiment.}
The covariate distribution is \(\nu_{\mathrm{cov}}\) for every sign realization and hypothesis. Conditional on these covariates, the original records have density \(g_\theta\) against \(\mu_{\mathrm{spin}}\). Thus, after grouping covariates and spins, the mixture laws have densities \(g_{+1}\) and \(g_{-1}\) against
\(\nu_{\mathrm{cov}}\otimes\mu_{\mathrm{spin}}\).

For two densities against a common dominating measure, substitution into \cref{obj:def:hellinger} gives the integral of their squared square-root difference. Applying this identity and then integrating over covariates yields
\[
H^2(M_{+1},M_{-1})
=
\int
\left[
\int(\sqrt{g_{+1}(\mathbf x,\mathbf s)}
-\sqrt{g_{-1}(\mathbf x,\mathbf s)})^2
\,d\mu_{\mathrm{spin}}(\mathbf s)
\right]d\nu_{\mathrm{cov}}(\mathbf x).
\]
The finite graph counts are measurable, and all functions in the finite weighted count sum are integrable. Integrating the conditional bound therefore gives
\[
H^2(M_{+1},M_{-1})
\le
K_{\mathrm{info}}a^2b^2
\sum_{p_{\mathrm{cmp}}=2}^{N}
(9/2)^{p_{\mathrm{cmp}}}p_{\mathrm{cmp}}^4
\int C_{p_{\mathrm{cmp}}}^{\mathrm L}\,d\nu_{\mathrm{cov}}.
\]
Every term is nonnegative, so extending the size sum to the shifted range used above gives
\[
\begin{aligned}
H^2(M_{+1},M_{-1})
&\le
K_{\mathrm{info}}a^2b^2
\sum_{j=0}^{N-1}
(9/2)^{j+2}(j+2)^4
\int C_{j+2}^{\mathrm L}\,d\nu_{\mathrm{cov}}\\
&\le
K_{\mathrm{info}}C_{\mathrm{tree}}
nNh^d\delta^da^2b^2.
\end{aligned}
\]
Define
\[
C:=K_{\mathrm{info}}C_{\mathrm{tree}}>0.
\]
This proves the required information estimate under the amplitude conditions with \(c_I\) and the occupancy condition with \(c_0\).
% lean: CausalSmith.Stat.TwosamplePointcateAnnotationFrontier.marked_information_bound

\item \textit{Common calibration and the remaining equalities.}
Choose
\[
c:=\min\{c_M,\min\{c_S,\min\{c_U,c_I\}\}\}.
\]
All four constants are positive and at most one, so
\[
0<c\le1,
\qquad
c\le c_M,\quad c\le c_S,\quad c\le c_U,\quad c\le c_I.
\]
Fix a tuple satisfying the theorem's hypotheses with this \(c\). Since \(\delta>0\) and \(h>0\), multiplication by their real powers preserves these comparisons. For each
\(c'\in\{c_M,c_S,c_U,c_I\}\),
\[
a\le c\delta^\alpha\le c'\delta^\alpha,\qquad
b\le c\delta^\beta\le c'\delta^\beta,\qquad
ab\le ch^\gamma\le c'h^\gamma.
\]
The membership, singleton, uninformative-component, and information estimates established above therefore apply simultaneously. Weight normalization makes the finite pushforward priors probability measures, and the pointwise membership conclusion supports them on \(\mathcal M\).

For the propensity distribution, the sign flip preserves the covariate-treatment marginal:
\[
(P_{+1,\lambda,\eta})_{XA}
=
(P_{-1,\lambda,-\eta})_{XA}.
\]
The designated propensity witnesses of these admitted laws are continuous and are versions of the same conditional treatment probability. Uniqueness of conditional margins followed by continuity under the full-support uniform design gives
\[
e_{P_{+1,\lambda,\eta}}(x)
=
e_{P_{-1,\lambda,-\eta}}(x)
\qquad(x\in[0,1]^d).
\]
Reindexing the finite pushforwards using
\(w_{+1}(\lambda,\eta)=w_{-1}(\lambda,-\eta)\) proves equality of the distributions of the entire propensity functions restricted to the cube.

At the target, \(B_h(x_0)=1\), and hence
\[
\tau_\theta(x_0)=-2\theta ab,
\qquad
\tau_{-1}(x_0)-\tau_{+1}(x_0)=4ab.
\]
Because \(a,b>0\) and \(c\le1\le4\),
\[
cab=c(ab)\le4(ab)=4ab.
\]

Finally, for every integer \(k_{\mathrm{aux}}\ge0\), reindexing by the same sign flip gives
\[
\begin{aligned}
\int P_{XA}^{\otimes k_{\mathrm{aux}}}\,d\varpi_{+1}(P)
&=
\sum_{\lambda,\eta}
w_{+1}(\lambda,\eta)
(P_{+1,\lambda,\eta})_{XA}^{\otimes k_{\mathrm{aux}}}\\
&=
\sum_{\lambda,\eta}
w_{-1}(\lambda,\eta)
(P_{-1,\lambda,\eta})_{XA}^{\otimes k_{\mathrm{aux}}}\\
&=
\int P_{XA}^{\otimes k_{\mathrm{aux}}}\,d\varpi_{-1}(P).
\end{aligned}
\]
For any fixed auxiliary covariates, their conditional mixture kernels are
\[
\sum_{\lambda,\eta}
w_\theta(\lambda,\eta)
\bigotimes_{i=1}^{k_{\mathrm{aux}}}
\operatorname{Bernoulli}(e_\lambda(x_i)).
\]
The factors are unchanged by the sign flip, so these kernels also agree for every covariate collection. When \(k_{\mathrm{aux}}=0\), the products are the unit probability measure on the empty tuple, and normalization gives the same equality.

The singleton calculation supplies the stated uniform labeled conditional laws. The conditional factorization and component equalities apply to every cube covariate collection and every spin assignment. Together with the calibrated information bound, these establish all assertions.
% lean: CausalSmith.Stat.TwosamplePointcateAnnotationFrontier.marked_component_certificate
\end{enumerate}
\end{proof}

\begin{proof}[Proof of \cref{obj:prop:annotation-threshold}]
Fix the public parameters in the stated domain. The frontier, oracle, and acquisition constants below are finite real numbers depending only on these parameters. In particular,
\[
\gamma>0,\qquad S>0,\qquad 2\gamma+d>0,\qquad \Delta>0.
\]
For an arbitrary auxiliary-size sequence, write
\[
N_n:=n+m_n\qquad(n\in\mathbb N).
\]
All asymptotic arguments may be restricted to the tail \(n\ge2\).

\begin{enumerate}
\item\label{proof:annotation-threshold-risk-rate}
By \cref{obj:thm:sharp-annotation-frontier,obj:thm:supplied-propensity}, choose constants satisfying
\[
0<c_F\le C_F<\infty,\qquad 0<c_O\le C_O<\infty
\]
such that, uniformly for \(n\ge2\) and \(m\ge0\),
\[
c_F r_*(n,m)\le R(n,m)\le C_F r_*(n,m),
\qquad
c_O r_o(n)\le R^e(n,m)\le C_O r_o(n).
\]
Since
\[
r_*(n,m)\ge r_o(n)=n^{-\gamma/(2\gamma+d)}>0,
\]
these bounds make both risks positive and finite. Dividing the upper bound for one risk by the lower bound for the other gives
\[
\begin{aligned}
0\le \frac{R(n,m)}{R^e(n,m)}
&\le \frac{C_F}{c_O}\frac{r_*(n,m)}{r_o(n)},\\
0\le \frac{r_*(n,m)}{r_o(n)}
&\le \frac{C_O}{c_F}\frac{R(n,m)}{R^e(n,m)}.
\end{aligned}
\]
The same calculation, using the frontier bounds at \((n,m)\) and \((n,0)\), gives
\[
\begin{aligned}
0\le \frac{R(n,m)}{R(n,0)}
&\le \frac{C_F}{c_F}\frac{r_*(n,m)}{r_*(n,0)},\\
0\le \frac{r_*(n,m)}{r_*(n,0)}
&\le \frac{C_F}{c_F}\frac{R(n,m)}{R(n,0)}.
\end{aligned}
\]
The first pair transfers eventual boundedness in both directions. The second pair transfers convergence to zero by squeezing between zero and a fixed multiple of the other ratio. Moreover, positivity of \(r_o(n)\) gives
\[
\frac{r_*(n,m_n)}{r_o(n)}=O(1)
\quad\Longleftrightarrow\quad
r_*(n,m_n)=O(r_o(n)).
\]
Consequently,
\[
\begin{aligned}
\frac{R(n,m_n)}{R^e(n,m_n)}=O(1)
&\quad\Longleftrightarrow\quad r_*(n,m_n)=O(r_o(n)),\\
\frac{R(n,m_n)}{R(n,0)}\longrightarrow0
&\quad\Longleftrightarrow\quad
\frac{r_*(n,m_n)}{r_*(n,0)}\longrightarrow0.
\end{aligned}
\]

If the rate comparison holds, there is a real constant \(K_{\mathrm{rate}}\) such that
\[
r_*(n,m_n)\le K_{\mathrm{rate}}r_o(n)
\qquad\text{eventually}.
\]
Set
\[
\zeta:=\max\{K_{\mathrm{rate}},1\}.
\]
Then \(1\le\zeta<\infty\), and the defining inequality in
\cref{obj:def:annotation-region} holds eventually. Conversely, eventual membership in that region supplies the rate comparison with constant \(\zeta\). Thus
\[
r_*(n,m_n)=O(r_o(n))
\quad\Longleftrightarrow\quad
\exists\,\zeta\in[1,\infty)\ 
\text{such that }(n,m_n)\in\mathcal B_*(\zeta)
\text{ eventually}.
\]
% lean: CausalSmith.Stat.TwosamplePointcateAnnotationFrontier.annotation_risk_rate_equivalences

\item
For the extended nonnegative lower limit, the tail characterization is
\[
0<\liminf_{n\to\infty}\frac{N_n}{n^{q_*}}
\quad\Longleftrightarrow\quad
\exists\,c_{\mathrm{count}}>0\
\text{such that }
c_{\mathrm{count}}\le\frac{N_n}{n^{q_*}}
\text{ eventually}.
\]
Indeed, a positive lower limit admits a smaller positive finite real number, which bounds every sufficiently late term from below; this also applies when the lower limit is \(+\infty\). Conversely, an eventual positive lower bound bounds the lower limit from below.

We now translate annotation-region membership into precisely such a bound. Directly from the definitions,
\[
q_*=\frac{S_{\mathrm{crit}}}{S},
\qquad
\Delta=(2\gamma+d)(1+q_*).
\]
For \(n\ge2\), this yields
\[
r_o(n)^{-\Delta/\gamma}
=n^{\Delta/(2\gamma+d)}
=n^{1+q_*}
=n\,n^{q_*}.
\]
Let \(\zeta\ge1\). Since \(r_o(n)\le\zeta r_o(n)\), the maximum defining \(r_*\) and
\cref{obj:def:annotation-region} give
\[
\begin{aligned}
(n,m)\in\mathcal B_*(\zeta)
&\quad\Longleftrightarrow\quad
(nN)^{-\gamma/\Delta}\le\zeta r_o(n)\\
&\quad\Longleftrightarrow\quad
(\zeta r_o(n))^{-\Delta/\gamma}\le nN\\
&\quad\Longleftrightarrow\quad
\zeta^{-\Delta/\gamma}\le\frac{N}{n^{q_*}}.
\end{aligned}
\]
Here the second equivalence uses the strictly decreasing positive-power map with exponent \(-\gamma/\Delta\), whose inverse has exponent \(-\Delta/\gamma\); all bases are positive.

Eventual membership therefore supplies the positive lower bound
\[
c_{\mathrm{count}}:=\zeta^{-\Delta/\gamma}>0.
\]
In the reverse direction, suppose an eventual lower bound \(c_{\mathrm{count}}>0\) is given, and choose
\[
\zeta:=\max\left\{1,c_{\mathrm{count}}^{-\gamma/\Delta}\right\}.
\]
Then
\[
1\le\zeta<\infty,
\qquad
\zeta^{-\Delta/\gamma}
\le
\left(c_{\mathrm{count}}^{-\gamma/\Delta}\right)^{-\Delta/\gamma}
=c_{\mathrm{count}}.
\]
The displayed count characterization gives eventual membership in
\(\mathcal B_*(\zeta)\). Combining this with \cref{proof:annotation-threshold-risk-rate} proves all four oracle-order equivalences.
% lean: CausalSmith.Stat.TwosamplePointcateAnnotationFrontier.annotationRegion_liminf_iff

\item
To characterize strict improvement, first suppose \(S<S_{\mathrm{crit}}\). The identities above imply
\[
q_*=\frac{S_{\mathrm{crit}}}{S}>1,
\qquad
\Delta=(2\gamma+d)(1+q_*)>2(2\gamma+d),
\]
and hence
\[
\frac{2\gamma}{\Delta}<\frac{\gamma}{2\gamma+d}.
\]
Define the positive decay exponent
\[
\eta_{\mathrm{dec}}
:=\frac{\gamma}{2\gamma+d}-\frac{2\gamma}{\Delta}>0.
\]
For \(n\ge2\), the inequality \(q_*>1\) gives
\[
n=n^1\le n^{q_*}.
\]
The lower-count branch in
\cref{obj:thm:sharp-annotation-frontier}, applied with \(m=0\), therefore gives
\[
r_*(n,0)=(n^2)^{-\gamma/\Delta}
=n^{-2\gamma/\Delta}.
\]
Dividing the two terms in the maximum defining \(r_*(n,m_n)\) by this positive rate yields
\[
\frac{r_o(n)}{n^{-2\gamma/\Delta}}
=n^{-\eta_{\mathrm{dec}}},
\qquad
\frac{(nN_n)^{-\gamma/\Delta}}{n^{-2\gamma/\Delta}}
=\left(\frac{N_n}{n}\right)^{-\gamma/\Delta}.
\]
Thus
\[
\frac{r_*(n,m_n)}{r_*(n,0)}
=
\max\left\{
n^{-\eta_{\mathrm{dec}}},
\left(\frac{N_n}{n}\right)^{-\gamma/\Delta}
\right\}.
\]

The second term tends to zero exactly when \(N_n/n\to+\infty\). To see the reverse implication explicitly, let
\[
b_{\mathrm{thr}}\in\mathbb R,
\qquad
x_{\mathrm{thr}}:=\max\{1,b_{\mathrm{thr}}\}>0.
\]
If the second term tends to zero, then eventually
\[
\left(\frac{N_n}{n}\right)^{-\gamma/\Delta}
\le x_{\mathrm{thr}}^{-\gamma/\Delta}.
\]
Inverting this decreasing power gives
\[
\frac{N_n}{n}
\ge
\left(x_{\mathrm{thr}}^{-\gamma/\Delta}\right)^{-\Delta/\gamma}
=x_{\mathrm{thr}}
\ge b_{\mathrm{thr}}.
\]
This proves divergence to \(+\infty\). The forward implication follows from the convergence of a strictly negative power to zero at \(+\infty\). Since
\[
\frac{N_n}{n}=1+\frac{m_n}{n}\qquad(n\ge2),
\]
divergence of the total-to-labeled ratio is equivalent to divergence of the auxiliary-to-labeled ratio: an eventual lower bound at any real threshold for either ratio gives the corresponding bound for the other after shifting the threshold by one.

The first term \(n^{-\eta_{\mathrm{dec}}}\) tends to zero. Both terms are nonnegative, so convergence of their maximum to zero implies convergence of the second term by squeezing; conversely, convergence of both terms implies convergence of their maximum. We have proved, in this regime,
\[
\frac{r_*(n,m_n)}{r_*(n,0)}\longrightarrow0
\quad\Longleftrightarrow\quad
\frac{m_n}{n}\longrightarrow+\infty.
\]

For the complementary case \(S\ge S_{\mathrm{crit}}\),
\cref{obj:thm:sharp-annotation-frontier} gives, for every \(n\ge2\),
\[
r_*(n,m_n)=r_*(n,0)=r_o(n)>0,
\qquad
\frac{r_*(n,m_n)}{r_*(n,0)}=1.
\]
The ratio therefore has limit one, which excludes convergence to zero. Combining the two cases and the risk-to-rate equivalence proves
\[
\frac{R(n,m_n)}{R(n,0)}\longrightarrow0
\quad\Longleftrightarrow\quad
\frac{r_*(n,m_n)}{r_*(n,0)}\longrightarrow0
\quad\Longleftrightarrow\quad
S<S_{\mathrm{crit}}
\ \text{and}\ 
\frac{m_n}{n}\longrightarrow+\infty.
\]
% lean: CausalSmith.Stat.TwosamplePointcateAnnotationFrontier.annotation_rate_improvement_iff

\item
The claimed exponent comparison follows from the same positive-denominator identity:
\[
S<S_{\mathrm{crit}}
\quad\Longrightarrow\quad
q_*=\frac{S_{\mathrm{crit}}}{S}>1,
\qquad S>0.
\]
% lean: CausalSmith.Stat.TwosamplePointcateAnnotationFrontier.qStar_gt_one

\item
Suppose \(m_n/n=O(1)\). Choose a real constant \(K_{\mathrm{aux}}\) satisfying
\[
\frac{m_n}{n}\le K_{\mathrm{aux}}
\qquad\text{eventually},
\]
and define
\[
K_{\mathrm{aux}}^+:=\max\{K_{\mathrm{aux}},0\},
\qquad
M_{\mathrm{aux}}
:=(1+K_{\mathrm{aux}}^+)^{\gamma/\Delta}\ge1.
\]
Eventually \(m_n\le K_{\mathrm{aux}}^+n\), and consequently
\[
n^2\le nN_n\le(1+K_{\mathrm{aux}}^+)n^2.
\]
Taking the negative power \(-\gamma/\Delta\) gives
\[
(nN_n)^{-\gamma/\Delta}\le(n^2)^{-\gamma/\Delta},
\]
and
\[
M_{\mathrm{aux}}^{-1}(n^2)^{-\gamma/\Delta}
\le(nN_n)^{-\gamma/\Delta},
\qquad
(n^2)^{-\gamma/\Delta}
\le M_{\mathrm{aux}}(nN_n)^{-\gamma/\Delta}.
\]
The oracle term is common to both rates and satisfies
\[
r_o(n)\le M_{\mathrm{aux}}r_o(n).
\]
Taking maxima therefore gives
\[
r_*(n,m_n)\le r_*(n,0)
\le M_{\mathrm{aux}}r_*(n,m_n)
\qquad\text{eventually}.
\]
Using the risk sandwiches,
\[
\begin{aligned}
R(n,m_n)
&\le C_F r_*(n,m_n)
\le C_F r_*(n,0)
\le \frac{C_F}{c_F}R(n,0),\\
R(n,0)
&\le C_F r_*(n,0)
\le C_F M_{\mathrm{aux}}r_*(n,m_n)
\le \frac{C_F M_{\mathrm{aux}}}{c_F}R(n,m_n).
\end{aligned}
\]
These eventual inequalities prove both stated big-\(O\) comparisons.
% lean: CausalSmith.Stat.TwosamplePointcateAnnotationFrontier.bounded_auxiliary_risk_order

\item
If \(S\ge S_{\mathrm{crit}}\), the frontier rate equals \(r_o(n)\). For every \(n\ge2\) and \(m\ge0\), the frontier upper bound and the oracle lower bound consequently give
\[
R(n,m)
\le C_F r_o(n)
=\frac{C_F}{c_O}\bigl(c_O r_o(n)\bigr)
\le\frac{C_F}{c_O}R^e(n,m).
\]
Dividing by the positive oracle risk yields
\[
0\le\frac{R(n,m)}{R^e(n,m)}
\le\frac{C_F}{c_O}.
\]
This uniform bound proves the oracle-order comparison for every auxiliary-size sequence.
% lean: CausalSmith.Stat.TwosamplePointcateAnnotationFrontier.smooth_oracle_risk_ratio

\item
For the finite-sample conclusions,
\cref{obj:thm:sharp-annotation-frontier} supplies, simultaneously for all \(n\ge2\) and \(m\ge0\),
\[
c_F r_*(n,m)\le R(n,m)
\le\sup_{P\in\mathcal M}\mathcal R_{n,m}(T_*,P)
\le C_F r_*(n,m).
\]
In particular, its upper bound follows by comparing the minimax value with the worst-case risk of \(T_*\). The same cited statement supplies exactly
\[
r_*(n,m)=
\begin{cases}
r_o(n),&S\ge S_{\mathrm{crit}},\\
(nN)^{-\gamma/\Delta},
&S<S_{\mathrm{crit}},\quad N\le n^{q_*},\\
r_o(n),
&S<S_{\mathrm{crit}},\quad N\ge n^{q_*},
\end{cases}
\]
together with
\[
N=n^{q_*}
\quad\Longrightarrow\quad
(nN)^{-\gamma/\Delta}=r_o(n).
\]
Thus \(c_F,C_F\) furnish the asserted finite-sample constants, including the equality of the two count branches at their boundary.
% lean: CausalSmith.Stat.TwosamplePointcateAnnotationFrontier.sharp_annotation_frontier

\item
For \(n\ge2\), specialize to \(m=0\). If \(S<S_{\mathrm{crit}}\), then \(q_*>1\), so
\[
N=n\le n^{q_*},
\qquad
r_*(n,0)=(n^2)^{-\gamma/\Delta}
=n^{-2\gamma/\Delta}.
\]
If \(S\ge S_{\mathrm{crit}}\), the first branch of
\cref{obj:thm:sharp-annotation-frontier} gives \(r_*(n,0)=r_o(n)\). Hence
\[
r_*(n,0)=
\begin{cases}
n^{-2\gamma/\Delta},&S<S_{\mathrm{crit}},\\
r_o(n),&S\ge S_{\mathrm{crit}}.
\end{cases}
\]
% lean: CausalSmith.Stat.TwosamplePointcateAnnotationFrontier.supervised_rate_branches

\item
Define the average nuisance smoothness by
\[
s_{\mathrm{pub}}:=\frac S2>0.
\]
By \cref{obj:lem:published-cate-benchmark}, algebra with positive denominators gives
\[
\frac{d/4}{1+d/(2\gamma)}
=\frac{S_{\mathrm{crit}}}{2},
\qquad
\frac{1}{1+d/(2\gamma)+d/(4s_{\mathrm{pub}})}
=\frac{2\gamma}{\Delta},
\qquad
\frac{1}{2+d/\gamma}
=\frac{\gamma}{2\gamma+d}.
\]
In particular,
\[
s_{\mathrm{pub}}<\frac{d/4}{1+d/(2\gamma)}
\quad\Longleftrightarrow\quad
S<S_{\mathrm{crit}}.
\]
For \(n\ge2\), these identities identify the numerical benchmark's strict split and its two exponents with the supervised branches just established.

The remaining sample sizes are \(n=0\) and \(n=1\). At \(n=1\), all power bases equal one, so both the benchmark and \(r_*(1,0)\) equal one. At zero, the numerical formulas use the convention
\[
0^{t_{\mathrm{pow}}}=0
\qquad
\bigl(t_{\mathrm{pow}}\in\mathbb R\setminus\{0\}\bigr).
\]
All exponents occurring here are strictly negative:
\[
-\frac{\gamma}{2\gamma+d}<0,
\qquad
-\frac{\gamma}{\Delta}<0,
\qquad
-\frac{2\gamma}{\Delta}<0.
\]
Thus both branches of the benchmark equal zero at \(n=0\), and the two terms in the maximum defining \(r_*(0,0)\) also equal zero. This proves the benchmark equality for every \(n\in\mathbb N\), with the stated formulas for \(n\ge1\).
% lean: CausalSmith.Stat.TwosamplePointcateAnnotationFrontier.published_benchmark_eq_supervised_rate

\item
For acquisition, define the positive exponents
\[
p_0:=\frac{2\gamma+d}{\gamma}
=2+\frac d\gamma,
\qquad
p_1:=\frac{\Delta}{\gamma}
=2+\frac d\gamma+\frac dS.
\]
For positive \(\eta\), the power-scaling identity is
\[
\left(\frac{\eta}{c_{\mathrm{scale}}}\right)^{-p_{\mathrm{scale}}}
=
c_{\mathrm{scale}}^{p_{\mathrm{scale}}}
\eta^{-p_{\mathrm{scale}}}
\qquad
\left(
c_{\mathrm{scale}}\in\{c_F,C_F\},\
p_{\mathrm{scale}}\in\{p_0,p_1\}
\right).
\]
Choose
\[
c_{\mathrm{acq}}
:=\min\{c_F^{p_0},c_F^{p_1}\}>0,
\qquad
C_{\mathrm{acq}}
:=\max\{C_F^{p_0},C_F^{p_1}\}>0.
\]

For any positive \(\eta\), \(n\ge2\), and \(m\ge0\), the defining maximum in
\cref{obj:def:sharp-rate} gives
\[
\begin{aligned}
r_*(n,m)\le\eta
&\quad\Longleftrightarrow\quad
n^{-\gamma/(2\gamma+d)}\le\eta
\ \text{and}\
(nN)^{-\gamma/\Delta}\le\eta\\
&\quad\Longleftrightarrow\quad
\eta^{-p_0}\le n
\ \text{and}\
\eta^{-p_1}\le nN.
\end{aligned}
\]
The second equivalence again inverts strictly decreasing powers of positive bases.

Now let \(0<\eta<1\), and suppose a Borel decision \(T\) satisfies
\[
\mathcal R_{n,m}(T,P)\le\eta
\qquad(P\in\mathcal M).
\]
By \cref{obj:def:minimax} and the frontier lower bound,
\[
c_F r_*(n,m)\le R(n,m)
\le\sup_{P\in\mathcal M}\mathcal R_{n,m}(T,P)
\le\eta.
\]
Therefore \(r_*(n,m)\le\eta/c_F\). Apply the exact rate equivalence at the positive scale \(\eta/c_F\), and then the scaling identity, to obtain
\[
c_F^{p_0}\eta^{-p_0}\le n,
\qquad
c_F^{p_1}\eta^{-p_1}\le nN.
\]
The choice of the minimum gives the common necessity constant:
\[
c_{\mathrm{acq}}\eta^{-p_0}
\le c_F^{p_0}\eta^{-p_0}\le n,
\qquad
c_{\mathrm{acq}}\eta^{-p_1}
\le c_F^{p_1}\eta^{-p_1}\le nN.
\]

Conversely, suppose
\[
C_{\mathrm{acq}}\eta^{-p_0}\le n,
\qquad
C_{\mathrm{acq}}\eta^{-p_1}\le nN.
\]
By the choice of the maximum and the scaling identity,
\[
\left(\frac{\eta}{C_F}\right)^{-p_0}
=C_F^{p_0}\eta^{-p_0}
\le C_{\mathrm{acq}}\eta^{-p_0}\le n,
\]
and
\[
\left(\frac{\eta}{C_F}\right)^{-p_1}
=C_F^{p_1}\eta^{-p_1}
\le C_{\mathrm{acq}}\eta^{-p_1}\le nN.
\]
The exact equivalence at scale \(\eta/C_F\) gives
\[
r_*(n,m)\le\frac{\eta}{C_F}.
\]
By \cref{obj:thm:sharp-annotation-frontier}, the decision \(T_*\) from
\cref{obj:def:sharp-decision} is Borel measurable, and, for every \(P\in\mathcal M\),
\[
\mathcal R_{n,m}(T_*,P)
\le\sup_{P'\in\mathcal M}\mathcal R_{n,m}(T_*,P')
\le C_F r_*(n,m)
\le\eta.
\]
Thus \(c_{\mathrm{acq}},C_{\mathrm{acq}}\) supply the asserted acquisition constants. Finally,
\[
n\le n+m=N
\]
because \(m\ge0\).
% lean: CausalSmith.Stat.TwosamplePointcateAnnotationFrontier.annotation_acquisition
\end{enumerate}
\end{proof}

\bibliographystyle{plainnat}

\bibliography{references}

\end{document}
